% I :Produire un plan détaillé d’une dissertation — Français Tle Tech — Cours signé OnBuch+
% Programme officiel MINESEC. Compiler avec : tectonic N13-i-produire-plan-detaille-dissertation.tex
\newcommand{\DOCMATIERE}{Français}
\newcommand{\DOCNIVEAU}{Tle Tech}
\newcommand{\DOCMODULE}{Français – classe de Terminale technique}
\newcommand{\DOCLECON}{N13}
\newcommand{\DOCTITRE}{I :Produire un plan détaillé d’une dissertation}
% OnBuch+ — préambule « cours signés » ENRICHI (compilé avec tectonic / XeLaTeX)
% Système visuel « Pop » d'OnBuch+ : fond crème (jamais blanc), bordures encre
% épaisses, ombres « dures » décalées, titres Archivo Black en capitales.
% Version MATHÉMATIQUES : graphes (pgfplots/tikz), boîtes pédagogiques
% (définition, propriété, méthode, attention, le savais-tu, exercice résolu,
% exercice, TP, bilan), page de garde et signature OnBuch+ ancrée en bas.
%
% Macros attendues dans main.tex AVANT \input{preamble} :
%   \DOCMATIERE  \DOCNIVEAU  \DOCMODULE  \DOCLECON (numéro)  \DOCTITRE
\documentclass[11pt]{article}

\usepackage{fontspec}
\usepackage[french]{babel}
\usepackage[a4paper,margin=2cm,top=3.4cm,bottom=2.8cm,headheight=1.4cm,footskip=1.3cm]{geometry}
\usepackage{amsmath,amssymb,amsfonts}
\usepackage{mathtools} % \xmapsto, \coloneqq... souvent utilisés par les modèles
\usepackage{cancel} % simplifications barrées (bilans de forces)
% \abs{z} (module, valeur absolue) et \norm{u} (norme), utilisés par les modèles
\providecommand{\abs}{}\let\abs\relax\DeclarePairedDelimiter{\abs}{\lvert}{\rvert}
\providecommand{\norm}{}\let\norm\relax\DeclarePairedDelimiter{\norm}{\lVert}{\rVert}
\usepackage[table]{xcolor} % [table] : fournit aussi \rowcolors (alternance de lignes)
\usepackage{pagecolor}
\usepackage{tikz}
\usetikzlibrary{arrows.meta,positioning,calc,shapes.geometric,shapes.misc,shapes.arrows,decorations.pathmorphing,decorations.pathreplacing,decorations.markings,patterns,shadows,mindmap,trees,fit,backgrounds,matrix,intersections,angles,quotes}
\usepackage{pgfplots}
\usepackage[version=4]{mhchem} % équations nucléaires \ce{^{235}_{92}U}
\usepackage[european,siunitx]{circuitikz} % schémas électriques
\pgfplotsset{compat=1.18}
\usepgfplotslibrary{fillbetween,groupplots} % aires entre courbes (calcul intégral)
\usepackage{enumitem}
\usepackage{fancyhdr}
% [most] charge toutes les bibliothèques tcolorbox (listings, minted, raster,
% documentation...) et gonfle inutilement la mémoire TeX sur un document long
% (~300 boîtes) : on ne charge que ce qui est réellement utilisé.
\usepackage[skins,breakable]{tcolorbox}
\usepackage{graphicx}
\usepackage{colortbl}
\usepackage{booktabs}
\usepackage{tabularx}
\usepackage{multirow}
\usepackage{diagbox} % case d'en-tête diagonale des tableaux à double entrée
% Colonnes extensibles centrée / gauche / droite (C, L, R), souvent utilisées par les modèles
\newcolumntype{C}{>{\centering\arraybackslash}X}
\newcolumntype{L}{>{\raggedright\arraybackslash}X}
\newcolumntype{R}{>{\raggedleft\arraybackslash}X}
\usepackage{array}
\usepackage{multicol}
\usepackage{siunitx}
% parse-numbers=false : accepte aussi \SI{2,5 \times 10^{-3}}{...}
\sisetup{locale=FR,output-decimal-marker={,},group-minimum-digits=4,parse-numbers=false}
\DeclareSIUnit\ohm{\ensuremath{\symup{\Omega}}} % U+2126 absent de la police
\DeclareSIUnit\division{div} % réglages d'oscilloscope (ms/div, V/div)
\DeclareSIUnit\tr{tr} % tours (vitesse de rotation, ex. \SI{1500}{\tr\per\minute})
\DeclareSIUnit\molar{M} % concentration molaire (ex. \SI{3}{\molar}, \SI{1}{\micro\molar})
\DeclareSIUnit\an{an} % année (le modèle confond parfois avec \year, un compteur TeX) — ex. \SI{5}{\kilo\meter\per\an}
\DeclareSIUnit\FCFA{FCFA} % franc CFA (ex. \SI{500}{\FCFA\per\kilo\gram})
\DeclareSIUnit\degreeBrix{\textdegree Brix} % teneur en sucre d'un sirop/jus (réfractométrie)
\DeclareSIUnit\month{mois} % le modèle utilise parfois \month, un compteur TeX, comme unité de durée
\DeclareSIUnit\microliter{\micro\liter} % le modèle écrit parfois \microliter en un seul token
\DeclareSIUnit\million{millions} % dénombrement (ex. \SI{15}{\million\per\mL} spermatozoïdes)
\usepackage{qrcode}
% \degree : commande fréquemment utilisée par les modèles pour un angle ou une
% température en dehors de \SI{}{} (non fournie par les paquets chargés).
\providecommand{\degree}{^{\circ}}
% \qed (fin de démonstration), utilisé par les modèles sans amsthm
\providecommand{\qed}{\hfill$\square$}
% \barre : double barre des valeurs interdites dans un tableau de variations (tkz-tab)
\providecommand{\barre}{\ensuremath{\|}}
% environnement proof (amsthm non chargé) utilisé par les modèles
\ifdefined\proof\else\newenvironment{proof}[1][Démonstration]{\par\noindent\textit{#1.}\ }{\qed\par}\fi
\usepackage{lastpage}
\usepackage{hyperref}
\hypersetup{hidelinks}

% ============================================================
% Palette « Pop » OnBuch+ — mêmes hex que la classe P (plus_ui.dart)
% ============================================================
\definecolor{poporange}{HTML}{F59321}
\definecolor{popdark}{HTML}{E07A0C}
\definecolor{popcream}{HTML}{FFF6E8}
\definecolor{popink}{HTML}{1C1714}
\definecolor{poppurple}{HTML}{7A5AE0}
\definecolor{popgreen}{HTML}{2E9E6A}
\definecolor{poppink}{HTML}{E0457B}
\definecolor{popblue}{HTML}{2F7DE1}
\definecolor{popgold}{HTML}{C98A12}
\definecolor{popmuted}{HTML}{8A7F76}
\definecolor{popline}{HTML}{EADFD2}
\definecolor{popgray}{HTML}{8A7F76}
\colorlet{popgrey}{popgray}
% Alias tolérants (le modèle nomme parfois les couleurs différemment)
\colorlet{popwhite}{white}
\colorlet{popblack}{popink}
\colorlet{popred}{poppink}
\colorlet{popyellow}{popgold}
% Teintes claires (fonds de tableaux/encadrés) — le modèle nomme parfois
% les couleurs avec un suffixe L (ex. popblueL) sans les avoir définies.
\colorlet{poporangeL}{poporange!20}
\colorlet{popdarkL}{popdark!20}
\colorlet{poppurpleL}{poppurple!20}
\colorlet{popgreenL}{popgreen!20}
\colorlet{poppinkL}{poppink!20}
\colorlet{popblueL}{popblue!20}
\colorlet{popgoldL}{popgold!20}
\colorlet{popredL}{popred!20}
\colorlet{popgrayL}{popgray!20}
% Teintes claires pour fonds de tableaux / zones
\colorlet{popblueL}{popblue!10!white}
\colorlet{poporangeL}{poporange!14!white}
\colorlet{popgreenL}{popgreen!12!white}
\colorlet{poppurpleL}{poppurple!12!white}
\colorlet{poppinkL}{poppink!12!white}
\colorlet{popgoldL}{popgold!14!white}

\pagecolor{popcream}

% ============================================================
% Typographie
% ============================================================
\defaultfontfeatures{Ligatures=TeX}
\setmainfont{PlusJakartaSans-Medium.ttf}[Path=../../../fonts/,BoldFont=PlusJakartaSans-Bold.ttf]
% Maths en sans-serif (Fira Math) avec chiffres et lettres droites de Plus
% Jakarta, pour que les formules se fondent dans le texte.
\usepackage[math-style=ISO,bold-style=ISO]{unicode-math}
\setmathfont{FiraMath-Regular.otf}[Path=../../../fonts/]
\setmathfont{PlusJakartaSans-Medium.ttf}[Path=../../../fonts/,range={up/num,up/{Latin,latin},\mathup}]
\usepackage{icomma} % virgule décimale sans espace en mode math
% Caractères mathématiques Unicode absents des polices de texte (Plus Jakarta,
% Archivo Black) : rendus par leur équivalent mathématique, en texte comme en
% formule — sinon ils s'affichent en carré vide (ex. « Arithmétique dans ℤ »).
\usepackage{newunicodechar}
\newunicodechar{ℝ}{\ensuremath{\mathbb{R}}}
\newunicodechar{ℤ}{\ensuremath{\mathbb{Z}}}
\newunicodechar{ℂ}{\ensuremath{\mathbb{C}}}
\newunicodechar{ℕ}{\ensuremath{\mathbb{N}}}
\newunicodechar{ℚ}{\ensuremath{\mathbb{Q}}}
\newunicodechar{𝔻}{\ensuremath{\mathbb{D}}}
\newunicodechar{α}{\ensuremath{\alpha}}
\newunicodechar{β}{\ensuremath{\beta}}
\newunicodechar{γ}{\ensuremath{\gamma}}
\newunicodechar{δ}{\ensuremath{\delta}}
\newunicodechar{ε}{\ensuremath{\varepsilon}}
\newunicodechar{θ}{\ensuremath{\theta}}
\newunicodechar{λ}{\ensuremath{\lambda}}
\newunicodechar{μ}{\ensuremath{\mu}}
\newunicodechar{σ}{\ensuremath{\sigma}}
\newunicodechar{τ}{\ensuremath{\tau}}
\newunicodechar{φ}{\ensuremath{\varphi}}
\newunicodechar{ψ}{\ensuremath{\psi}}
\newunicodechar{ω}{\ensuremath{\omega}}
\newunicodechar{Σ}{\ensuremath{\Sigma}}
% majuscules produites par \MakeUppercase dans les titres (α-aminés -> Α)
\newunicodechar{Α}{\ensuremath{\alpha}}
\newunicodechar{Β}{\ensuremath{\beta}}
\newunicodechar{Γ}{\ensuremath{\Gamma}}
\newunicodechar{Δ}{\ensuremath{\Delta}}
\newunicodechar{Ε}{\ensuremath{\varepsilon}}
\newunicodechar{Θ}{\ensuremath{\Theta}}
\newunicodechar{Λ}{\ensuremath{\Lambda}}
\newunicodechar{Μ}{\ensuremath{\mu}}
\newunicodechar{Τ}{\ensuremath{\tau}}
\newunicodechar{Φ}{\ensuremath{\Phi}}
\newunicodechar{Ψ}{\ensuremath{\Psi}}
\newunicodechar{Ω}{\ensuremath{\Omega}}
\newunicodechar{↦}{\ensuremath{\mapsto}}
\newunicodechar{∈}{\ensuremath{\in}}
\newunicodechar{∉}{\ensuremath{\notin}}
\newunicodechar{∧}{\ensuremath{\wedge}}
\newunicodechar{⇔}{\ensuremath{\Leftrightarrow}}
\newunicodechar{⟺}{\ensuremath{\Longleftrightarrow}}
\newunicodechar{⇒}{\ensuremath{\Rightarrow}}
\newunicodechar{⟹}{\ensuremath{\Longrightarrow}}
\newunicodechar{∘}{\ensuremath{\circ}}
\newunicodechar{∪}{\ensuremath{\cup}}
\newunicodechar{∩}{\ensuremath{\cap}}
\newunicodechar{≡}{\ensuremath{\equiv}}
\newunicodechar{⊂}{\ensuremath{\subset}}
\newunicodechar{⊥}{\ensuremath{\perp}}
\newunicodechar{∃}{\ensuremath{\exists}}
\newunicodechar{∀}{\ensuremath{\forall}}
\newunicodechar{∛}{\ensuremath{\sqrt[3]{\,}}}
\newunicodechar{∜}{\ensuremath{\sqrt[4]{\,}}}
\newunicodechar{⃗}{\ensuremath{{}^{\to}}}
\newunicodechar{ⁿ}{\ifmmode^{n}\else\textsuperscript{\ensuremath{n}}\fi}
\newunicodechar{ˣ}{\ifmmode^{x}\else\textsuperscript{\ensuremath{x}}\fi}
\newunicodechar{ᵇ}{\ifmmode^{b}\else\textsuperscript{\ensuremath{b}}\fi}
\newunicodechar{ʳ}{\ifmmode^{r}\else\textsuperscript{\ensuremath{r}}\fi}
\newunicodechar{ᵘ}{\ifmmode^{u}\else\textsuperscript{\ensuremath{u}}\fi}
\newunicodechar{ʸ}{\ifmmode^{y}\else\textsuperscript{\ensuremath{y}}\fi}
\newunicodechar{ᵃ}{\ifmmode^{a}\else\textsuperscript{\ensuremath{a}}\fi}
\newunicodechar{ᵉ}{\ifmmode^{e}\else\textsuperscript{\ensuremath{e}}\fi}
\newunicodechar{ᵏ}{\ifmmode^{k}\else\textsuperscript{\ensuremath{k}}\fi}
\newunicodechar{ᵖ}{\ifmmode^{p}\else\textsuperscript{\ensuremath{p}}\fi}
\newunicodechar{ᴺ}{\ifmmode^{N}\else\textsuperscript{\ensuremath{N}}\fi}
\newunicodechar{ᶜ}{\ifmmode^{c}\else\textsuperscript{\ensuremath{c}}\fi}
\newunicodechar{ʰ}{\ifmmode^{h}\else\textsuperscript{\ensuremath{h}}\fi}
\newunicodechar{ᵠ}{\ifmmode^{q}\else\textsuperscript{\ensuremath{q}}\fi}
\newunicodechar{ₙ}{\ifmmode_{n}\else\textsubscript{\ensuremath{n}}\fi}
\newunicodechar{ₐ}{\ifmmode_{a}\else\textsubscript{\ensuremath{a}}\fi}
\newunicodechar{ᵢ}{\ifmmode_{i}\else\textsubscript{\ensuremath{i}}\fi}
\newunicodechar{ᵣ}{\ifmmode_{r}\else\textsubscript{\ensuremath{r}}\fi}
\newunicodechar{ₖ}{\ifmmode_{k}\else\textsubscript{\ensuremath{k}}\fi}
\newunicodechar{ᵦ}{\ifmmode_{\beta}\else\textsubscript{\ensuremath{\beta}}\fi}
\newunicodechar{ₓ}{\ifmmode_{x}\else\textsubscript{\ensuremath{x}}\fi}
\newfontfamily\popdisplay{ArchivoBlack-Regular.ttf}[Path=../../../fonts/]
\newfontfamily\poplogo{SpaceGrotesk-Bold.ttf}[Path=../../../fonts/]
\newfontfamily\popsemi{PlusJakartaSans-SemiBold.ttf}[Path=../../../fonts/]
\newfontfamily\popxbold{PlusJakartaSans-ExtraBold.ttf}[Path=../../../fonts/]
\newfontfamily\popxboldls{PlusJakartaSans-ExtraBold.ttf}[Path=../../../fonts/,LetterSpace=3.0]

\setlist[enumerate]{leftmargin=1.7em,itemsep=5pt,topsep=4pt}
\setlist[itemize]{leftmargin=1.7em,itemsep=4pt,topsep=4pt,label={\color{poporange}\textbullet}}
\setlength{\parindent}{0pt}
\setlength{\parskip}{5pt}
\renewcommand{\arraystretch}{1.25}

% pgfplots : style Pop par défaut
\pgfplotsset{
  every axis/.append style={axis line style={popink,line width=1pt},tick style={popink},
    grid style={popline,line width=0.5pt},label style={font=\small},tick label style={font=\footnotesize}},
  popcurve/.style={poporange,line width=1.8pt,no markers,smooth},
}
% Même style disponible dans un \draw TikZ simple (hors pgfplots)
\tikzset{popcurve/.style={poporange,line width=1.8pt}}
\tikzset{popgrid/.style={popline,very thin}}
\colorlet{popcurve}{poporange} % le nom du style est parfois utilisé comme couleur

% ============================================================
% Pill / squiggle
% ============================================================
\newcommand{\poppill}[3]{%
  \tikz[baseline=-0.55ex]\node[draw=popink,line width=1.2pt,rounded corners=2.6pt,%
    fill=#2,inner xsep=6.5pt,inner ysep=3pt]{%
    \popxboldls\fontsize{7}{7}\selectfont\color{#3}\MakeUppercase{#1}};%
}
\newcommand{\popsquiggle}[1]{%
  \tikz[baseline=-0.4ex]{
    \draw[#1,line width=2.6pt,line cap=round]
      plot [smooth] coordinates {(0,0.09) (0.34,0.01) (0.68,0.21) (1.02,0.01) (1.36,0.10)};
  }%
}
% Mot-clé surligné (marqueur orange sous le texte)
\newcommand{\cle}[1]{{\popxbold\color{popdark}#1}}

% ============================================================
% En-tête / pied
% ============================================================
\pagestyle{fancy}
\fancyhf{}
\renewcommand{\headrulewidth}{0pt}
\renewcommand{\footrulewidth}{0pt}
\lhead{\raisebox{-0.15ex}{\poplogo\fontsize{13}{13}\selectfont\color{popink}OnBuch\color{poporange}+}}
\rhead{\poppill{\DOCMATIERE\ \textbullet\ \DOCNIVEAU\ \textbullet\ Leçon \DOCLECON}{white}{popink}}
\lfoot{\popxbold\fontsize{7.6}{7.6}\selectfont\color{popmuted}\MakeUppercase{Cours signé OnBuch+}\ \textbullet\ {\popsemi\DOCTITRE}}
\rfoot{\tikz[baseline=-0.6ex]\node[rounded corners=8.5pt,fill=popink,minimum height=17pt,minimum width=17pt,inner xsep=5pt,inner sep=0]{\popxbold\fontsize{8}{8}\selectfont\color{popcream}\thepage\,/\,\pageref*{LastPage}};}

% ============================================================
% Titres
% ============================================================
\newcommand{\chaptitre}[2]{%
  \begin{tcolorbox}[enhanced,colback=poporange,colframe=popink,boxrule=2.6pt,arc=11pt,
      drop shadow=popink,left=15pt,right=15pt,top=12pt,bottom=11pt,before skip=3pt,after skip=18pt]
    \poppill{#2}{popink}{popcream}\par\vspace{9pt}
    {\raggedright\hyphenpenalty=10000\exhyphenpenalty=10000\popdisplay\fontsize{22}{24}\selectfont\color{popcream}\MakeUppercase{#1}\par}
  \end{tcolorbox}
}
% Section numérotée : pastille orange avec le numéro + titre Archivo
\newcounter{popsec}
\newcommand{\coursec}[1]{%
  \stepcounter{popsec}\setcounter{popsub}{0}%
  \par\vspace{16pt}\noindent
  \begin{minipage}{\linewidth}
  \tikz[baseline=-0.7ex]\node[draw=popink,line width=1.6pt,fill=poporange,rounded corners=4pt,minimum size=22pt,inner sep=0]{\popdisplay\fontsize{12}{12}\selectfont\color{popcream}\thepopsec};%
  \hspace{8pt}{\popdisplay\fontsize{15}{17}\selectfont\color{popink}\MakeUppercase{#1}}\par
  \vspace{4pt}\noindent{\color{poporange}\rule{36pt}{3.6pt}}
  \end{minipage}\par\nopagebreak\vspace{8pt}
}
\newcounter{popsub}
\newcommand{\courssub}[1]{%
  \stepcounter{popsub}%
  \par\vspace{9pt}\noindent{\popxbold\fontsize{11.5}{13}\selectfont\color{popink}{\color{poporange}\thepopsec.\thepopsub}\hspace{6pt}#1}\par\nopagebreak\vspace{4pt}
}

% ============================================================
% Boîtes pédagogiques — même recette : fond blanc, bordure épaisse colorée,
% ombre dure encre, badge plein. Titre optionnel [#1].
% ============================================================
\tcbset{popbase/.style={breakable,enhanced,colback=white,boxrule=2.3pt,arc=9pt,
  shadow={3pt}{-3pt}{0pt}{popink},left=14pt,right=14pt,top=11pt,bottom=11pt,before skip=11pt,after skip=15pt,
  fontupper=\popsemi\fontsize{10.6}{14.5}\selectfont\color{popink},
  fontlower=\popsemi\fontsize{10.6}{14.5}\selectfont\color{popink}}}
% Coche dessinée (le glyphe ✓ n'existe pas dans les polices du document)
\newcommand{\popcheck}[1]{\tikz[baseline=-0.5ex]\draw[#1,line width=1.6pt,line cap=round,line join=round](-0.13,0.0)--(-0.04,-0.09)--(0.13,0.1);}
\providecommand{\coche}{\popcheck{popgreen}}
\providecommand{\resultat}[1]{\par\smallskip\noindent\fcolorbox{popgreen}{popgreenL}{\parbox{\dimexpr\linewidth-2\fboxsep-2\fboxrule\relax}{\textbf{Résultat.} #1}}\par\smallskip}
\newcommand{\popboxhead}[3]{\poppill{#1}{#2}{white}\ifx\relax#3\relax\else\hspace{6pt}{\popxbold\fontsize{10.6}{12}\selectfont\color{popink}#3}\fi\par\vspace{7pt}}

\newtcolorbox{aretenir}[1][]{popbase,colframe=popgreen,before upper={\popboxhead{À retenir}{popgreen}{#1}}}
\newtcolorbox{exemplebox}[1][]{popbase,colframe=poporange,before upper={\popboxhead{Exemple}{poporange}{#1}}}
\newtcolorbox{definition}[1][]{popbase,colframe=popblue,before upper={\popboxhead{Définition}{popblue}{#1}}}
\newtcolorbox{propriete}[1][]{popbase,colframe=poppurple,before upper={\popboxhead{Propriété}{poppurple}{#1}}}
\newtcolorbox{methode}[1][]{popbase,colframe=popgold,before upper={\popboxhead{Méthode}{popgold}{#1}}}
\newtcolorbox{attention}[1][]{popbase,colframe=poppink,before upper={\popboxhead{Attention}{poppink}{#1}}}
\newtcolorbox{experience}[1][]{popbase,colframe=popblue,colback=popblueL,before upper={\popboxhead{Expérience}{popblue}{#1}}}
% « Le savais-tu ? » : carte violette pleine (contexte, culture, Cameroun)
\newtcolorbox{savaistu}[1][]{popbase,colback=poppurple,colframe=popink,
  fontupper=\popsemi\fontsize{10.4}{14.2}\selectfont\color{white},
  before upper={\poppill{Le savais-tu\,?}{white}{poppurple}\ifx\relax#1\relax\else\hspace{6pt}{\popxbold\color{white}#1}\fi\par\vspace{7pt}}}
% Exercice résolu : énoncé en haut, \tcblower puis la solution sur fond crème
\newcounter{popexo}\newcounter{popexr}
\newtcolorbox{exoresolu}[1][]{popbase,colframe=poporange,colbacklower=popcream,
  segmentation style={popink,line width=1.2pt,dashed},
  before upper={\stepcounter{popexr}\popboxhead{Exercice résolu \thepopexr}{poporange}{#1}},
  before lower={\poppill{Solution}{popink}{popcream}\par\vspace{6pt}}}
% Exercice (non résolu en place) — la correction est dans la partie Corrigés
\newtcolorbox{exercice}[1][]{popbase,colframe=popink,
  before upper={\stepcounter{popexo}\popboxhead{Exercice \thepopexo}{popink}{#1}}}
% Corrigé
\newtcolorbox{corrige}[1][]{popbase,colframe=popgreen,colback=popgreenL,
  before upper={\popboxhead{Corrigé}{popgreen}{#1}}}
% Objectifs / prérequis / bilan
\newtcolorbox{objectifs}{popbase,colframe=popink,colback=poporangeL,
  before upper={\popboxhead{Objectifs de la leçon}{popink}{}}}
\newtcolorbox{prerequis}{popbase,colframe=popink,colback=popblueL,
  before upper={\popboxhead{Prérequis}{popblue}{}}}
\newtcolorbox{bilan}{popbase,colframe=popink,colback=white,boxrule=2.8pt,
  before upper={\popboxhead{Fiche bilan}{poporange}{L'essentiel à connaître pour l'examen}}}
% Figure encadrée avec légende
\newtcolorbox{popfigure}[1][]{enhanced,breakable=false,colback=white,colframe=popline,boxrule=1.4pt,arc=8pt,
  left=10pt,right=10pt,top=10pt,bottom=8pt,before skip=10pt,after skip=12pt,halign=center,
  lower separated=false,halign lower=center,
  fontlower=\popsemi\fontsize{9}{11.5}\selectfont\color{popmuted},
  #1}
\newcommand{\legende}[1]{\par\vspace{6pt}{\popsemi\fontsize{9}{11.5}\selectfont\color{popmuted}#1}}

% ============================================================
% Signature OnBuch+
% ============================================================
\newcommand{\poplogoinline}[1]{{\poplogo\fontsize{#1}{#1}\selectfont\color{popink}OnBuch\color{poporange}+}}
\newcommand{\signatureonbuch}{%
  \begin{tcolorbox}[enhanced,colback=popink,colframe=popink,boxrule=2.2pt,arc=10pt,
      drop shadow=poporange,left=14pt,right=14pt,top=10pt,bottom=10pt,before skip=0pt,after skip=0pt]
    \tikz[baseline=-0.7ex]\node[circle,fill=popgreen,draw=popcream,line width=1.3pt,minimum size=22pt,inner sep=0]{\popcheck{white}};%
    \hspace{9pt}\begin{minipage}[c]{0.62\linewidth}
      {\popxbold\fontsize{12}{13}\selectfont\color{popcream}Signé\ \textbullet\ {\poplogo OnBuch\color{poporange}+}}\par\vspace{2pt}
      {\raggedright\popsemi\fontsize{8}{10.5}\selectfont\color{popline}\MakeUppercase{Équipe pédagogique OnBuch+}\\\MakeUppercase{Conforme au programme officiel MINESEC}\par}
    \end{minipage}\hfill
    {\poplogo\fontsize{22}{22}\selectfont\color{popcream}OnBuch\color{poporange}+}
  \end{tcolorbox}%
}

% ============================================================
% Page de garde
% #1 titre · #2 matière · #3 niveau · #4 module · #5 n° leçon · #6 durée ·
% #7 description / objectifs (texte ou itemize)
% ============================================================
\newcommand{\pagegarde}[7]{%
  \thispagestyle{empty}
  \newgeometry{a4paper,margin=1.7cm,top=1.6cm,bottom=1.4cm}%
  \begin{tikzpicture}[remember picture,overlay]
    \fill[poporange!18] ([xshift=-3.2cm,yshift=-3.6cm]current page.north east) circle (4.6cm);
    \fill[poppurple!14] ([xshift=2.4cm,yshift=7.6cm]current page.south west) circle (3.8cm);
  \end{tikzpicture}%
  {\poplogo\fontsize{30}{30}\selectfont\color{popink}OnBuch\color{poporange}+}\hfill
  \raisebox{0.6ex}{\poppill{Cours signé \textbullet\ Premium}{popink}{popcream}}\par
  \vspace{1pt}\popsquiggle{poppurple}\par
  \vspace{8pt}
  \begin{tcolorbox}[enhanced,colback=poporange,colframe=popink,boxrule=2.8pt,arc=13pt,
      drop shadow=popink,left=18pt,right=18pt,top=12pt,bottom=13pt,after skip=10pt]
    \poppill{#2 \textbullet\ #3}{popink}{popcream}\hspace{5pt}\poppill{#4}{white}{popink}\par\vspace{8pt}
    {\popdisplay\fontsize{14}{14}\selectfont\color{popink}LEÇON #5}\par\vspace{4pt}
    {\raggedright\hyphenpenalty=10000\exhyphenpenalty=10000\popdisplay\fontsize{25}{27}\selectfont\color{popcream}\MakeUppercase{#1}\par}
    \vspace{8pt}
    {\popxbold\fontsize{9.5}{9.5}\selectfont\color{popink}\MakeUppercase{Durée conseillée : #6}}
  \end{tcolorbox}
  \begin{tcolorbox}[enhanced,colback=white,colframe=popink,boxrule=2.2pt,arc=10pt,
      drop shadow=popink,left=16pt,right=16pt,top=10pt,bottom=10pt,after skip=8pt]
    \poppill{Dans cette leçon}{poppurple}{white}\par\vspace{5pt}
    \popsemi\fontsize{9.3}{12.6}\selectfont\color{popink}#7
  \end{tcolorbox}
  \noindent\begin{minipage}[t]{0.487\linewidth}
    \begin{tcolorbox}[enhanced,colback=popgreen,colframe=popink,boxrule=2.2pt,arc=10pt,
        drop shadow=popink,left=11pt,right=11pt,top=10pt,bottom=10pt,height=3.1cm,valign=center]
      \tcbox[colback=white,colframe=popink,boxrule=1.3pt,arc=5pt,boxsep=0pt,nobeforeafter,
        left=3pt,right=3pt,top=3pt,bottom=3pt,valign=center]{\qrcode[height=1.9cm]{https://play.google.com/store/apps/details?id=com.luvvix.onbuch&hl=fr}}%
      \hspace{8pt}\begin{minipage}[c]{0.5\linewidth}
        {\popdisplay\fontsize{11}{12.5}\selectfont\color{white}\MakeUppercase{Télécharge OnBuch}}\par\vspace{4pt}
        {\raggedright\popsemi\fontsize{8.4}{11}\selectfont\color{white}Gratuit \textbullet\ Google Play\\Tuteur IA, annales, concours, résultats\par}
      \end{minipage}
    \end{tcolorbox}
  \end{minipage}\hfill
  \begin{minipage}[t]{0.487\linewidth}
    \begin{tcolorbox}[enhanced,colback=poppurple,colframe=popink,boxrule=2.2pt,arc=10pt,
        drop shadow=popink,left=11pt,right=11pt,top=10pt,bottom=10pt,height=3.1cm,valign=center]
      \tikz[baseline=-0.7ex]\node[circle,fill=white,draw=popink,line width=1.3pt,minimum size=1.6cm,inner sep=0]{\popdisplay\fontsize{24}{24}\selectfont\color{poppurple}+};%
      \hspace{8pt}\begin{minipage}[c]{0.6\linewidth}
        {\popdisplay\fontsize{11}{12.5}\selectfont\color{white}\MakeUppercase{Passe à OnBuch+}}\par\vspace{4pt}
        {\raggedright\popsemi\fontsize{8.4}{11}\selectfont\color{white}Tous les cours signés, Tuteur IA illimité, fiches PDF — onglet OnBuch+ de l'app\par}
      \end{minipage}
    \end{tcolorbox}
  \end{minipage}
  \vspace{8pt}
  \popcontenu
  \vfill
  \signatureonbuch
  \restoregeometry
  \newpage
}

% Rangée « Ce cours contient » (page de garde)
\newcommand{\poptile}[3]{% #1 couleur #2 gros texte #3 légende
  \begin{tcolorbox}[enhanced,colback=#1,colframe=popink,boxrule=2pt,arc=9pt,shadow={2.5pt}{-2.5pt}{0pt}{popink},
      left=6pt,right=6pt,top=9pt,bottom=9pt,halign=center,valign=center,height=1.75cm,width=\linewidth,nobeforeafter]
    {\popdisplay\fontsize{12.5}{13}\selectfont\color{white}\MakeUppercase{#2}}\par\vspace{3pt}
    {\centering\popsemi\fontsize{7.8}{9.5}\selectfont\color{white}#3\par}
  \end{tcolorbox}}
\newcommand{\popcontenu}{%
  \par\noindent{\popxboldls\fontsize{7.5}{7.5}\selectfont\color{popmuted}CE COURS SIGNÉ CONTIENT}\par\vspace{7pt}
  \noindent\begin{minipage}[t]{0.235\linewidth}\poptile{popblue}{Cours}{complet, illustré, pas à pas}\end{minipage}\hfill
  \begin{minipage}[t]{0.235\linewidth}\poptile{poporange}{Méthodes}{et exercices résolus}\end{minipage}\hfill
  \begin{minipage}[t]{0.235\linewidth}\poptile{poppink}{Exercices}{type examen + corrigés}\end{minipage}\hfill
  \begin{minipage}[t]{0.235\linewidth}\poptile{popgreen}{Fiche bilan}{l'essentiel à retenir}\end{minipage}\par}

% Sommaire
\newtcolorbox{sommaire}{enhanced,colback=white,colframe=popink,boxrule=2.2pt,arc=10pt,
  drop shadow=popink,left=18pt,right=18pt,top=13pt,bottom=13pt,before skip=6pt,after skip=20pt,
  before upper={\poppill{Sommaire}{poppurple}{white}\par\vspace{9pt}\popsemi\fontsize{10.6}{14.5}\selectfont\color{popink}}}

% Signature de fin de cours — poussée tout en bas de la dernière page
\newcommand{\findecours}{%
  \par\vfill\nopagebreak
  \begin{center}\popsquiggle{poporange}\end{center}\vspace{4pt}
  \signatureonbuch
}
\AtBeginDocument{\def\llbracket{\mathopen{[\mkern-3.5mu[}}\def\rrbracket{\mathclose{]\mkern-3.5mu]}}} % intervalles d'entiers (glyphe ⟦ absent de la police)
% Glyphes absents de Fira Math : redéfinis à partir de glyphes disponibles
\AtBeginDocument{%
  \def\setminus{\mathbin{\backslash}}\let\smallsetminus\setminus
  \def\vdots{\mathord{\vbox{\baselineskip3.2pt\lineskiplimit0pt\kern4pt\hbox{.}\hbox{.}\hbox{.}}}}%
  \def\ddots{\mathinner{\mkern1mu\raise6.5pt\hbox{.}\mkern2mu\raise3.6pt\hbox{.}\mkern2mu\raise0.7pt\hbox{.}\mkern1mu}}%
  \def\triangle{\mathord{\scalebox{0.9}{$\symup{\Delta}$}}}%
  \def\nabla{\mathord{\raisebox{\depth}{\scalebox{1}[-1]{$\symup{\Delta}$}}}}%
  \def\checkmark{\popcheck{popdark}}%
  \def\star{\mathbin{\ast}}\def\bigstar{\mathord{\ast}}%
  \def\diamond{\mathbin{\scalebox{0.6}{$\Diamond$}}}%
  \def\bigcup{\mathop{\mathchoice{\raisebox{-0.3em}{\scalebox{1.7}{$\cup$}}}{\scalebox{1.25}{$\cup$}}{\cup}{\cup}}\displaylimits}%
  \def\bigcap{\mathop{\mathchoice{\raisebox{-0.3em}{\scalebox{1.7}{$\cap$}}}{\scalebox{1.25}{$\cap$}}{\cap}{\cap}}\displaylimits}%
  \def\lesseqqgtr{\mathrel{\lessgtr}}\def\bigoplus{\mathop{\oplus}\displaylimits}%
}
\providecommand{\ssi}{si et seulement si} % abréviation fréquente
\AtBeginDocument{\providecommand{\wideparen}{\overparen}} % arc de cercle

\begin{document}

\pagegarde{I :Produire un plan détaillé d’une dissertation}{Français}{Tle Tech}{Français – classe de Terminale technique}{N13}{16 séances (fiche de progression harmonisée)}{Cette leçon guide l'élève dans la maîtrise complète de la dissertation littéraire et philosophique : de l'analyse fine du sujet à la rédaction aboutie, en passant par l'élaboration d'un plan détaillé cohérent, le tout ancré dans l'étude de l'œuvre "Le Vieux Nègre et la Médaille" et le renforcement des outils linguistiques (cohésion, référents, lexique).
\vspace{4pt}
\begin{multicols}{2}\small
\begin{itemize}
\item À la fin de cette leçon, je sais analyser un sujet de dissertation (termes clés, présupposés, champ de réflexion) et en formuler une problématique pertinente.
\item À la fin de cette leçon, je sais mobiliser des techniques de recherche d'idées (brainstorming, mind-mapping, fiches de lecture) pour alimenter ma démonstration.
\item À la fin de cette leçon, je sais construire un plan détaillé (dialectique, thématique ou analytique) avec axes, sous-parties, arguments, exemples et transitions.
\item À la fin de cette leçon, je sais rédiger une introduction (accroche, présentation sujet, problématique, annonce de plan) et une conclusion (bilan, ouverture) conformes aux attendus du Baccalauréat.
\item À la fin de cette leçon, je sais rédiger un paragraphe argumenté structuré (thèse, argument, exemple/citation, lien) en assurant la cohésion par les substituts du référent et la progression thématique.
\item À la fin de cette leçon, je sais identifier et expliquer les sigles, emprunts et figures de style dans "Le Vieux Nègre et la Médaille" pour enrichir mon analyse littéraire.
\end{itemize}\end{multicols}}

\begin{sommaire}
\begin{multicols}{2}
\begin{enumerate}
\item 1. Entrée dans l'œuvre \& Outils linguistiques fondamentaux
\item 2. La dissertation : De l'analyse du sujet à la problématique
\item 3. L'architecture maîtrisée : Élaboration du plan détaillé
\item 4. Approfondissement littéraire \& Cohérence textuelle (Lectures méthodiques)
\item 5. De la plume à la copie : Rédaction, Intégration \& Remédiation
\item Activité d'intégration
\item Méthodes et exercices résolus
\item Exercices
\item Corrigés des exercices
\item Fiche bilan
\end{enumerate}
\end{multicols}
\end{sommaire}

\begin{objectifs}
\begin{itemize}
\item À la fin de cette leçon, je sais analyser un sujet de dissertation (termes clés, présupposés, champ de réflexion) et en formuler une problématique pertinente.
\item À la fin de cette leçon, je sais mobiliser des techniques de recherche d'idées (brainstorming, mind-mapping, fiches de lecture) pour alimenter ma démonstration.
\item À la fin de cette leçon, je sais construire un plan détaillé (dialectique, thématique ou analytique) avec axes, sous-parties, arguments, exemples et transitions.
\item À la fin de cette leçon, je sais rédiger une introduction (accroche, présentation sujet, problématique, annonce de plan) et une conclusion (bilan, ouverture) conformes aux attendus du Baccalauréat.
\item À la fin de cette leçon, je sais rédiger un paragraphe argumenté structuré (thèse, argument, exemple/citation, lien) en assurant la cohésion par les substituts du référent et la progression thématique.
\item À la fin de cette leçon, je sais identifier et expliquer les sigles, emprunts et figures de style dans "Le Vieux Nègre et la Médaille" pour enrichir mon analyse littéraire.
\item À la fin de cette leçon, je sais produire un compte rendu de lecture critique et négocier un projet d'étude sur l'œuvre au programme.
\item À la fin de cette leçon, je sais auto-évaluer ma production écrite à l'aide de grilles de critères (cohérence, syntaxe, orthographe, pertinence) pour la remédiation.
\end{itemize}
\end{objectifs}

\begin{prerequis}
\begin{itemize}
\item Structure de base du paragraphe argumentatif (thèse, argument, exemple) acquise en Première.
\item Notions de cohésion (connecteurs logiques, référents) et de cohérence (progression thématique) vues en début d'année.
\item Connaissance du contexte historique colonial au Cameroun (indigénat, travail forcé) pour l'étude de l'œuvre.
\item Maîtrise des temps de la narration (passé simple, imparfait, plus-que-parfait) pour l'analyse textuelle.
\item Technique de la prise de notes et du surlignage lors de la lecture d'œuvres intégrales.
\end{itemize}
\end{prerequis}

\newpage

\coursec{Entrée dans l'œuvre \& Outils linguistiques fondamentaux}

\textbf{Situation d'accroche.} Au lycée technique de Bamenda, Armande fixe le sujet du Bac blanc : « La tradition est-elle un frein au progrès ? ». Elle a les idées, elle connaît des exemples, elle a même lu \emph{Le Vieux Nègre et la Médaille} de Ferdinand Oyono. Pourtant, sa copie précédente a été sanctionnée par deux mentions rouges : « Plan incohérent » et « Pas de problématique ». Son professeur lui tend alors une grille d'analyse : « Avant de bâtir la maison, on trace le plan ; avant d'écrire, on structure sa pensée. » Comment Armande peut-elle transformer ses intuitions en une architecture rigoureuse qui impressionnera le correcteur ?

\textbf{Problématique de la section.} Toute bonne dissertation commence par une bonne lecture. Avant d'apprendre à construire un plan (sections suivantes), il faut savoir \emph{entrer dans une œuvre} : décrypter ses paratextes, identifier ses personnages et son cadre, suivre la chaîne des référents qui assure la cohérence du récit, et comprendre les mots empruntés qui ancrent le roman dans la réalité camerounaise. C'est l'objet de cette première section : donner à Armande — et à vous — les outils de lecture et de langue indispensables.

\courssub{Lire avant de lire : l'analyse des paratextes}

Quand vous tenez un livre entre les mains, vous ne commencez jamais par la première ligne du chapitre 1. Votre œil rencontre d'abord le titre, le nom de l'auteur, l'image de couverture, le résumé au dos. Ces éléments ne sont pas des décorations : ce sont des \emph{messages} que l'éditeur et l'auteur vous adressent pour orienter votre lecture.

\begin{definition}[Le paratexte]
Le \cle{paratexte} (notion forgée par Gérard Genette) est l'ensemble des éléments qui entourent le texte proprement dit et qui guident la lecture : le \textbf{titre}, le \textbf{nom de l'auteur}, la \textbf{couverture} (image, couleurs), la \textbf{préface}, la \textbf{dédicace}, la \textbf{quatrième de couverture}. On distingue le \emph{péritexte} (éléments dans le livre : titre, préface, dédicace) et l'\emph{épitexte} (éléments hors du livre : interviews, articles).
\end{definition}

Appliquons cette grille au \emph{Vieux Nègre et la Médaille} :

\begin{center}
\begin{tabularx}{\linewidth}{|l|X|}
\hline
\rowcolor{popblueL} \textbf{Paratexte} & \textbf{Hypothèses de lecture qu'il permet de formuler} \\
\hline
Titre : \emph{Le Vieux Nègre et la Médaille} & Deux pôles s'opposent : un homme âgé, défini par sa « négritude » (mot chargé de mépris colonial), et un objet, la médaille, symbole de la récompense coloniale. Antithèse annonciatrice d'un conflit. \\
\hline
Auteur : Ferdinand Oyono & Écrivain camerounais, auteur d'\emph{Une vie de boy} : on peut s'attendre à un roman de dénonciation de la colonisation. \\
\hline
Couverture (selon éditions) & Un vieil homme africain, parfois une médaille militaire : le contraste visuel confirme l'ironie du titre. \\
\hline
Dédicace & « À tous ceux qui... » : elle révèle souvent l'intention de l'auteur, son engagement. \\
\hline
Quatrième de couverture & Annonce Meka, vieux cultivateur, invité à Douala pour recevoir une médaille : hypothèse d'un récit à la fois comique et tragique. \\
\hline
\end{tabularx}
\end{center}

\begin{popfigure}
\begin{tikzpicture}[mindmap, grow cyclic, every node/.style={concept, text=white}, root concept/.append style={concept color=popdark, font=\large\bfseries}, level 1/.append style={level distance=3.6cm, sibling angle=72}, level 1 concept/.append style={concept color=popblue, font=\small\bfseries}, level 2/.append style={level distance=2.6cm, sibling angle=50}, level 2 concept/.append style={concept color=poporange, font=\scriptsize, text=white}]
\node{Roman}
  child { node {Titre} child { node {opposition homme / objet} } }
  child { node {Auteur} child { node {Oyono, engagé anticolonial} } }
  child { node {Couverture} child { node {image : ironie visuelle} } }
  child { node {Préface} child { node {intention de l'auteur} } }
  child { node {4\textsuperscript{e} de couverture} child { node {résumé, ton, genre} } };
\end{tikzpicture}
\legende{Carte mentale des paratextes : chaque élément autour du roman fournit des indices pour formuler des hypothèses de lecture.}
\end{popfigure}

\begin{aretenir}[La démarche en trois temps]
L'analyse des paratextes suit toujours le même mouvement : \textbf{observer} (que vois-je ?), \textbf{interpréter} (que signifie cet indice ?), \textbf{anticiper} (quelle lecture cela prépare-t-il ?). Ces hypothèses seront vérifiées ou corrigées pendant la lecture : c'est toute la méthode de la lecture méthodique.
\end{aretenir}

\courssub{La lecture ouvroir : première découverte du texte}

La \cle{lecture ouvroir} (lecture « qui ouvre » l'œuvre) est la première lecture globale, celle qui permet de situer le texte sans encore l'analyser en détail. Elle répond à quatre questions : \emph{quel genre ? qui raconte ? qui sont les personnages ? où et quand ?}

\begin{exemplebox}[Fiche d'identité de l'œuvre]
\begin{itemize}
\item \textbf{Titre} : \emph{Le Vieux Nègre et la Médaille}.
\item \textbf{Auteur} : Ferdinand Oyono (Cameroun, 1929-2010), écrivain et homme d'État.
\item \textbf{Date de publication} : 1956.
\item \textbf{Genre} : roman réaliste à visée satirique, proche du témoignage par son ancrage dans la réalité coloniale.
\item \textbf{Narrateur} : narrateur hétérodiégétique (extérieur à l'histoire, « il »), omniscient, qui adopte souvent le point de vue de Meka (focalisation interne).
\item \textbf{Personnages principaux} : \textbf{Meka}, vieux cultivateur, le « vieux nègre » du titre ; \textbf{Kelara}, son épouse ; le \textbf{Commandant}, représentant de l'administration coloniale ; le père Vandermayer, Engamba, les notables.
\item \textbf{Cadre spatio-temporel} : le village d'Eloumazock, puis la ville de \textbf{Douala}, au Cameroun, dans les \textbf{années 1950}, à la fin de la période coloniale.
\end{itemize}
\end{exemplebox}

\begin{savaistu}[Un roman prophétique ?]
Ferdinand Oyono a écrit ce roman en \textbf{1956}, soit deux ans avant l'indépendance du Cameroun (1960). C'est une charge satirique contre l'\textbf{assimilationnisme} français : l'idée que l'Africain devait devenir « un Français à peau noire » pour mériter la considération. La médaille offerte à Meka symbolise cette fausse récompense : on décore l'homme qu'on a spolié de ses terres. L'ironie d'Oyono a fait de ce roman un classique des programmes scolaires africains.
\end{savaistu}

\courssub{Outil de langue n°1 : le référent et ses substituts}

Pourquoi un texte est-il facile à suivre ? Parce que l'auteur ne répète jamais le même nom vingt fois de suite. Il le \emph{remplace} par des pronoms, des descriptions, des périphrases. C'est le système de la \cle{référence}.

\begin{definition}[Référent et substituts]
Le \cle{référent} est l'être ou l'objet dont on parle (la première mention, appelée \emph{antécédent}). Les \cle{substituts} sont tous les mots qui le reprennent ensuite : \textbf{pronoms} (il, le, celui-ci), \textbf{noms synonymes} (le vieillard), \textbf{périphrases} (le père de Kelara). Quand le substitut renvoie en arrière, on parle d'\cle{anaphore} ; quand il annonce ce qui vient, de \cle{cataphore}.
\end{definition}

\begin{exemplebox}[Chaîne de référence dans le roman]
« \textbf{Meka} se leva de bonne heure. \emph{Le vieux cultivateur} regarda ses mains calleuses. \emph{Il} pensa à la médaille qu'on allait lui remettre. \emph{Le récipiendaire}, comme l'appelait le Commandant, sourit malgré lui. \emph{Le père de Kelara} se sentait soudain important. »
\end{exemplebox}

\begin{popfigure}
\begin{tikzpicture}[node distance=0.9cm]
\node[draw=popdark, fill=popdark, text=white, rounded corners, font=\bfseries, minimum width=3.4cm] (ref) {Meka (référent)};
\node[draw=popblue, fill=popblueL, rounded corners, below left=1.1cm and 0.4cm of ref] (s1) {il};
\node[draw=popblue, fill=popblueL, rounded corners, below=1.1cm of ref] (s2) {le vieux cultivateur};
\node[draw=popblue, fill=popblueL, rounded corners, below right=1.1cm and 0.4cm of ref] (s3) {le récipiendaire};
\node[draw=popblue, fill=popblueL, rounded corners, below=3.4cm of ref] (s4) {le père de Kelara};
\draw[-{Stealth}, thick, poporange] (ref.south west) -- (s1.north) node[midway, left, font=\scriptsize\itshape, text=popmuted] {anaphore};
\draw[-{Stealth}, thick, poporange] (ref.south) -- (s2.north);
\draw[-{Stealth}, thick, poporange] (ref.south east) -- (s3.north);
\draw[-{Stealth}, thick, poporange] (s2.south) -- (s4.north);
\node[draw=popgreen, fill=popgreenL, rounded corners, right=2.6cm of s2, font=\small, text width=4.2cm] (cat) {Cataphore : « \emph{Il} était là, \textbf{Meka}, immobile... »};
\end{tikzpicture}
\legende{Chaîne de référence : le référent « Meka » est repris par des substituts pronominaux (anaphore) et nominaux. La cataphore, plus rare, annonce le référent.}
\end{popfigure}

\begin{methode}[Identifier les substituts en quatre étapes]
\begin{enumerate}
\item \textbf{Entourer} le référent initial (la première mention du personnage ou de l'objet).
\item \textbf{Souligner} chaque pronom, nom ou périphrase qui le remplace dans les phrases suivantes.
\item \textbf{Relier} chaque substitut au référent par une flèche.
\item \textbf{Vérifier l'accord} : genre et nombre doivent correspondre (« il » = masculin singulier = Meka). Si l'accord est impossible, le pronom renvoie à un autre référent.
\end{enumerate}
\end{methode}

\begin{exemplebox}[Exemple chiffré]
Dans les dix premières pages du roman, on compte \textbf{14 occurrences} du référent « Meka » : \textbf{6 substituts pronominaux} (il, le, lui) et \textbf{8 substituts nominaux} (le vieux, le cultivateur, le vieillard, le récipiendaire...). Cette alternance évite la répétition et donne au portrait de Meka des éclairages successifs : tantôt son âge, tantôt son métier, tantôt son rôle dans la cérémonie.
\end{exemplebox}

\begin{attention}[Piège de l'ambiguïté référentielle]
Un pronom mal placé rend une phrase incompréhensible : « Meka rencontra le Commandant quand \emph{il} arriva » — qui arrive ? Dans vos dissertations, vérifiez toujours que chaque pronom ne peut renvoyer qu'à \emph{un seul} référent possible. En cas de doute, répétez le nom.
\end{attention}

\courssub{Outil de langue n°2 : sigles et emprunts}

Le roman d'Oyono est écrit en français, mais un français traversé par la réalité camerounaise et coloniale. Deux phénomènes lexicaux méritent l'attention.

\begin{definition}[Sigle et acronyme]
Un \cle{sigle} est formé des initiales de plusieurs mots et se prononce \textbf{lettre par lettre} : O-N-U (Organisation des Nations Unies), R-D-P-C. Un \cle{acronyme} est un sigle qui se lit \textbf{comme un mot ordinaire} : FIDES (Fonds d'investissement pour le développement économique et social), UNESCO. Un \cle{emprunt} est un mot pris à une autre langue : « mvet » (instrument et art oratoire), « ngondo » (fête traditionnelle douala), « boy », « commandant ».
\end{definition}

\begin{attention}[Erreur fréquente au Bac]
Ne confondez pas \textbf{sigle} (initiales prononcées lettre par lettre : O-N-U) et \textbf{acronyme} (initiales lues comme un mot : FIDES, UNESCO). Tous les acronymes sont des sigles, mais tous les sigles ne sont pas des acronymes.
\end{attention}

\textbf{Fonction des emprunts dans le roman.} Les mots empruntés aux langues locales (bassa, douala, ewondo) et à l'anglais ne sont pas des accidents : ils ont une double fonction. D'abord une \textbf{fonction réaliste} : ils ancrent le récit dans le Cameroun des années 1950, avec ses fêtes (le ngondo), sa musique (le mvet), ses réalités administratives (le FIDES, l'ONU qui supervise le territoire sous tutelle). Ensuite une \textbf{fonction identitaire} : en faisant entrer ces mots dans la langue française, Oyono affirme que la culture camerounaise a sa place dans la littérature, et il souligne le \emph{choc des cultures} au cœur du roman.

\begin{popfigure}
\begin{center}
\begin{tikzpicture}
\begin{axis}[ybar, bar width=0.55cm, width=0.9\linewidth, height=6.2cm, xlabel={Chapitres}, ylabel={Nombre d'emprunts}, symbolic x coords={Ch.1, Ch.2, Ch.3, Ch.4, Ch.5}, xtick=data, ymin=0, ymax=10, legend style={at={(0.5,-0.28)}, anchor=north, legend columns=3, font=\small}, nodes near coords, every node near coord/.append style={font=\scriptsize}]
\addplot[fill=popblue] coordinates {(Ch.1,2) (Ch.2,4) (Ch.3,3) (Ch.4,6) (Ch.5,5)};
\addplot[fill=poporange] coordinates {(Ch.1,3) (Ch.2,2) (Ch.3,5) (Ch.4,4) (Ch.5,2)};
\addplot[fill=popgreen] coordinates {(Ch.1,1) (Ch.2,3) (Ch.3,2) (Ch.4,5) (Ch.5,4)};
\legend{Anglais, Bassa, Douala}
\end{axis}
\end{tikzpicture}
\end{center}
\legende{Fréquence fictive des emprunts linguistiques par chapitre (données inventées pour l'exercice) : on observe une concentration des emprunts au chapitre 4, celui de la cérémonie de remise de la médaille à Douala, où le choc des cultures est maximal.}
\end{popfigure}

\courssub{Construire sa fiche de lecture synthétique}

La fiche de lecture est l'outil qui transforme une lecture passive en lecture active. Elle servira de réservoir d'exemples et de citations pour vos dissertations.

\begin{methode}[Les quatre rubriques d'une fiche de lecture efficace]
\begin{enumerate}
\item \textbf{Fiche d'identité} : titre, auteur, date, genre, cadre spatio-temporel, narrateur.
\item \textbf{Résumé par chapitres} : deux ou trois phrases par chapitre, au présent de narration, sans commentaire personnel.
\item \textbf{Personnages} : pour chacun, une ligne d'identité et sa fonction dans le récit.
\item \textbf{Thèmes majeurs} : relever les grandes problématiques avec une citation courte pour chacune.
\end{enumerate}
\end{methode}

\begin{exemplebox}[Les thèmes majeurs du \emph{Vieux Nègre et la Médaille}]
\begin{itemize}
\item \textbf{L'aliénation} : Meka, décoré par ceux qui ont pris ses terres, incarne l'homme coupé de lui-même par le système colonial.
\item \textbf{La dignité} : malgré l'humiliation, le vieil homme conserve une fierté que la médaille ne peut ni acheter ni effacer.
\item \textbf{Le choc des cultures} : la cérémonie de Douala met face à face traditions africaines et rituels coloniaux, dans une confrontation souvent comique, toujours cruelle.
\item \textbf{La satire de l'assimilation} : la médaille, symbole de la « reconnaissance » coloniale, est démasquée comme un leurre.
\end{itemize}
\end{exemplebox}

\begin{exoresolu}[Repérer référents et substituts]
\textbf{Énoncé.} Dans l'extrait suivant : « Meka entra dans la case. Le vieil homme s'assit près du feu. Il ferma les yeux. Kelara regarda son mari avec inquiétude : elle savait qu'il pensait à la médaille. » 1) Identifiez les deux référents. 2) Classez les substituts de « Meka » en pronoms et noms. 3) Justifiez que « elle » ne peut pas renvoyer à Meka.
\tcblower
\textbf{Solution détaillée.} 1) Les deux référents sont \textbf{Meka} (première mention : « Meka ») et \textbf{Kelara} (première mention : « Kelara »). 2) Substituts de Meka : un substitut \emph{nominal} (« le vieil homme », « son mari ») et des substituts \emph{pronominaux} (« il » dans « Il ferma les yeux » et « il pensait »). 3) « Elle » est un pronom féminin singulier ; or Meka est un référent masculin. L'accord en genre impose donc que « elle » renvoie à Kelara, seul référent féminin du passage. C'est la vérification de l'accord qui lève l'ambiguïté.
\end{exoresolu}

\begin{exercice}[À vous de jouer]
1) Relevez dans la quatrième de couverture de votre édition trois indices paratextuels et formulez pour chacun une hypothèse de lecture. 2) Dans un extrait de deux pages du roman, construisez la chaîne de référence de « la médaille » (référent et substituts). 3) Classez les mots suivants en sigle, acronyme ou emprunt : ONU, FIDES, ngondo, UNESCO, mvet.
\end{exercice}

\begin{corrige}[Exercice : sigles, acronymes, emprunts]
\textbf{Sigles} (lettre par lettre) : ONU. \textbf{Acronymes} (lus comme un mot) : FIDES, UNESCO. \textbf{Emprunts} (langues locales) : ngondo, mvet. Pour les questions 1 et 2, la correction se fera en classe à partir de votre édition : l'essentiel est de justifier chaque classement par la prononciation ou par la langue d'origine.
\end{corrige}

\begin{aretenir}[Bilan de la section]
Entrer dans une œuvre, c'est : \textbf{(1)} interroger les paratextes pour formuler des hypothèses ; \textbf{(2)} identifier genre, narrateur, personnages et cadre lors de la lecture ouvroir ; \textbf{(3)} suivre les chaînes de référence qui assurent la cohérence du récit ; \textbf{(4)} décoder sigles et emprunts qui ancrent le texte dans son époque ; \textbf{(5)} consigner le tout dans une fiche de lecture. Armande dispose désormais d'une lecture structurée de l'œuvre : la prochaine étape sera d'appliquer la même rigueur à l'analyse d'un sujet de dissertation et à la formulation d'une problématique.
\end{aretenir}

\coursec{La dissertation : De l'analyse du sujet à la problématique}

Dans la section précédente, nous sommes entrés dans l'œuvre de Ferdinand Oyono, \emph{Le Vieux nègre et la médaille}, et nous avons consolidé nos outils linguistiques (référents et substituts, sigles et emprunts). Ces acquis vont maintenant servir un objectif décisif pour le Baccalauréat technique : \cle{la dissertation}. Or, toute dissertation réussie commence bien avant la rédaction : elle commence par une \textbf{analyse rigoureuse du sujet} et par la construction d'une \textbf{problématique}. C'est l'objet de cette section.

\courssub{L'analyse du sujet : première étape obligatoire}

\textbf{Intuition.} Imagine un mécanicien à qui l'on apporte une machine en panne. S'il se met à démonter au hasard, il perdra son temps. Il commence par \emph{écouter} la machine, \emph{identifier} chaque pièce, \emph{comprendre} ce qu'on attend de lui. Devant un sujet de dissertation, l'élève est ce mécanicien : il ne faut jamais commencer à écrire avant d'avoir démonté le sujet pièce par pièce.

\begin{definition}[Le sujet de dissertation]
Le \cle{sujet} est l'énoncé (citation, question ou affirmation) proposé à la réflexion du candidat. Il contient toujours : des \textbf{mots-clés} (les termes porteurs du sens), une \textbf{relation} entre ces termes (opposition, cause, condition) et des \textbf{présupposés} (ce que le sujet admet sans le dire).
\end{definition}

\courssub{L'analyse terminologique : définir les mots-clés}

La première opération consiste à définir chaque mot-clé, d'abord \textbf{au dictionnaire}, puis \textbf{dans le contexte du sujet}. Les deux définitions ne coïncident pas toujours : c'est justement cet écart qui fait naître le problème.

\begin{exemplebox}[Analyse des termes du sujet « La tradition est-elle un frein au progrès ? »]
\begin{itemize}
\item \textbf{Tradition} : au dictionnaire, « transmission de doctrines, de coutumes, d'usages, de génération en génération ». Dans le sujet, elle désigne l'ensemble des pratiques culturelles héritées (rites, chefferies, langues, valeurs).
\item \textbf{Progrès} : au dictionnaire, « amélioration, avancement vers un état meilleur ». Dans le sujet, il vise le développement technique, économique et social (médecine moderne, école, industrie, démocratie).
\item \textbf{Frein} : au dictionnaire, « dispositif qui ralentit ou arrête un mouvement ». Par métaphore, dans le sujet : « obstacle, ce qui empêche d'avancer ». C'est le terme \emph{tensionnel} : toute la question tient dans ce mot.
\end{itemize}
\end{exemplebox}

\begin{attention}[Piège de la définition unique]
Ne jamais se contenter d'une seule définition. Un mot comme « progrès » peut être technique, moral, social ou politique. Si tu ne précises pas le sens retenu, ton devoir sera flou. Annonce toujours, dans ton brouillon, le sens que tu choisis : « Dans ce devoir, nous entendrons par progrès... ».
\end{attention}

\courssub{Les présupposés et l'opposition fondamentale}

\begin{definition}[Présupposé]
Un \cle{présupposé} est une idée que le sujet admet comme vraie sans la démontrer, et que le candidat doit repérer pour la discuter. Le repérer, c'est déjà commencer à problématiser.
\end{definition}

Dans le sujet « La tradition est-elle un frein au progrès ? », les présupposés sont :
\begin{enumerate}
\item la tradition et le progrès sont deux réalités \emph{distinctes}, voire opposées ;
\item la tradition pourrait avoir un effet \emph{négatif} (freiner) ;
\item le progrès est présenté comme quelque chose de \emph{désirable} qu'on pourrait empêcher.
\end{enumerate}

De ces présupposés naît l'\textbf{opposition fondamentale}, c'est-à-dire le couple thèse / antithèse :

\begin{center}
\begin{tabularx}{\linewidth}{|l|X|X|}
\hline
\rowcolor{popblueL} & \textbf{Thèse (oui)} & \textbf{Antithèse (non)} \\
\hline
Position & La tradition freine le progrès : rites figés, refus de l'école, poids du patriarcat. & La tradition n'est pas un frein : elle transmet des valeurs éthiques et peut s'adapter (chefferies modernes, Ngondo). \\
\hline
Fondement & Le passé immobile s'oppose au mouvement. & L'identité est la condition d'un progrès authentique. \\
\hline
\end{tabularx}
\end{center}

\courssub{La formulation de la problématique}

\begin{definition}[La problématique]
La \cle{problématique} est la \textbf{question tensionnelle} qui structure toute la dissertation. Elle ne se trouve pas \emph{dans} le sujet : elle se \emph{construit} à \emph{partir} du sujet, en reformulant l'opposition fondamentale sous forme d'une question ouverte qui appelle une réponse nuancée — ni un simple « oui », ni un simple « non ».
\end{definition}

\textbf{Intuition.} Une bonne problématique ressemble à une bonne question de débat entre amis : si la réponse est évidente, il n'y a rien à discuter ; si la question oblige à peser, à comparer, à nuancer, alors il y a matière à dissertation.

\begin{methode}[La technique du « Pourquoi ? / Comment ? / Dans quelle mesure ? »]
Pour tester la qualité d'une problématique, procède en trois étapes :
\begin{enumerate}
\item \textbf{Transforme} le sujet en question commençant par « Pourquoi ? », « Comment ? » ou « Dans quelle mesure ? ». Ces interrogations interdisent la réponse fermée (oui/non).
\item \textbf{Vérifie} que la question contient une \emph{tension} : deux forces qui s'opposent (tradition/progrès, argent/bonheur, liberté/règle).
\item \textbf{Contrôle} qu'elle appelle une démonstration en 2 ou 3 temps : si tu peux y répondre en une phrase, elle est mauvaise ; si elle exige d'examiner d'abord une thèse, puis son contraire, puis une synthèse, elle est bonne.
\end{enumerate}
\end{methode}

\begin{exemplebox}[Bonnes et mauvaises problématiques]
Sujet : « L'argent fait-il le bonheur ? »
\begin{itemize}
\item \textbf{Mauvaise problématique} : « L'argent est-il utile ? » — trop simple, la réponse est évidemment oui ; aucune tension, aucun débat possible.
\item \textbf{Bonne problématique} : « Dans quelle mesure la possession de l'argent conditionne-t-elle l'accès au bonheur, au risque de le corrompre ? » — question ouverte, tension entre \emph{conditionner} et \emph{corrompre}, réponse nécessairement nuancée en plusieurs temps.
\end{itemize}
\end{exemplebox}

\begin{attention}[Le piège du « plan catalogue »]
Une mauvaise problématique entraîne un mauvais plan. Le \textbf{plan catalogue} — I. La tradition, II. Le progrès, III. Leur rapport — ne démontre rien : il juxtapose des notions comme des articles sur une étagère. Exige-toi un \textbf{plan dialectique} (Thèse / Antithèse / Synthèse) ou un \textbf{plan analytique} (Causes / Conséquences / Solutions). Le plan détaillé sera étudié dans la section suivante ; retiens dès maintenant qu'il naît de la problématique.
\end{attention}

\begin{popfigure}
\begin{center}
\begin{tikzpicture}[node distance=0.55cm and 0.7cm,
etape/.style={rectangle, rounded corners, draw=popblue, fill=popblueL, text width=3.1cm, align=center, minimum height=0.9cm, font=\small},
fleche/.style={-{Stealth[length=2.5mm]}, thick, popdark}]
\node[etape, fill=popgoldL, draw=popgold, align=center] (sujet){\textbf{Sujet}\\ énoncé proposé};
\node[etape, right=of sujet, align=center] (termes){\textbf{Analyse des termes}\\ dictionnaire + contexte};
\node[etape, right=of termes, align=center] (presup){\textbf{Présupposés}\\ opposition thèse/antithèse};
\node[etape, below=of presup, align=center] (problema){\textbf{Problématique}\\ question ouverte et tensionnelle};
\node[etape, left=of problema, align=center] (brain){\textbf{Brainstorming}\\ carte heuristique, fiches Pour/Contre};
\node[etape, left=of brain, align=center] (tri){\textbf{Tri des arguments}\\ pertinence, originalité, exemples};
\node[etape, below=of tri, fill=popgreenL, draw=popgreen, align=center] (plan){\textbf{Plan détaillé}\\ dialectique ou analytique};
\draw[fleche] (sujet) -- (termes);
\draw[fleche] (termes) -- (presup);
\draw[fleche] (presup) -- (problema);
\draw[fleche] (problema) -- (brain);
\draw[fleche] (brain) -- (tri);
\draw[fleche] (tri) -- (plan);
\end{tikzpicture}
\end{center}
\legende{Diagramme de flux : les sept étapes qui conduisent du sujet au plan détaillé. Aucune étape ne doit être sautée.}
\end{popfigure}

\courssub{Les techniques de recherche d'idées}

Une fois la problématique posée, il faut trouver de la matière. Trois techniques complémentaires, à pratiquer au brouillon le jour de l'examen.

\textbf{1. Le brainstorming chronométré (5 minutes).} Pendant cinq minutes, sans t'arrêter et sans juger, note sur le brouillon \emph{tout} ce que la problématique évoque : mots, exemples, souvenirs de cours, faits d'actualité, proverbes, scènes de l'œuvre étudiée. La règle d'or : \textbf{quantité d'abord, qualité ensuite}. On ne censure rien pendant le chronométrage.

\textbf{2. La carte heuristique (mind-map).} On organise ensuite la matière autour de la problématique placée au centre : chaque branche principale correspond à un axe de réflexion (pour, contre, exemples), chaque sous-branche à une idée précise.

\begin{popfigure}
\begin{center}
\begin{tikzpicture}[mindmap, grow cyclic, every node/.style={concept, font=\small},
root concept/.append style={concept color=popblue, font=\small\bfseries, text=white},
level 1/.append style={level distance=3.4cm, sibling angle=120},
level 2/.append style={level distance=2.1cm, sibling angle=45}]
\node[root concept]{Tradition : frein au progrès ?}
  child[concept color=poporange] { node {Pour (frein)}
    child { node {Rites figés, refus du changement} }
    child { node {Patriarcat, poids des aînés} }
  }
  child[concept color=popgreen] { node {Contre (valeurs)}
    child { node {Valeurs éthiques, solidarité} }
    child { node {Capacité d'adaptation} }
  }
  child[concept color=popgold] { node {Exemples Cameroun}
    child { node {Ngondo, rites rénovés} }
    child { node {Chefferies modernes} }
  };
\end{tikzpicture}
\end{center}
\legende{Carte heuristique centrée sur la problématique « La tradition est-elle un frein au progrès ? » : trois branches (arguments pour, arguments contre, exemples camerounais) et leurs sous-branches.}
\end{popfigure}

\textbf{3. Les fiches « Pour / Contre ».} Alimentées pendant l'année par le cours, les lectures et l'actualité, ces fiches fournissent des arguments et des exemples prêts à l'emploi. Pour notre sujet, l'œuvre au programme offre une fiche précieuse : dans \emph{Le Vieux nègre et la médaille}, Meka reçoit une médaille des colonisateurs pour avoir donné ses terres et ses fils ; sa « récompense » révèle la destruction de l'homme traditionnel par un prétendu progrès colonial. Voilà un exemple littéraire directement mobilisable.

\begin{savaistu}[L'œil du correcteur au Baccalauréat technique]
Au Bac technique, le correcteur lit en priorité l'\textbf{introduction} et la \textbf{conclusion}. Une problématique absente ou floue dans l'introduction plafonne souvent la copie autour de 8/20, même si le développement est correct. Cinq minutes investies dans la problématique valent plus qu'une page de rédaction hâtive.
\end{savaistu}

\courssub{Sélection et hiérarchisation des arguments}

Le brainstorming produit trop d'idées : il faut trier. Trois critères permettent de retenir les meilleurs arguments.

\begin{propriete}[Critères de sélection d'un argument]
Un argument est retenu s'il satisfait trois critères :
\begin{enumerate}
\item la \textbf{pertinence} : il répond directement à la problématique, et non au sujet en général ;
\item l'\textbf{originalité} : il dépasse les évidences que tous les candidats donneront ;
\item la \textbf{disponibilité d'exemples} : il peut être illustré par l'œuvre au programme, l'actualité camerounaise ou la culture générale. Un argument sans exemple est une affirmation sans preuve.
\end{enumerate}
\end{propriete}

\textbf{Hiérarchiser}, c'est ensuite ordonner les arguments retenus : on va généralement du plus simple au plus fort, ou du plus attendu au plus personnel, afin que la démonstration gagne en puissance. Dans un plan dialectique, la thèse présente les arguments les plus immédiats, l'antithèse les renverse, la synthèse les dépasse.

\begin{center}
\begin{tabularx}{\linewidth}{|X|X|X|}
\hline
\rowcolor{popblueL} \textbf{Sujet type Bac technique} & \textbf{Problématique attendue} & \textbf{Problématique « hors-sujet » fréquente} \\
\hline
La tradition est-elle un frein au progrès ? & Dans quelle mesure la tradition, loin de s'opposer au progrès, peut-elle le conditionner ou l'entraver ? & Qu'est-ce que la tradition ? (définition sans tension) \\
\hline
L'argent fait-il le bonheur ? & Comment l'argent, condition du bien-être, peut-il en même temps corrompre le bonheur qu'il prétend servir ? & L'argent est-il utile ? (réponse évidente) \\
\hline
L'école est-elle la seule voie de la réussite ? & Dans quelle mesure la réussite dépend-elle de l'école, et quelles autres voies l'expérience camerounaise révèle-t-elle ? & Faut-il aller à l'école ? (question fermée) \\
\hline
\end{tabularx}
\end{center}

\begin{exoresolu}[Construire une problématique]
\textbf{Énoncé.} Sujet : « Le travail est-il une aliénation ? » 1) Définis les mots-clés. 2) Dégage les présupposés. 3) Formule une problématique de qualité.
\tcblower
\textbf{Solution détaillée.}
1) \emph{Travail} : activité humaine organisée visant la production de biens ou de services ; dans le sujet, il vise surtout le travail salarié. \emph{Aliénation} : perte de soi-même, état où l'individu devient étranger à ce qu'il fait et à ce qu'il est.
2) Présupposés : le travail pourrait déposséder l'homme de sa liberté ; le travail et l'épanouissement seraient opposés ; il existerait un « soi » que le travail menacerait.
3) Problématique : « Dans quelle mesure le travail, qui asservit l'homme à des tâches imposées, peut-il être en même temps la voie de son accomplissement ? » — question ouverte, tension entre \emph{asservir} et \emph{accomplir}, réponse exigeant une démonstration en trois temps (thèse : le travail aliène ; antithèse : le travail libère et socialise ; synthèse : tout dépend des conditions et du sens du travail).
\end{exoresolu}

\begin{exercice}[À toi de jouer]
Sujet : « Faut-il craindre les réseaux sociaux ? » 1) Définis les mots-clés au dictionnaire puis dans le contexte. 2) Relève deux présupposés. 3) Formule une problématique avec « Dans quelle mesure... ». 4) Propose deux arguments « pour » et deux arguments « contre », chacun accompagné d'un exemple (œuvre, actualité camerounaise ou culture générale).
\end{exercice}

\begin{corrige}[Exercice 1]
1) \emph{Réseaux sociaux} : plateformes numériques d'échange entre utilisateurs (WhatsApp, Facebook, TikTok) ; dans le sujet, ils désignent un espace social nouveau qui transforme les relations. \emph{Craindre} : éprouver de la peur devant un danger ; le terme présuppose une menace.
2) Présupposés : les réseaux sociaux présenteraient un danger ; on pourrait choisir de s'en méfier ou de s'en passer.
3) Problématique possible : « Dans quelle mesure les réseaux sociaux, outils de lien et d'information, constituent-ils une menace pour l'individu et la société ? »
4) Pour : ils diffusent rumeurs et violences (ex. : fausses informations en période électorale au Cameroun) ; ils enferment dans la dépendance (ex. : jeunes scotchés à TikTok au détriment des études). Contre : ils informent et mobilisent (ex. : campagnes de santé publique relayées sur WhatsApp) ; ils créent des activités économiques (ex. : jeunes entrepreneurs camerounais vendant via Facebook). Plan attendu : dialectique (thèse : un danger réel ; antithèse : un outil précieux ; synthèse : tout dépend de l'usage et de l'éducation aux médias).
\end{corrige}

\begin{aretenir}[L'essentiel de la section]
\begin{itemize}
\item On n'écrit jamais avant d'avoir \textbf{analysé les mots-clés} (dictionnaire + contexte) et dégagé les \textbf{présupposés}.
\item La \textbf{problématique} se construit : question ouverte, tensionnelle, qui appelle une réponse nuancée en 2 ou 3 temps.
\item La recherche d'idées combine \textbf{brainstorming chronométré}, \textbf{carte heuristique} et \textbf{fiches Pour/Contre}.
\item On trie les arguments selon trois critères : \textbf{pertinence}, \textbf{originalité}, \textbf{disponibilité d'exemples}.
\item Interdit : le plan catalogue. Exigé : un plan dialectique ou analytique — ce sera l'objet de la section suivante.
\end{itemize}
\end{aretenir}

\coursec{L'architecture maîtrisée : Élaboration du plan détaillé}

Dans la section précédente, nous avons appris à analyser un sujet de dissertation et à formuler une \cle{problématique} claire. Mais une question, même bien posée, ne fait pas encore une démonstration. Entre la problématique et la rédaction, il faut construire l'\cle{architecture} du devoir : c'est l'objet du \cle{plan détaillé}. Imaginez un maçon qui se lancerait dans la construction d'une maison sans plan d'architecte : les murs risqueraient de s'écrouler. Il en va de même de la dissertation au Probatoire et au Baccalauréat technique.

\begin{definition}[Le plan détaillé]
Le \cle{plan détaillé} est le \emph{squelette} de la dissertation. Il présente, avant toute rédaction, les grandes parties (I, II, III), les sous-parties (A, B, C) et, pour chacune, les arguments, les exemples et les transitions. Il doit être lisible comme un \emph{sommaire} : un correcteur doit comprendre toute la démonstration sans lire la rédaction.
\end{definition}

\courssub{Choisir le type de plan adapté à la problématique}

Il n'existe pas un plan unique, mais trois grands types de plans. Le bon plan est celui qui \emph{découle naturellement} de la problématique. C'est la problématique qui commande, jamais l'inverse.

\begin{definition}[Les trois types de plans]
\begin{itemize}
\item Le \cle{plan dialectique} : \emph{Thèse / Antithèse / Synthèse}. Il convient aux sujets qui opposent deux points de vue (débat, contradiction apparente). Exemple : « Faut-il rejeter la tradition au nom du progrès ? »
\item Le \cle{plan analytique} : \emph{Constat ou Causes / Manifestations / Conséquences ou Solutions}. Il convient aux sujets qui décrivent un phénomène à expliquer. Exemple : « Pourquoi les jeunes s'éloignent-ils des coutumes ? »
\item Le \cle{plan thématique} : \emph{Aspect 1 / Aspect 2 / Aspect 3}. Il convient aux sujets qui invitent à étudier une notion sous plusieurs angles. Exemple : « Quelles fonctions la tradition remplit-elle dans la société ? »
\end{itemize}
\end{definition}

\begin{methode}[Choisir son plan en trois questions]
\begin{enumerate}
\item Le sujet oppose-t-il deux thèses ? $\rightarrow$ plan \textbf{dialectique}.
\item Le sujet demande-t-il d'expliquer un fait social ? $\rightarrow$ plan \textbf{analytique}.
\item Le sujet demande-t-il de décrire les facettes d'une notion ? $\rightarrow$ plan \textbf{thématique}.
\end{enumerate}
\end{methode}

\begin{popfigure}
\begin{tikzpicture}[>=Stealth, thick, scale=1.05]
\node[draw=popblue, fill=popblueL, rounded corners, minimum width=3.4cm, minimum height=1.1cm, align=center] (these) at (0,0) {\textbf{THÈSE}\\La tradition protège};
\node[draw=poporange, fill=poporangeL, rounded corners, minimum width=3.4cm, minimum height=1.1cm, align=center] (antithese) at (7,0) {\textbf{ANTITHÈSE}\\La tradition enferme};
\node[draw=popgreen, fill=popgreenL, rounded corners, minimum width=3.6cm, minimum height=1.1cm, align=center] (synthese) at (3.5,3.4) {\textbf{SYNTHÈSE}\\Une tradition vivante,\\ouverte au progrès};
\draw[->, popdark, line width=1.2pt] (these) -- (antithese) node[midway, below, align=center, font=\small]{Négation de la thèse\\(critique)};
\draw[->, popdark, line width=1.2pt] (antithese) -- (synthese) node[midway, right, align=center, font=\small]{Dépassement};
\draw[->, popdark, line width=1.2pt] (these) -- (synthese) node[midway, left, align=center, font=\small]{Dépassement};
\end{tikzpicture}
\legende{Le triangle dialectique : la synthèse ne supprime pas la thèse et l'antithèse, elle les \emph{dépasse} en conservant ce que chacune a de juste.}
\end{popfigure}

\courssub{Rédiger les intitulés d'axes et les sous-parties}

Une fois le type de plan choisi, il faut rédiger les \cle{intitulés}. Les grandes parties (I, II, III) s'écrivent en \emph{phrases nominales courtes}, parallèles entre elles, qui annoncent l'argument fort de la partie. Les sous-parties (A, B, C) s'écrivent en \emph{phrases affirmatives} : chacune est une petite thèse partielle.

\begin{exemplebox}[Intitulés pour le sujet « Tradition et progrès »]
Plan dialectique :
\begin{itemize}
\item \textbf{I.} La tradition, fondement de l'identité et de la cohésion sociale.
\begin{itemize}
\item \textbf{A.} La tradition transmet les valeurs et structure la société.
\item \textbf{B.} Les rites et les langues locales maintiennent le lien communautaire.
\end{itemize}
\item \textbf{II.} Cependant, une tradition figée freine l'émancipation et le progrès.
\begin{itemize}
\item \textbf{A.} Certaines coutumes bloquent l'épanouissement individuel.
\item \textbf{B.} Le refus du changement isole la communauté.
\end{itemize}
\item \textbf{III.} Il faut donc bâtir une tradition vivante, réconciliée avec le progrès.
\begin{itemize}
\item \textbf{A.} Trier dans l'héritage : conserver l'essentiel, abandonner l'archaïque.
\item \textbf{B.} L'éducation comme lieu de rencontre entre mémoire et modernité.
\end{itemize}
\end{itemize}
Remarquez le parallélisme : chaque axe commence par un nom (« La tradition », « une tradition figée », « une tradition vivante »).
\end{exemplebox}

\begin{methode}[La règle des 3 x 3]
Pour un devoir équilibré, visez : \textbf{3 grandes parties} (I, II, III), \textbf{2 à 3 sous-parties} par grande partie (A, B, C), et \textbf{2 à 3 arguments illustrés} par sous-partie. Cette régularité garantit l'équilibre et la progressivité de la démonstration.
\end{methode}

\begin{attention}[Argument et exemple : ne pas confondre]
L'\cle{argument} est une raison \emph{abstraite} ; l'\cle{exemple} est un fait \emph{précis, daté ou cité}. Exemple : Argument = « La tradition structure la société » ; Exemple = « Le Ngondo chez les Sawa régule la vie de la cité ». Un argument sans exemple est une affirmation gratuite ; un exemple sans argument est une anecdote.
\end{attention}

\courssub{Remplir la grille du plan détaillé}

Le plan détaillé se présente sous forme d'une grille à quatre colonnes : l'axe, la sous-partie, l'argument accompagné de son exemple précis (citation de l'œuvre, fait social camerounais, référence historique), et la transition. C'est cette grille que vous construisez au brouillon pendant l'épreuve.

\begin{popfigure}
\begin{center}
\small
\begin{tabularx}{\linewidth}{|l|X|X|X|}
\hline
\rowcolor{popblueL}\textbf{Axes} & \textbf{Sous-parties} & \textbf{Arguments + Exemples} & \textbf{Transitions} \\
\hline
\textbf{I.} La tradition, fondement de l'identité & A. Elle transmet les valeurs & Arg. : la coutume donne des repères. Ex. : le respect des anciens dans \emph{Le Vieux Nègre et la médaille} (le vieillard honoré par son village). & « La tradition fonde, mais fonde-t-elle toujours justement ? » \\
\hline
 & B. Elle soude la communauté & Arg. : les rites créent le lien. Ex. : le Ngondo des Sawa ; les rites villageois chez Ferdinand Oyono. & « Ce lien peut cependant devenir une chaîne. » \\
\hline
\textbf{II.} Une tradition figée freine le progrès & A. Elle étouffe l'individu & Arg. : la coutume impose des conduites. Ex. : dans l'œuvre d'Oyono, le vieux nègre découvre que la médaille coloniale, symbole admis par tradition, est une humiliation. & « Dès lors, comment concilier mémoire et avenir ? » \\
\hline
 & B. Elle refuse le changement & Arg. : le rejet du nouveau isole. Ex. : refus de l'école dans certains villages ; méfiance envers les élites « modernes ». & « La solution n'est ni la soumission ni la rupture. » \\
\hline
\textbf{III.} Une tradition vivante, ouverte au progrès & A. Trier l'héritage & Arg. : tout héritage n'est pas à garder. Ex. : abandon progressif des pratiques contraires à la dignité, maintien des langues et des rites fédérateurs. & « Ce tri suppose un instrument : l'éducation. » \\
\hline
 & B. Éduquer pour réconcilier & Arg. : l'école transmet et renouvelle. Ex. : les littératures africaines (Oyono, Mongo Beti) font dialoguer mémoire et modernité. & Vers la conclusion : « Ainsi, tradition et progrès ne s'opposent que si l'on fige l'une ou l'autre. » \\
\hline
\end{tabularx}
\end{center}
\legende{Grille type du plan détaillé, pré-remplie pour le sujet « Tradition et progrès » à partir du \emph{Vieux Nègre et la médaille} de Ferdinand Oyono.}
\end{popfigure}

\begin{propriete}[La double fonction de la transition]
Toute \cle{transition} explicite fait double emploi : elle \emph{referme} l'idée précédente (résumé bref) et \emph{ouvre} la suivante (annonce + lien logique). Les connecteurs adaptés sont par exemple : \emph{Cependant}, \emph{Par ailleurs}, \emph{Dès lors}, \emph{Ainsi}, \emph{En effet}.
\end{propriete}

\courssub{Vérifier le plan : équilibre, progressivité, réponse à la problématique}

Avant de rédiger, relisez votre grille en vous posant trois questions de contrôle :

\begin{enumerate}
\item \textbf{Équilibre} : les trois parties ont-elles un volume comparable (même nombre de sous-parties et d'arguments) ? Une partie trop courte déséquilibre tout le devoir.
\item \textbf{Progressivité} : va-t-on du simple au complexe, du concret à l'abstrait ? Chaque partie doit marquer un progrès dans la réflexion, non une répétition.
\item \textbf{Réponse à la problématique} : la dernière partie répond-elle vraiment à la question posée dans l'introduction ? Si non, le plan est hors sujet.
\end{enumerate}

\begin{exemplebox}[Un plan détaillé réussi]
Au Baccalauréat 2023, une copie notée 18/20 présentait un plan détaillé de 3 axes, 9 sous-parties et 27 arguments illustrés, dont 4 citations exactes de l'œuvre au programme. Le correcteur a souligné la lisibilité du sommaire : la démonstration se comprenait avant même la lecture de la rédaction.
\end{exemplebox}

\begin{exoresolu}[Construire un axe complet]
\textbf{Énoncé.} Sujet : « L'école détruit-elle les valeurs traditionnelles ? » Rédigez l'axe I (thèse) d'un plan dialectique, avec deux sous-parties, un argument et un exemple pour chacune, et la transition vers l'axe II.
\tcblower
\textbf{Solution détaillée.}
\begin{itemize}
\item \textbf{I.} L'école peut sembler menacer les valeurs traditionnelles.
\item \textbf{A.} Elle détourne l'enfant de la coutume. \emph{Argument} : l'école impose une langue et des savoirs étrangers. \emph{Exemple} : dans \emph{Le Vieux Nègre et la médaille}, l'école et l'administration coloniales ont éloigné les villageois de leurs repères.
\item \textbf{B.} Elle crée une rupture entre générations. \emph{Argument} : l'élève instruit méprise parfois les anciens. \emph{Exemple} : les « évolués » des romans de Mongo Beti tournent le dos au village.
\item \textbf{Transition} : « Cependant, cette rupture n'est pas fatale : l'école peut aussi transmettre la tradition. » (elle résume l'axe I et annonce l'axe II).
\end{itemize}
\end{exoresolu}

\begin{popfigure}
\begin{center}
\begin{tikzpicture}
\begin{axis}[
    ybar, bar width=22pt,
    width=0.85\linewidth, height=6cm,
    ymin=0, ymax=140,
    ylabel={Durée (minutes)},
    symbolic x coords={Analyse, Plan, Rédaction, Relecture},
    xtick=data,
    nodes near coords,
    every node near coord/.append style={font=\small},
    axis line style={popdark},
    tick label style={font=\small},
]
\addplot[fill=popblue, draw=popdark] coordinates {(Analyse,15) (Plan,20) (Rédaction,120) (Relecture,15)};
\end{axis}
\end{tikzpicture}
\end{center}
\legende{Répartition idéale du temps d'examen pour une dissertation de 3 heures : analyse du sujet (15 min), élaboration du plan détaillé (20 min), rédaction (2 h), relecture (15 min).}
\end{popfigure}

\begin{attention}[Pièges fréquents]
\begin{itemize}
\item Rédiger les sous-parties en phrases nominales vagues (« A. Les traditions ») au lieu de phrases affirmatives qui énoncent une idée directrice.
\item Multiplier les exemples sans argument qui les commande.
\item Oublier les transitions : le correcteur doit sentir la \emph{progression} entre les parties.
\item Consacrer plus de 25 minutes au plan : c'est un outil, pas une fin en soi.
\end{itemize}
\end{attention}

\begin{savaistu}[L'annonce de plan dans l'introduction]
L'annonce de plan, dernière phrase de l'introduction, doit reprendre \emph{exactement} les libellés des axes (I, II, III) du plan détaillé, sans les sous-parties. Si votre axe I s'intitule « La tradition, fondement de l'identité », l'annonce dira : « Nous verrons d'abord que la tradition est le fondement de l'identité, puis... ». Cette cohérence entre introduction et plan est un critère de notation explicite au Baccalauréat technique.
\end{savaistu}

\begin{exercice}[À vous de jouer]
Sujet : « Faut-il préserver les langues locales à l'école ? »
\begin{enumerate}
\item Choisissez le type de plan adapté et justifiez votre choix.
\item Rédigez les trois intitulés d'axes en phrases nominales parallèles.
\item Développez l'axe II en deux sous-parties, chacune avec un argument, un exemple camerounais précis et une transition.
\item Vérifiez votre plan à l'aide des trois questions de contrôle (équilibre, progressivité, réponse à la problématique).
\end{enumerate}
\end{exercice}

\begin{corrige}[Exercice]
\begin{enumerate}
\item Le sujet oppose deux points de vue (préserver / ne pas préserver) : le plan \textbf{dialectique} s'impose.
\item Exemple d'axes : \textbf{I.} L'enseignement des langues locales semble freiner l'unité et l'efficacité scolaires. \textbf{II.} Pourtant, les langues locales portent l'identité et facilitent les apprentissages. \textbf{III.} Il convient donc d'articuler langues locales et langues officielles dans un bilinguisme maîtrisé.
\item Axe II : \textbf{A.} Les langues locales transmettent la culture. \emph{Argument} : la langue maternelle véhicule les valeurs et les récits. \emph{Exemple} : les contes et proverbes ewondo, duala ou fulfuldé transmis oralement. \emph{Transition} : « Par ailleurs, cette transmission n'est pas seulement culturelle : elle est aussi pédagogique. » \textbf{B.} Elles facilitent les apprentissages. \emph{Argument} : on apprend mieux dans une langue que l'on maîtrise. \emph{Exemple} : les écoles bilingues expérimentales du Cameroun montrent de meilleurs résultats en lecture. \emph{Transition} : « Dès lors, loin de s'opposer, langues locales et langues officielles peuvent se compléter. »
\item Vérification : les trois axes ont deux sous-parties chacun (équilibre) ; on va du préjugé défavorable à la solution (progressivité) ; l'axe III répond à la question « faut-il ? » (réponse à la problématique).
\end{enumerate}
\end{corrige}

\begin{aretenir}[L'essentiel]
Le plan détaillé est le squelette de la dissertation : il se choisit selon la problématique (dialectique, analytique ou thématique), se construit selon la règle des 3 x 3, associe à chaque argument un exemple précis, relie les parties par des transitions à double fonction, et se vérifie par l'équilibre, la progressivité et la réponse à la question posée. C'est lui qui, dans la section suivante, guidera la rédaction de l'introduction, des paragraphes et de la conclusion.
\end{aretenir}

\coursec{Approfondissement littéraire \& Cohérence textuelle (Lectures méthodiques)}

Dans la section précédente, nous avons appris à bâtir l'\cle{architecture} d'une dissertation : du sujet à la problématique, de la problématique au plan détaillé. Mais un plan, aussi solide soit-il, ne vaut que par la matière qu'il organise. Or cette matière, au Probatoire comme au Baccalauréat technique, c'est le texte qui la fournit. Il faut donc savoir \emph{lire méthodiquement} une œuvre : non pas la parcourir des yeux, mais la démonter pièce par pièce pour comprendre comment elle est fabriquée. C'est ce travail que nous allons accomplir sur \emph{Le Vieux Nègre et la médaille} de Ferdinand Oyono, tout en acquérant les outils linguistiques de la \cle{progression thématique}, de la \cle{cohésion} et de la \cle{cohérence}, qui serviront ensuite à rédiger vos propres paragraphes.

\begin{attention}[Lecture méthodique ou lecture cursive ?]
Ne confondez pas les deux exercices. La \cle{lecture cursive} est une lecture personnelle, rapide, faite à la maison : elle donne une vision globale de l'œuvre. La \cle{lecture méthodique} est une analyse guidée, lente, faite en classe avec le professeur : elle s'arrête sur un extrait, interroge la langue, les procédés, les enjeux. Les deux sont évaluées : une question de contrôle de lecture peut porter sur l'ensemble du roman (cursive), un commentaire ou une question d'analyse porte sur un passage précis (méthodique).
\end{attention}

\courssub{Lecture méthodique 1 : les chapitres 1 à 3 (la médaille, la fête, l'attente)}

\textbf{Observation de départ.} Relisons l'incipit : Meka, vieux paysan de Doum, apprend que les Blancs vont lui remettre une médaille pour « services rendus à la France » (il a donné ses deux fils à la guerre et sa terre à la mission). Premier constat : le narrateur ne nous dit jamais de l'extérieur ce que vaut cette médaille ; nous la découvrons \emph{à travers les yeux de Meka}.

\begin{definition}[Focalisation interne]
La \cle{focalisation interne} est un mode de narration où le lecteur perçoit les événements à travers la conscience d'un personnage : nous savons ce qu'il sait, nous ignorons ce qu'il ignore. Dans les chapitres 1 à 3, la focalisation est fixée sur \cle{Meka} : sa joie naïve, sa fierté, son attente fébrile deviennent les nôtres.
\end{definition}

\textbf{L'ironie narrative.} L'\cle{ironie} naît d'un écart entre ce que croit le personnage et ce que le lecteur comprend. Meka croit que la médaille est un honneur ; le lecteur, lui, mesure le prix réel : deux fils morts, une terre perdue. Plus le village félicite Meka, plus l'écart se creuse. La fête du chapitre 2, l'attente du chapitre 3 (les préparatifs, le voyage à Nkongsamba, les habits neufs) accumulent les signes d'une illusion collective. Oyono n'a pas besoin de commenter : c'est la \emph{mise en scène} qui dénonce.

\textbf{Le cadre colonial.} Relevons les indices : le commandant, le père Vandermayer, l'église qui domine le village, les « évolués » comme Engamba qui traduisent et servent d'intermédiaires. Le cadre est celui du Cameroun sous administration française : une hiérarchie où la parole vient d'en haut (les Blancs) et où les colonisés reçoivent, obéissent, espèrent. La médaille est le symbole parfait de cette relation : un objet brillant, sans valeur marchande, donné en échange de sacrifices immenses.

\begin{exemplebox}[Analyse d'un procédé : le style indirect libre]
Phrase type du chapitre 3 : Meka attend, et soudain la narration glisse : « Était-ce donc cela, la médaille ? » Pas de guillemets, pas de « il se demanda ». La pensée du personnage s'insère directement dans le récit au passé simple. C'est le \cle{style indirect libre} : la voix du narrateur et celle du personnage se confondent, et le lecteur entre dans la tête de Meka sans transition.
\end{exemplebox}

\begin{savaistu}[Pourquoi Oyono choisit-il le style indirect libre ?]
Ferdinand Oyono publie le roman en 1956, au moment où les nationalismes africains s'affirment. En fondant la pensée de Meka dans la narration, il évite le discours militant frontal : c'est le lecteur lui-même qui, placé dans la conscience du vieillard, découvre l'absurdité de la situation coloniale. La forme du récit est déjà une prise de position politique.
\end{savaistu}

\courssub{Lecture méthodique 2 : les chapitres 4 à 6 (la remise, le discours, le retour, la chute)}

\textbf{La bascule tragique.} Le chapitre 4 (la cérémonie de remise) est le sommet apparent de la gloire de Meka : fanfare, discours du Commandant, médaille épinglée. Mais dès le chapitre 5, tout bascule : le retour sous la pluie, la médaille qui devient un poids, l'errance nocturne de Meka, son arrestation et sa nuit au poste. Le chapitre 6 achève la chute : humilié, Meka jette la médaille. Le roman est construit comme une \cle{courbe dramatique} : montée de l'illusion (chapitres 1 à 4), rupture (chapitre 5), effondrement (chapitre 6).

\textbf{L'échec de la communication.} Toute la cérémonie repose sur des paroles qui ne communiquent pas : le Commandant parle français, Meka ne comprend pas ; Engamba traduit, mais déforme ; Meka répond, mais on ne l'écoute pas. La médaille elle-même est un message brouillé : pour les Blancs, elle signifie « la France vous remercie » ; pour Meka, elle finit par signifier « on s'est moqué de toi ». C'est ce qu'on appelle un \cle{malentendu} : le même signe reçoit deux interprétations opposées.

\textbf{La symbolique de la médaille jetée.} Jeter la médaille, c'est rejeter le contrat colonial tout entier : l'objet censé couronner une vie de soumission devient un déchet. Ce geste final transforme Meka : de vieillard crédule, il devient lucide. L'ironie du début se retourne en \emph{dénonciation} : le lecteur, qui souriait de la naïveté de Meka, partage désormais sa colère.

\begin{popfigure}
\begin{center}
\begin{tikzpicture}
\begin{axis}[
width=0.9\linewidth, height=6.2cm,
xlabel={Chapitres}, ylabel={Intensité (échelle 0 à 10)},
xmin=0.5, xmax=6.5, ymin=0, ymax=10.5,
xtick={1,2,3,4,5,6},
xticklabels={Ch.1, Ch.2, Ch.3, Ch.4, Ch.5, Ch.6},
legend style={at={(0.02,0.98)},anchor=north west,font=\small},
grid=both, grid style={popline!40},
]
\addplot[popcurve, color=popblue, mark=*, very thick] coordinates {(1,4) (2,6) (3,7) (4,10) (5,3) (6,1)};
\addlegendentry{Honneur / Fierté}
\addplot[popcurve, color=poporange, mark=square*, very thick, dashed] coordinates {(1,1) (2,1) (3,2) (4,2) (5,7) (6,10)};
\addlegendentry{Honte / Colère}
\end{axis}
\end{tikzpicture}
\legende{Évolution des indices « Honneur/Fierté » (bleu) et « Honte/Colère » (orange) au fil des six chapitres du roman. Le croisement des courbes entre les chapitres 4 et 5 matérialise la bascule tragique : la cérémonie de remise est le sommet de l'illusion, le retour sous la pluie et l'arrestation déclenchent l'effondrement. (Échelle indicative construite à partir du relevé des chaînes lexicales.)}
\end{center}
\end{popfigure}

\courssub{La progression thématique : comment le texte fait avancer son sujet}

\textbf{Intuition.} Quand vous racontez une journée, vous ne dites pas n'importe quoi dans n'importe quel ordre : chaque phrase reprend un élément de la précédente et ajoute une information nouvelle. C'est exactement ce que fait un texte bien construit.

\begin{definition}[La progression thématique]
La \cle{progression thématique} est l'organisation de l'information qui fait avancer le thème d'un texte. Dans chaque phrase, on distingue le \cle{thème} (ce dont on parle, l'information déjà connue) et le \cle{rhème} (ce qu'on en dit, l'information nouvelle). Le rhème d'une phrase peut devenir le thème de la suivante : Thème $\rightarrow$ Rhème $\rightarrow$ Nouveau Thème.
\end{definition}

\begin{exemplebox}[Application au roman]
« Meka [Thème] reçut une médaille [Rhème]. La médaille [Nouveau Thème] brillait au soleil [Rhème]. Ce soleil [Nouveau Thème] était celui de la grande fête [Rhème]. » Chaque information nouvelle devient le point de départ de la phrase suivante : le texte avance comme une chaîne. Rompez la chaîne (parlez de la médaille, puis soudain de la pluie, puis du Commandant sans lien) et le texte devient incohérent.
\end{exemplebox}

\textbf{Les chaînes lexicales et leur évolution.} Une \cle{chaîne lexicale} est l'ensemble des mots qui, dans un texte, renvoient au même champ de sens. Relevons dans le roman deux grandes chaînes antagonistes :

\begin{center}
\begin{tabularx}{\linewidth}{|l|X|X|}
\hline
\rowcolor{popblueL} \textbf{Chaîne lexicale} & \textbf{Chapitres 1 à 4 (montée)} & \textbf{Chapitres 5 à 6 (chute)} \\
\hline
Honneur / honte & médaille, honneur, gloire, fête, fierté, briller & honte, ridicule, boue, pluie, poste, jeter \\
\hline
Blanc / noir & les Blancs, le Commandant, le père, amis & les Blancs, rire, mépris, mensonge \\
\hline
Argent / terre & terre donnée, fils partis, sacrifice, récompense & terre perdue, fils morts, rien, pauvreté \\
\hline
\end{tabularx}
\end{center}

Le constat est décisif : le \emph{champ lexical de la fierté} domine la première moitié du roman, puis cède la place au \emph{champ lexical de l'humiliation}. Cette évolution n'est pas un décor : elle \emph{est} le sens de l'œuvre. Le roman raconte le passage d'un lexique à l'autre.

\begin{popfigure}
\begin{center}
\begin{tikzpicture}[scale=1]
% Colonisateur (gauche, bleu)
\node[font=\bfseries\small, popblue] at (-3.2,2.6) {Lexique du colonisateur};
\node[font=\Large\bfseries, popblue] at (-3.2,1.5) {honneur};
\node[font=\large, popblue!80] at (-4.6,0.7) {médaille};
\node[font=\normalsize, popblue!70] at (-2.0,0.8) {reconnaissance};
\node[font=\normalsize, popblue!60] at (-3.4,0.0) {France};
\node[font=\small, popblue!50] at (-4.8,-0.5) {amitié};
\node[font=\small, popblue!50] at (-1.9,-0.4) {gratitude};
\node[font=\footnotesize, popblue!40] at (-3.2,-1.0) {civilisation};
% Colonisé (droite, orange)
\node[font=\bfseries\small, poporange] at (3.2,2.6) {Lexique du colonisé};
\node[font=\Large\bfseries, poporange] at (3.2,1.5) {honte};
\node[font=\large, poporange!80] at (4.7,0.7) {pluie};
\node[font=\normalsize, poporange!70] at (1.9,0.8) {boue};
\node[font=\normalsize, poporange!60] at (3.5,0.0) {poste};
\node[font=\small, poporange!50] at (4.8,-0.5) {rire};
\node[font=\small, poporange!50] at (1.8,-0.4) {vieillard};
\node[font=\footnotesize, poporange!40] at (3.2,-1.0) {jetée};
\draw[popline, dashed] (0,-1.4) -- (0,2.8);
\end{tikzpicture}
\legende{Nuage de mots simplifié : fréquence relative des chaînes lexicales dans l'extrait de la cérémonie (chapitre 4) et dans le retour (chapitre 5). La taille des mots figure leur poids dans le passage : « honneur » et « médaille » structurent le discours officiel ; « honte », « pluie » et « boue » envahissent la réalité vécue par Meka.}
\end{center}
\end{popfigure}

\courssub{Cohésion et cohérence : la texture du texte}

\textbf{Intuition.} Un texte n'est pas une pile de phrases, c'est un tissu. La \cle{cohérence} est l'unité de sens globale (tout le texte parle de la même chose et va dans la même direction) ; la \cle{cohésion} est l'ensemble des liens visibles qui cousent les phrases entre elles.

\begin{definition}[Les outils de la cohésion]
La cohésion repose sur trois familles d'outils : (1) les \cle{référents et leurs substituts} (reprendre « Meka » par « il », « le vieillard », « le vieux nègre ») ; (2) les \cle{connecteurs logiques} (\emph{donc, mais, cependant, alors, en effet}) qui annoncent la relation entre les idées ; (3) la \cle{chaîne des temps verbaux} (passé simple et imparfait pour le récit, présent pour les vérités générales).
\end{definition}

\begin{methode}[Colorier la texture d'un texte]
Pour analyser la cohésion d'un extrait, procédez par couleurs :
\begin{enumerate}
\item \textbf{Bleu} : soulignez les référents et leurs substituts (Meka $\rightarrow$ il $\rightarrow$ le vieillard). Vérifiez qu'on peut toujours identifier de qui l'on parle.
\item \textbf{Vert} : entourez les connecteurs (\emph{mais, donc, cependant}) et demandez-vous quelle relation chacun introduit (opposition, conséquence, concession).
\item \textbf{Rouge} : surlignez les chaînes lexicales (tous les mots de l'honneur, tous les mots de la honte).
\item Observez le résultat : un texte cohésif présente un maillage régulier des trois couleurs ; un trou dans le maillage signale une rupture de cohésion.
\end{enumerate}
\end{methode}

\textbf{La structure énonciative.} Qui parle ? À qui ? Le roman alterne le discours direct (dialogues avec guillemets), le discours indirect (« le Commandant déclara que... ») et le discours indirect libre (pensée de Meka fondue dans la narration). Ce dernier procédé est capital : il installe le lecteur dans la conscience du héros et rend l'ironie possible.

\begin{attention}[Pièges fréquents en analyse]
\begin{itemize}
\item Ne confondez pas \emph{cohésion} (liens de surface : connecteurs, reprises) et \emph{cohérence} (unité de sens en profondeur). Un texte peut être cohésif sans être cohérent s'il enchaîne correctement des idées contradictoires.
\item N'affirmez pas « l'auteur utilise le style indirect libre » sans citer la phrase exacte et montrer l'absence de guillemets et de verbe de parole.
\item Un connecteur n'a pas de sens en soi : \emph{mais} n'exprime l'opposition que si les deux segments s'opposent réellement.
\end{itemize}
\end{attention}

\courssub{Le texte explicatif et argumentatif : les discours du Commandant et du Chef}

La cérémonie du chapitre 4 met en scène deux discours officiels. Un \cle{texte explicatif} vise à faire comprendre ; un \cle{texte argumentatif} vise à convaincre. Les discours officiels font les deux : ils expliquent la médaille et justifient la colonisation. Tout discours argumentatif s'analyse en trois étages : la \cle{thèse} (ce qu'on veut faire admettre), les \cle{arguments} (les raisons), les \cle{exemples} (les preuves concrètes).

\begin{popfigure}
\begin{center}
\begin{tikzpicture}[
grow=down,
level 1/.style={sibling distance=6.5cm, level distance=1.7cm},
level 2/.style={sibling distance=3.4cm, level distance=1.7cm},
every node/.style={draw, rounded corners, align=center, font=\small, fill=popcream},
edge from parent/.style={draw=popline, -{Stealth}, thick}
]
\node[fill=popblueL, font=\small\bfseries] {Thèse : « La France vous honore »}
  child { node[fill=popblue!15, align=center]{Argument 1 :\\ la médaille =\\ reconnaissance}
    child { node[fill=white, align=center]{Exemple :\\ les anciens\\ combattants} }
    child { node[fill=white, align=center]{Exemple :\\ Meka, « ami \\ de la France »} }
  }
  child { node[fill=popblue!15, align=center]{Argument 2 :\\ la France protège\\ et civilise}
    child { node[fill=white, align=center]{Exemple :\\ la mission,\\ l'église} }
    child { node[fill=white, align=center]{Exemple :\\ la paix\\ « apportée »} }
  };
\end{tikzpicture}
\legende{Arbre simplifié de la structure argumentative du discours du Commandant (chapitre 4) : une thèse générale, deux arguments, chacun étayé d'exemples. L'analyse révèle un discours parfaitement construit — mais au service d'une thèse que le roman dément.}
\end{center}
\end{popfigure}

\textbf{Le ton.} Le discours du Commandant est \cle{péremptoire} : affirmations sans nuance, vocabulaire de la certitude, aucune place pour la réponse de l'autre. Mais Oyono double ce ton d'une \emph{ironie} de narration : plus le Commandant est solennel, plus la scène est grotesque, car le lecteur sait ce que la médaille a coûté à Meka.

\begin{exemplebox}[Analyse énonciative : le « nous » qui masque le « je »]
Dans son discours, le Commandant multiplie les « nous » : « nous sommes réunis », « nous honorons ». Ce \cle{nous} se présente comme inclusif (la France et le Cameroun ensemble), mais il masque en réalité le \emph{je} du pouvoir colonial : c'est lui, le Commandant, qui décide, parle et distribue les honneurs. L'\cle{énonciateur} (celui qui parle réellement : l'autorité coloniale) se cache derrière un \cle{énonciataire} collectif imaginaire. Repérer qui dit « nous » et qui est exclu de ce « nous » est un outil d'analyse redoutable pour tout discours politique.
\end{exemplebox}

\begin{exoresolu}[Analyser la cohésion d'un court extrait]
Énoncé. Soit l'extrait reconstitué : « Meka reçut la médaille. Celle-ci brillait au soleil. Cependant, le vieillard ne souriait plus, car il comprenait enfin. » (a) Relevez les substituts du référent « la médaille » et du référent « Meka ». (b) Identifiez les connecteurs et précisez leur valeur. (c) Montrez comment la progression thématique fait avancer le texte.
\tcblower
Solution détaillée. (a) Substituts : « la médaille » est reprise par le pronom démonstratif « Celle-ci » ; « Meka » est repris par la périphrase « le vieillard » puis par le pronom « il ». La chaîne anaphorique est complète : aucune ambiguïté sur les référents. (b) Connecteurs : « Cependant » introduit une opposition (la brillance de l'objet contre la tristesse de l'homme) ; « car » introduit une explication (la cause du changement : la lucidité). (c) Progression thématique : « Meka » est le thème de la phrase 1, « la médaille » son rhème ; « la médaille » (reprise par « Celle-ci ») devient le thème de la phrase 2 ; la phrase 3 opère un retour au thème initial « Meka » (via « le vieillard ») avec un rhème nouveau : « il comprenait enfin ». Le texte boucle son mouvement : de l'homme à l'objet, puis de l'objet à l'homme transformé. C'est exactement la structure du roman entier, en miniature.
\end{exoresolu}

\begin{exercice}[À vous de jouer]
Dans le chapitre 5 (le retour sous la pluie), relevez : (1) cinq mots de la chaîne lexicale de l'humiliation ; (2) trois substituts du référent « Meka » ; (3) deux connecteurs et leur valeur ; (4) une phrase au style indirect libre, en justifiant votre identification. Rédigez ensuite un paragraphe de dix lignes montrant comment la cohésion de l'extrait sert la bascule tragique du roman.
\end{exercice}

\begin{aretenir}[L'essentiel de la section]
\begin{itemize}
\item La lecture méthodique démonte un extrait : focalisation, ironie, cadre, énonciation.
\item Le roman est construit sur une bascule : chapitres 1 à 4 (illusion, champ lexical de la fierté), chapitres 5 à 6 (lucidité, champ lexical de l'humiliation).
\item La progression thématique enchaîne thème et rhème ; les chaînes lexicales dessinent le sens global.
\item La cohésion s'analyse en trois couleurs : référents (bleu), connecteurs (vert), chaînes lexicales (rouge).
\item Tout discours argumentatif se décompose en thèse, arguments, exemples ; le « nous » officiel peut masquer un « je » de pouvoir.
\end{itemize}
\end{aretenir}

Ces outils d'analyse ne sont pas une fin en soi : ils préparent la section suivante, où nous passerons \emph{de la plume à la copie} — rédiger l'introduction, insérer les citations, construire les transitions — en appliquant à votre propre production ce que nous venons de découvrir dans celle d'Oyono.

\coursec{De la plume à la copie : Rédaction, Intégration \& Remédiation}

Dans la section précédente, les lectures méthodiques du \emph{Vieux Nègre et la médaille} de Ferdinand Oyono nous ont permis de vérifier la cohérence d'un texte et d'affiner notre compréhension de l'œuvre. Nous disposons désormais d'un plan détaillé solide : sujet analysé, problématique formulée, idées hiérarchisées. Mais un plan, même excellent, ne vaut rien s'il reste à l'état de squelette. Il faut maintenant lui donner chair : rédiger. Cette dernière section vous conduit de la plume à la copie : rédaction de l'introduction, construction des paragraphes, articulation des transitions, rédaction de la conclusion, puis intégration finale et remédiation. C'est ici que se joue la note au Probatoire et au Baccalauréat technique.

\courssub{La rédaction de l'introduction : l'entonnoir}

\begin{definition}[Introduction et texte argumentatif]
L'\cle{introduction} est la première partie rédigée de la dissertation. Elle conduit le lecteur du général vers le particulier, jusqu'à la question que le devoir va traiter. Rappelons la distinction essentielle : le texte \cle{explicatif} répond à « Comment ? Pourquoi ? » de manière neutre (il informe) ; le texte \cle{argumentatif} répond à « Doit-on ? Faut-il ? » par une prise de position justifiée. La dissertation est un texte \textbf{argumentatif} : vous défendez une thèse, vous ne vous contentez pas d'expliquer.
\end{definition}

Imaginez un \cle{entonnoir} : large en haut, étroit en bas. L'introduction fonctionne exactement ainsi. On part du large (le monde, le siècle, une vérité générale), on resserre (l'œuvre, le sujet précis), et on aboutit à la pointe (la problématique), avant d'annoncer l'itinéraire (le plan).

\begin{methode}[Les quatre étapes obligatoires de l'introduction]
\begin{enumerate}
\item \textbf{L'accroche (amorce contextuelle)} : une phrase générale liée au thème du sujet (une vérité sur l'homme, la société, la colonisation, la tradition...). Jamais de définition de dictionnaire recopiée.
\item \textbf{La présentation du sujet et de l'œuvre} : on resserre. On cite l'auteur, le titre, le genre, le contexte de l'œuvre, et on reformule le sujet avec ses propres mots.
\item \textbf{La problématique} : une question (ou un ensemble de questions liées) qui naît naturellement de ce qui précède. C'est le fil conducteur de tout le devoir.
\item \textbf{L'annonce du plan} : on présente les deux (ou trois) parties du développement, sans détailler les paragraphes.
\end{enumerate}
\end{methode}

\begin{popfigure}
\begin{center}
\begin{tikzpicture}[font=\small]
\draw[fill=popblueL, draw=popblue, rounded corners] (0,0) rectangle (10,1.4);
\node[align=center] at (5,0.7) {\textbf{1. Accroche} : vérité générale (le monde, l'homme, l'époque)};
\draw[fill=popgreenL, draw=popgreen, rounded corners] (1,1.7) rectangle (9,3.1);
\node[align=center] at (5,2.4) {\textbf{2. Sujet / œuvre} : auteur, titre, reformulation du sujet};
\draw[fill=poporangeL, draw=poporange, rounded corners] (2,3.4) rectangle (8,4.8);
\node[align=center] at (5,4.1) {\textbf{3. Problématique} : la question directrice};
\draw[fill=poppinkL, draw=poppink, rounded corners] (3,5.1) rectangle (7,6.5);
\node[align=center] at (5,5.8) {\textbf{4. Annonce du plan} : I... II...};
\draw[-{Stealth}, very thick, popdark] (11.2,0.7) -- (11.2,5.8) node[midway, right, align=center]{du large\\vers\\l'étroit};
\end{tikzpicture}
\end{center}
\legende{Modèle de l'« introduction parfaite » en entonnoir : chaque étape resserre le propos, de l'accroche générale (bleu) jusqu'à l'annonce du plan (rose).}
\end{popfigure}

\begin{exemplebox}[Introduction réussie (sujet : « La médaille récompense-t-elle le mérite ? »)]
« Toute société distingue ses membres par des honneurs censés sanctionner leur valeur. [\textbf{Accroche}] Dans \emph{Le Vieux Nègre et la médaille}, roman publié en 1956, Ferdinand Oyono raconte l'histoire de Meka, vieillard camerounais décoré par l'administration coloniale pour sa « fidélité ». [\textbf{Sujet/œuvre}] Mais cette récompense est-elle la reconnaissance d'un mérite réel, ou l'instrument d'une domination déguisée ? [\textbf{Problématique}] Nous verrons d'abord que la médaille apparaît comme une récompense flatteuse, puis qu'elle se révèle être un piège aliénant. [\textbf{Plan}] »
\end{exemplebox}

\begin{attention}[Les fautes qui tuent une introduction]
\begin{itemize}
\item Commencer par « De tout temps, l'homme... » sans lien avec le sujet : accroche creuse.
\item Raconter toute l'intrigue de l'œuvre : l'introduction n'est pas un résumé.
\item Oublier l'annonce du plan : le correcteur ne sait pas où vous allez.
\item Poser une problématique sans rapport avec le sujet posé.
\end{itemize}
\end{attention}

\begin{exoresolu}[Réécrire une introduction faible]
\textbf{Énoncé.} Voici une introduction faible : « La médaille est un objet en métal. Dans le roman, Meka reçoit une médaille. Nous allons parler de ce roman. » Relevez trois défauts, puis réécrivez cette introduction en quatre étapes.
\tcblower
\textbf{Solution.} Défauts : (1) accroche-définition sans portée ; (2) aucune présentation de l'auteur ni du contexte ; (3) pas de problématique ni d'annonce de plan. Réécriture : « Les honneurs publics visent, en principe, à couronner le mérite d'un individu. [\textbf{Accroche}] C'est cette logique que Ferdinand Oyono met à l'épreuve dans \emph{Le Vieux Nègre et la médaille} (1956), où Meka, vieux paysan camerounais, est décoré par l'administration coloniale. [\textbf{Sujet/œuvre}] Cette décoration honore-t-elle réellement le vieillard, ou sert-elle les intérêts du colonisateur ? [\textbf{Problématique}] Après avoir montré que la médaille flatte et réconforte Meka, nous démontrerons qu'elle le coupe en réalité de sa communauté et de sa dignité. [\textbf{Plan}] »
\end{exoresolu}

\courssub{Le développement : la méthode A.E.C. et l'insertion des citations}

Chaque paragraphe du développement est une mini-dissertation. Il suit la méthode \cle{A.E.C.} : \textbf{A}ffirmer, \textbf{E}xpliquer-Étayer, \textbf{C}onclure (et transiter). Un paragraphe = une idée = un argument développé.

\begin{methode}[Construire un paragraphe avec A.E.C.]
\begin{enumerate}
\item \textbf{Affirmer} : la première phrase énonce clairement l'argument (phrase-thèse du paragraphe).
\item \textbf{Expliquer et étayer} : on développe l'idée, puis on la prouve par une citation de l'œuvre, un exemple culturel ou un fait d'actualité camerounaise.
\item \textbf{Conclure et transiter} : une phrase de mini-bilan qui confirme l'argument et prépare le paragraphe suivant.
\end{enumerate}
\end{methode}

\begin{popfigure}
\begin{center}
\begin{tikzpicture}[font=\small, node distance=0.35cm]
\node[draw=popblue, fill=popblueL, rounded corners, minimum width=9.5cm, minimum height=0.9cm] (t) {\textbf{Thèse du paragraphe} : j'affirme l'argument};
\node[draw=popgreen, fill=popgreenL, rounded corners, minimum width=9.5cm, minimum height=0.9cm, below=of t] (a1) {\textbf{Argument 1 + citation} : je prouve par le texte (guillemets + analyse)};
\node[draw=poporange, fill=poporangeL, rounded corners, minimum width=9.5cm, minimum height=0.9cm, below=of a1] (a2) {\textbf{Argument 2 + exemple} : j'illustre (réalité camerounaise, autre œuvre)};
\node[draw=poppink, fill=poppinkL, rounded corners, minimum width=9.5cm, minimum height=0.9cm, below=of a2] (b) {\textbf{Mini-bilan + transition} : je conclus et j'annonce la suite};
\draw[-{Stealth}, thick, popdark] (t) -- (a1);
\draw[-{Stealth}, thick, popdark] (a1) -- (a2);
\draw[-{Stealth}, thick, popdark] (a2) -- (b);
\end{tikzpicture}
\end{center}
\legende{Schéma du paragraphe type « A.E.C. » : de l'affirmation au mini-bilan, en passant par la preuve textuelle et l'exemple concret.}
\end{popfigure}

\begin{attention}[La citation « orpheline »]
Une citation \cle{orpheline} est une citation jetée dans la copie sans analyse : « Meka donne ses terres. « Il avait donné ses terres à la mission. » Cela montre sa générosité. » C'est insuffisant. Règle d'or : \textbf{1 citation = minimum 2 lignes de commentaire}. Il faut (a) expliquer le sens de la citation et (b) la relier à la thèse du paragraphe. La citation s'introduit (avec deux-points ou un verbe de parole), se met entre guillemets, puis se commente.
\end{attention}

\begin{exemplebox}[Citation correctement intégrée]
La générosité de Meka tourne à la dépossession. Oyono écrit : « Il avait donné ses terres à la mission ». Ce don, présenté comme un acte de foi, révèle en vérité l'emprise de l'administration coloniale sur les esprits : en cédant sa terre, source de vie et d'identité pour un paysan, Meka aliène son autonomie. La médaille ne fera que couronner cette dépossession déjà accomplie.
\end{exemplebox}

\courssub{Les transitions : les coutures du devoir}

Une dissertation sans transitions est un empilement de blocs ; avec des transitions, c'est un raisonnement. Il existe deux niveaux de transitions : entre paragraphes d'une même partie (transition simple : « De plus », « Cependant », une phrase de relance) et entre les grandes parties (transition forte).

\begin{propriete}[La transition forte entre deux parties]
La transition entre la partie I et la partie II doit obligatoirement : (1) \textbf{reprendre le mot-clé} de la partie qui se termine (synthèse partielle) et (2) \textbf{annoncer le mot-clé} de la partie qui commence (relance problématique). Modèle : « Si [mot-clé de I]..., il n'en reste pas moins que [mot-clé de II]... » Exemple : « Si la tradition freine parfois l'individu (I), elle peut aussi impulser son épanouissement (II). »
\end{propriete}

\begin{exemplebox}[Transition forte réussie]
« Ainsi, la médaille apparaît d'abord à Meka comme une consécration : elle le comble de fierté et lui vaut l'admiration des siens. Mais cette lueur est trompeuse. Derrière l'éclat du ruban se cache en réalité un instrument de domination, comme la suite du roman va le révéler. » (Reprise du mot-clé « consécration » de la partie I, annonce du mot-clé « domination » de la partie II.)
\end{exemplebox}

\courssub{La conclusion : l'entonnoir inversé}

Si l'introduction est un entonnoir (du large vers la pointe), la conclusion est un \cle{entonnoir inversé} : on part de la pointe (la réponse précise à la problématique) et on s'élargit vers le large (l'ouverture).

\begin{methode}[Rédiger la conclusion]
\begin{enumerate}
\item \textbf{Le bilan synthétique} : répondez clairement à la problématique en reformulant les acquis des deux parties, sans copier-coller les phrases du développement. Interdit : introduire une idée neuve.
\item \textbf{L'ouverture pertinente} : élargissez vers une autre œuvre, un fait d'actualité, une question philosophique plus vaste. L'ouverture doit rester en lien logique avec le sujet.
\end{enumerate}
\end{methode}

\begin{exemplebox}[Ouverture réussie]
« Au terme de cette étude, Meka incarne le refus de l'aliénation : derrière l'honneur de la médaille se dissimulait l'asservissement. Ce combat trouve un écho aujourd'hui dans la revendication des langues nationales à l'école au Cameroun : affirmer sa culture, c'est déjà résister. »
\end{exemplebox}

\begin{attention}[Pièges de la conclusion]
\begin{itemize}
\item L'ouverture hors sujet (« parlons maintenant de la guerre dans le monde ») : elle doit prolonger le sujet, pas le quitter.
\item Le bilan qui introduit un argument nouveau : tout argument neuf doit être dans le développement.
\item La conclusion d'une seule ligne : elle vaut presque zéro point.
\end{itemize}
\end{attention}

\courssub{Activité d'intégration finale et remédiation}

\begin{exercice}[Sujet inédit type Bac — production autonome]
Sujet : « L'école est-elle le seul chemin de la réussite ? » En vous appuyant sur votre expérience, sur \emph{Le Vieux Nègre et la médaille} et sur la réalité camerounaise, produisez : (1) un plan détaillé en deux parties (titres + sous-parties + exemples) ; (2) l'introduction complète en quatre étapes ; (3) un paragraphe type construit selon A.E.C. avec une citation analysée ; (4) la conclusion avec ouverture. Durée : 1 h 30.
\end{exercice}

\begin{corrige}[Exercice d'intégration — éléments attendus]
Plan possible : I. L'école, voie privilégiée de la réussite (1. elle transmet savoirs et diplômes — ex. : accès aux emplois publics au Cameroun ; 2. elle forme l'esprit critique — ex. : Meka, sans instruction, se laisse abuser par le discours colonial). II. Mais d'autres chemins mènent à la réussite (1. la formation professionnelle et l'artisanat — ex. : les techniciens formés dans les lycées techniques ; 2. l'expérience et l'entrepreneuriat — ex. : commerçants et agriculteurs prospères). Introduction : accroche sur la réussite comme aspiration universelle, présentation du débat, problématique (« L'instruction scolaire est-elle la condition unique de l'accomplissement de l'individu ? »), annonce du plan. Paragraphe A.E.C. attendu : affirmation, citation d'Oyono analysée en deux lignes minimum, exemple camerounais, mini-bilan. Conclusion : bilan nuancé (l'école est un chemin majeur mais non exclusif) + ouverture (la place grandissante des formations techniques et professionnelles au Cameroun).
\end{corrige}

\begin{methode}[L'auto-évaluation croisée]
Échangez votre copie avec un camarade. Chacun note la copie de l'autre sur les six critères de la grille ci-dessous, puis justifie chaque note par une observation écrite. Comparez ensuite avec la grille du professeur : l'écart entre votre notation et la sienne est un indicateur précieux de votre lucidité d'apprenti correcteur.
\end{methode}

\begin{popfigure}
\begin{center}
\begin{tikzpicture}[scale=0.85, font=\small]
\foreach \r in {1,2,3,4} \draw[popline] (0,0) circle (\r);
\foreach \a/\t in {90/{Analyse/Problématique}, 30/{Plan/Structure}, -30/{Contenu/Exemples}, -90/{Langue/Style}, -150/{Intro/Conclusion}, 150/{Orthographe}}
  \node[align=center] at (\a:5.1) {\t};
\foreach \a in {90,30,-30,-90,-150,150} \draw[popline] (0,0) -- (\a:4);
\draw[thick, popblue, fill=popblue!25] (90:3.4) -- (30:3.0) -- (-30:2.6) -- (-90:2.2) -- (-150:3.2) -- (150:3.6) -- cycle;
\draw[thick, poporange, dashed] (90:4) -- (30:4) -- (-30:4) -- (-90:4) -- (-150:4) -- (150:4) -- cycle;
\node[popblue] at (0,-4.9) {Copie évaluée (cercle 4 = 20/20)};
\node[poporange] at (0,-5.5) {Objectif : 20/20 (pointillés)};
\end{tikzpicture}
\end{center}
\legende{Grille radar d'auto-évaluation : six axes notés sur 20 (le cercle de rayon 4 vaut 20/20, chaque cercle intérieur retire 5 points). Reportez vos notes et visualisez vos points forts et vos axes de remédiation.}
\end{popfigure}

\begin{savaistu}[La remédiation n'est pas une punition]
La \cle{remédiation} est l'étape où l'on réécrit \emph{un seul} paragraphe noté faible, en appliquant exactement la correction du professeur. C'est là que la note monte : les élèves qui progressent le plus vite ne sont pas ceux qui refont toute la dissertation, mais ceux qui corrigent chirurgicalement leur point faible — une introduction refaite, une citation enfin analysée, une transition enfin construite. Une copie corrigée et retravaillée vaut dix copies jamais relues.
\end{savaistu}

\begin{aretenir}[L'essentiel de la section]
\begin{itemize}
\item Introduction = entonnoir en 4 étapes : accroche, sujet/œuvre, problématique, annonce du plan.
\item Paragraphe = méthode A.E.C. : Affirmer, Expliquer-Étayer, Conclure-Transiter.
\item 1 citation = minimum 2 lignes d'analyse (sens + lien avec la thèse). Jamais de citation orpheline.
\item Transition forte = mot-clé de la partie finie + mot-clé de la partie annoncée.
\item Conclusion = entonnoir inversé : bilan qui répond à la problématique + ouverture pertinente.
\item La remédiation (réécriture ciblée) est le moteur réel de la progression.
\end{itemize}
\end{aretenir}

Vous voilà désormais équipé de la chaîne complète : analyser, problématiser, planifier, rédiger, corriger. Au Baccalauréat technique, le correcteur ne demande pas la perfection : il demande une démarche visible, maîtrisée et honnête. C'est exactement celle que vous venez de construire.

\newpage

\coursec{Activité d'intégration}

\coursec{Produire un plan détaillé d'une dissertation}

\courssub{Objectifs de l'activité}

À l'issue de cet atelier « Architecte du Bac », tu seras capable de :

\begin{itemize}
\item analyser un sujet de dissertation en identifiant les \cle{mots-clés}, les \cle{présupposés} et en formulant une \cle{problématique} pertinente ;
\item rechercher et organiser des idées, des arguments et des exemples précis tirés de l'œuvre étudiée, \emph{Le Vieux Nègre et la Médaille} de Ferdinand Oyono ;
\item élaborer un \cle{plan détaillé} complet (axes, sous-parties, arguments, citations, transitions) ;
\item rédiger une introduction, un développement et une conclusion cohérents, en insérant correctement des citations et en assurant la \cle{cohésion} et la \cle{cohérence} du texte grâce aux substituts du référent et à la progression thématique ;
\item évaluer ta propre production et celle d'un camarade à l'aide d'une grille critériée, dans un esprit de remédiation.
\end{itemize}

Cette activité mobilise l'ensemble des savoirs de l'unité : analyse du sujet, recherche des idées, élaboration du plan détaillé, rédaction de l'introduction, des paragraphes avec citations et transitions, rédaction de la conclusion, ainsi que les outils de langue (référent et substituts, sigles et emprunts, progression thématique, texte explicatif).

\courssub{Situation}

Tu es élève en classe de Terminale technique au Lycée technique de Nkongsamba. À trois semaines du Probatoire, ton professeur de français organise, un samedi matin, un \textbf{examen blanc} de trois heures dans la salle polyvalente. L'épreuve reproduit les conditions réelles : sujet inédit, copie double, aucun document autorisé sauf les extraits distribués, barème officiel sur 20 points.

Le sujet tombé ce jour-là est le suivant :

\begin{definition}[Sujet de l'examen blanc]
\textbf{Sujet :} Le personnage de Meka dans \emph{Le Vieux Nègre et la Médaille} incarne-t-il la résistance ou la soumission ?
\end{definition}

Pour t'aider, le professeur te remet trois documents :

\begin{itemize}
\item \textbf{Document 1 :} extrait du chapitre 1 — Meka apprend qu'il va recevoir une médaille ; il s'agenouille, remercie le commandant, offre sa terre aux missionnaires.
\item \textbf{Document 2 :} extrait du chapitre 3 — la veillée au village : les anciens, Engamba, Amalia et les femmes commentent avec ironie cette médaille ; Meka, grisé par les éloges et les boissons, se laisse porter par l'orgueil.
\item \textbf{Document 3 :} extrait du chapitre 6 — après la cérémonie, Meka, oublié sous la pluie, humilié, arrêté puis libéré, rentre au village et déclare à ses frères que la médaille n'est rien, que « le Blanc est le Blanc » et que les siens sont sa seule richesse.
\end{itemize}

Le barème indicatif est affiché au tableau : analyse du sujet et problématique (4 points), plan détaillé (4 points), introduction et conclusion (4 points), développement avec citations (6 points), langue et présentation (2 points).

\courssub{Consignes}

\textbf{Étape 1 — Analyse du sujet (20 minutes).}
\begin{enumerate}
\item Souligne les mots-clés du sujet et dégage leur champ sémantique (résistance / soumission).
\item Repère le présupposé contenu dans l'alternative « résistance ou soumission » : peut-on l'accepter, la nuancer, la dépasser ?
\item Formule une problématique sous forme de question ouverte qui dépasse le simple « oui ou non ».
\item Sur ta fiche de brainstorming, note toutes les idées, scènes et citations qui te viennent à l'esprit, sans les trier encore.
\end{enumerate}

\textbf{Étape 2 — Plan détaillé (30 minutes).}
\begin{enumerate}
\setcounter{enumi}{4}
\item Regroupe tes idées en deux ou trois axes et choisis le type de plan adapté (thématique, dialectique ou analytique) en le justifiant.
\item Remplis la grille officielle : pour chaque partie, indique l'axe, deux sous-parties, les arguments, une citation précise (chapitre) et la transition annoncée.
\end{enumerate}

\textbf{Étape 3 — Rédaction (1 heure 30).}
\begin{enumerate}
\setcounter{enumi}{6}
\item Rédige l'introduction complète (amorce, présentation de l'œuvre et de l'auteur, problématique, annonce du plan).
\item Rédige au moins deux parties du développement, chaque paragraphe comportant une idée, un argument, une citation introduite et commentée.
\item Rédige la conclusion (bilan, réponse à la problématique, ouverture).
\end{enumerate}

\textbf{Étape 4 — Auto-correction guidée (20 minutes).}
\begin{enumerate}
\setcounter{enumi}{9}
\item Relis-toi avec la grille critériée : cohérence entre plan et rédaction, qualité des citations, correction de la langue, problématique réellement traitée.
\end{enumerate}

\textbf{Étape 5 — Échange en binôme (20 minutes).}
\begin{enumerate}
\setcounter{enumi}{10}
\item Échange ta copie avec ton voisin : note une force et un axe de progrès, avec bienveillance et précision.
\end{enumerate}

\courssub{Ressources à mobiliser}

\begin{aretenir}[Boîte à outils de la dissertation]
\begin{itemize}
\item \textbf{Analyse du sujet :} mots-clés, champ sémantique, présupposés, problématique.
\item \textbf{Types de plans :} thématique (aspects d'une question), dialectique (thèse / antithèse / synthèse), analytique (constat / causes / conséquences).
\item \textbf{Introduction :} amorce — présentation de l'œuvre — problématique — annonce du plan.
\item \textbf{Paragraphe :} idée — argument — citation introduite (« Comme le montre Oyono au chapitre \ldots ») — commentaire.
\item \textbf{Conclusion :} bilan — réponse nuancée — ouverture.
\item \textbf{Langue :} reprise du référent par des substituts (Meka, le vieillard, le héros, il), connecteurs logiques (d'abord, cependant, enfin), progression thématique assurant la cohérence.
\item \textbf{Sigles et emprunts :} identifier les sigles (ex. : ONU, CEMAC) et emprunts lexicaux (ex. : \og{}chef de poste\fg{}, \og{}commandant\fg{}) dans les extraits pour en analyser la portée réaliste ou ironique.
\item \textbf{Texte explicatif :} structurer l'explication de la citation (contexte, sens littéral, sens figuré, portée argumentative).
\end{itemize}
\end{aretenir}

\courssub{Démarche et corrigé commenté}

\begin{exoresolu}[Corrigé commenté de l'examen blanc]
\textbf{Énoncé.} Le personnage de Meka dans \emph{Le Vieux Nègre et la Médaille} incarne-t-il la résistance ou la soumission ? Produire l'analyse du sujet, le plan détaillé et les éléments de rédaction attendus.

\tcblower

\textbf{Étape 1 — Analyse du sujet.}

Les mots-clés sont « Meka », « incarne », « résistance », « soumission ». Le champ sémantique de la soumission : obéissance, agenouillement, gratitude, naïveté, aliénation. Celui de la résistance : refus, révolte, lucidité, retour aux siens, dignité retrouvée. Le présupposé du sujet est que Meka serait \emph{soit} l'un \emph{soit} l'autre : c'est une fausse alternative, car le personnage évolue au fil du roman. Problématique possible : \emph{Dans quelle mesure le parcours de Meka illustre-t-il le passage d'une soumission naïve à une prise de conscience qui s'apparente à une résistance lucide ?}

\textbf{Étape 2 — Plan détaillé (plan dialectique justifié : le sujet pose une alternative que le parcours du personnage dépasse par une synthèse).}

\begin{center}
\begin{tabularx}{\linewidth}{|l|X|X|}
\hline
\rowcolor{popblueL} \textbf{Partie} & \textbf{Arguments} & \textbf{Citations (chapitres)} \\
\hline
I. Une soumission apparente et aliénée & 1. Meka accepte tout des Blancs (terre, religion, médaille) par naïveté. 2. Il se laisse griser par l'orgueil et les honneurs factices. & Chap. 1 : l'agenouillement, le don de la terre aux missionnaires. Chap. 3 : la veillée, son ivresse d'orgueil face aux éloges. \\
\hline
II. Une prise de conscience douloureuse & 1. La cérémonie révèle le mépris colonial (attente sous la pluie). 2. L'humiliation et l'arrestation ouvrent les yeux de Meka sur la réalité du pouvoir. & Chap. 6 : l'attente interminable sous la pluie, la nuit au poste de police, l'indifférence du commandant. \\
\hline
III. Une résistance lucide et identitaire & 1. Meka rejette la médaille et les valeurs coloniales. 2. Il se réfugie dans la communauté villageoise, seule richesse véritable. & Chap. 6 : « le Blanc est le Blanc » ; déclaration à ses frères : « mes frères sont ma seule richesse ». \\
\hline
\end{tabularx}
\end{center}

Transitions annoncées :
\begin{itemize}
\item Entre I et II : « Pourtant, cette soumission n'est que le point de départ d'un parcours : l'épreuve de la cérémonie va transformer le vieillard en lui arrachant ses illusions. »
\item Entre II et III : « Fort de cette lucidité nouvelle, acquise dans la souffrance, Meka opère un choix radical : il tourne le dos aux leurres coloniaux pour retrouver l'essentiel, sa communauté. »
\end{itemize}

\textbf{Étape 3 — Éléments de rédaction.}

\emph{Introduction (modèle).} Amorce : la littérature africaine francophone interroge souvent le rapport complexe entre colonisé et colonisateur, oscillant entre aliénation et révolte. Présentation : Ferdinand Oyono, écrivain et diplomate camerounais, publie \emph{Le Vieux Nègre et la Médaille} en 1956 ; le roman raconte l'histoire de Meka, vieillard de Doum, « récompensé » par l'administration coloniale pour sa docilité. Problématique : Meka incarne-t-il la résistance ou la soumission ? Annonce du plan : après avoir montré sa soumission apparente (I), nous analyserons sa prise de conscience douloureuse (II), pour enfin mettre en lumière sa résistance lucide et identitaire (III).

\emph{Paragraphe rédigé (exemple, partie I).} D'abord, Meka apparaît comme l'image même de la soumission consentie. Il accueille la médaille comme une grâce divine et ne remet en cause ni l'autorité du commandant, ni la présence des missionnaires. Comme le montre Oyono au chapitre 1, le vieillard s'agenouille devant le commandant et offre même sa terre aux missionnaires : « Meka s'agenouilla, baisa les pieds du commandant » et « donna sa terre aux missionnaires ». Ce geste, loin d'être anodin, symbolise l'aliénation d'un homme qui a intériorisé la domination coloniale au point de donner spontanément ce que le colonisateur aurait pu prendre par la force.

\emph{Paragraphe rédigé (exemple, partie III).} Enfin, Meka accomplit un acte de résistance silencieuse mais radicale. Devant ses frères, il dénonce la vanité de la médaille et affirme son attachement aux valeurs traditionnelles. Au chapitre 6, il lance : « Le Blanc est le Blanc [...] Mes frères sont ma seule richesse. » Cette réplique marque le retour à l'authenticité : le héros refuse l'assimilation et choisit la solidarité villageoise comme unique rempart contre l'oppression.

\emph{Conclusion (modèle).} En définitive, Meka n'est ni un soumis définitif, ni un héros de la révolte armée : son parcours va de la soumission naïve à la résistance intérieure et lucide. La médaille, symbole de la mystification coloniale, finit rejetée comme un objet sans valeur. On peut se demander si cette résistance par le retrait et le retour aux sources, partagée par de nombreux personnages de la littérature africaine des indépendances (comme dans \emph{L'Aventure ambiguë} de Cheikh Hamidou Kane), n'est pas, à sa manière, plus durable que la confrontation ouverte.

\textbf{Étape 4 — Auto-correction.} Vérifier : la problématique est-elle traitée jusqu'au bout ? Chaque citation est-elle introduite (contexte) et commentée (portée) ? Les substituts de « Meka » (le vieillard, le héros, le récipiendaire, il) évitent-ils les répétitions ? Les connecteurs (d'abord, cependant, enfin) assurent-ils la progression logique ? L'orthographe des noms propres (Oyono, Doum, Meka, Engamba) est-elle respectée ?

\textbf{Étape 5 — Échange.} Exemple de retour bienveillant : « Force : problématique claire, évolution du personnage bien montrée, citations bien choisies et commentées. Axe de progrès : la transition entre les parties II et III pourrait être plus explicite sur le lien entre lucidité et acte de résistance ; attention à l'accord des participes passés (ex. : "la médaille qu'il a reçue"). »
\end{exoresolu}

\courssub{Bilan de l'activité}

Cet atelier a permis de mettre en évidence trois acquis essentiels pour l'examen. D'abord, la réussite d'une dissertation repose sur une \cle{analyse du sujet} rigoureuse : dépasser la fausse alternative « résistance ou soumission » en identifiant la dynamique temporelle du personnage était la clé du sujet. Ensuite, le \cle{plan détaillé} apparaît comme l'architecture invisible de la copie : une copie cohérente est une copie où chaque paragraphe découle logiquement du plan annoncé. Enfin, la \cle{relecture critériée} et l'échange entre pairs montrent que l'amélioration d'un texte est un processus itératif : identifier une force et un axe de progrès précis prépare efficacement à la remédiation avant l'examen réel. Tu es désormais prêt à affronter l'épreuve de dissertation au Probatoire et au Baccalauréat technique avec méthode et sérénité.

\newpage

\coursec{Méthodes et exercices résolus}

\courssub{Méthodes et exercices résolus}

\begin{methode}[Analyser un sujet de dissertation et formuler la problématique]
    \textbf{Objectif :} Comprendre l'exigence du sujet, en délimiter les termes et poser la question centrale qui guidera toute la réflexion.
    \begin{enumerate}
        \item \textbf{Dégager les mots-clés :} Souligner les termes importants (notions, verbes d'action, connecteurs). Les définir (sens propre, sens figuré, sens technique).
        \item \textbf{Identifier le type de sujet :} Sujet \cle{thèse} (affirmation à discuter), sujet \cle{question} (interrogation directe), sujet \cle{citation} (propos d'auteur à commenter), sujet \cle{corpus} (ensemble de textes à confronter).
        \item \textbf{Analyser les relations logiques :} Repérer les oppositions, concessions, causes, conséquences, conditions implicites dans l'énoncé.
        \item \textbf{Formuler la problématique :} Transformer le sujet en une \cle{question ouverte} (commençant par \emph{En quoi}, \emph{Comment}, \emph{Dans quelle mesure}, \emph{Pourquoi}) qui contient une \cle{tension} (un paradoxe, un conflit de valeurs, une complexité). Elle ne doit pas être une simple reformulation du sujet.
        \item \textbf{Vérifier la pertinence :} La problématique doit permettre une réponse nuancée (ni tout à fait oui, ni tout à fait non) et englober l'ensemble du sujet.
    \end{enumerate}
\end{methode}

\begin{exoresolu}[Analyse de sujet et problématique — Type Bac Probatoire]
\textbf{Sujet :} \og La littérature africaine n'a de sens que si elle est engagée. \fg{} Discutez cette affirmation.
\tcblower
\textbf{Solution détaillée :}

\textbf{1. Dégagement des mots-clés :}
\begin{itemize}
    \item \textbf{Littérature africaine :} Ensemble des œuvres littéraires produites par des Africains ou sur l'Afrique (roman, théâtre, poésie, oralité). Notion large : diversité des époques (négritude, post-indépendance, contemporain), des langues (français, anglais, langues nationales), des thèmes.
    \item \textbf{Sens :} Signification, finalité, portée, utilité, valeur esthétique ou éthique.
    \item \textbf{Engagée :} Au sens sartrien : littérature qui prend parti, qui dénonce, qui participe à la lutte politique/sociale (colonialisme, néocolonialisme, dictatures, condition féminine, écologie). \textbf{Attention :} L'engagement n'est pas seulement politique ; il peut être esthétique (défense de la langue), culturel (réhabilitation de l'histoire), humain.
    \item \textbf{N'a de sens que si :} Restriction exclusive (condition nécessaire). La phrase affirme : \og Sans engagement $\Rightarrow$ pas de sens \fg. C'est une thèse forte, totalisante. La particule restrictive \og que \fg (correlatif de \og ne \fg) est l'opérateur logique clé.
\end{itemize}

\textbf{2. Type de sujet :} Sujet \textbf{thèse} (affirmation péremptoire à discuter). Verbe opérateur : \og Discutez \fg $\Rightarrow$ Examiner les pour et les contre, nuancer, hiérarchiser.

\textbf{3. Relations logiques / Tensions :}
\begin{itemize}
    \item Tension entre \textbf{utilitarisme} (littérature comme arme/outil) et \textbf{esthétisme} (littérature comme art pour l'art, recherche de la beauté formelle).
    \item Tension entre \textbf{devoir de mémoire/témoignage} (engagement historique) et \textbf{liberté créatrice} (refus de l'assignation à résidence identitaire ou politique).
    \item L'adverbe \og seulement \fg (sous-entendu dans la restriction \og n'a de sens \textbf{que} si \fg) est le point de rupture : il nie toute valeur à la littérature non engagée (divertissement, exploration psychologique, formalisme, poésie pure).
\end{itemize}

\textbf{4. Formulation de la problématique :}
\og \textbf{Dans quelle mesure l'engagement (politique, social, culturel) constitue-t-il l'unique horizon de sens de la littérature africaine, au risque de nier sa dimension esthétique universelle et la liberté de l'écrivain ?} \fg

\textbf{Justification :}
\begin{itemize}
    \item \emph{Dans quelle mesure} : ouvre la nuance (partiellement oui, mais pas uniquement).
    \item \emph{Engagement... horizon de sens} : reprend le terme du sujet en l'élargissant.
    \item \emph{Au risque de nier...} : introduit la contre-partie (la tension), annonçant le plan dialectique (Thèse : l'engagement comme moteur historique ; Antithèse : les limites de l'engagement exclusif / valeur de l'esthétique ; Synthèse : dépassement, littérature monde, engagement esthétique).
\end{itemize}

\textbf{Conclusion :} La problématique est validée car elle permet de structurer un plan en trois temps dialectiques tout en respectant l'exigence du \og Discutez \fg.
\end{exoresolu}

\begin{methode}[Rechercher, sélectionner et hiérarchiser les idées (Inventio)]
    \textbf{Objectif :} Passer de la problématique à la matière brute de la dissertation (arguments, exemples, citations) sans écrire de phrases définitives.
    \begin{enumerate}
        \item \textbf{Brainstorming / Mind-mapping :} Noter \textbf{tout} ce qui vient à l'esprit (œuvres lues en classe : \emph{Le Vieux Nègre et la Médaille}, \emph{Une Si Longue Lettre}, \emph{Les Bouts de Bois de Dieu}, \emph{Vol de Nuit}, auteurs : Senghor, Césaire, Oyono, Sembène, Bâ, Kourouma, etc., connaissances personnelles, actualité) autour des deux pôles de la problématique (Engagement vs Esthétique/Liberté).
        \item \textbf{Classification par pôles (Thèse / Antithèse / Synthèse) :} Ranger chaque idée dans la colonne correspondante.
            \begin{itemize}
                \item \textbf{Thèse (Pour l'engagement) :} Nécessité historique (décolonisation), fonction sociale (éveil des consciences), littérature de combat (Ousmane, Beti, Oyono).
                \item \textbf{Antithèse (Limites de l'engagement exclusif) :} Risque de propagande / pamphlet, mort de l'art (Senghor : \og L'émotion est nègre, la raison hellène \fg mais l'art est universel), droit à l'intime, à l'imaginaire (romans psychologiques, fantasy africaine, poésie pure).
                \item \textbf{Synthèse (Dépassement) :} L'engagement \emph{est} une esthétique (style, langue, forme). \og Littérature-monde \fg (Glissant, Mabanckou). L'écrivain engagé \emph{par la forme}.
            \end{itemize}
        \item \textbf{Sélection et hiérarchisation :} Ne garder que les arguments \textbf{forts}, \textbf{variés} (différents genres, époques, aires géographiques) et \textbf{illustrables} par des citations précises. Éliminer le hors-sujet, les redites, les généralités non étayées.
        \item \textbf{Ordonnancement logique :} À l'intérieur de chaque partie, ordonner les arguments : du plus général au plus spécifique, ou du chronologique au thématique, ou de l'argument d'autorité (citation) à l'analyse personnelle.
    \end{enumerate}
\end{methode}

\begin{exoresolu}[Recherche d'idées — Application au sujet précédent]
\textbf{Consigne :} À partir de la problématique dégagée (\og Dans quelle mesure l'engagement constitue-t-il l'unique horizon de sens... \fg), établir la liste structurée des arguments et exemples (œuvres, citations) pour les trois temps de la dissertation.
\tcblower
\textbf{Solution détaillée :}

\textbf{Partie I : L'engagement, fondement historique et fonction sociale majeure de la littérature africaine (Thèse)}

\begin{center}
\begin{tabularx}{\linewidth}{|l|X|X|}\hline
\textbf{Argument} & \textbf{Explication / Analyse} & \textbf{Exemples / Citations (Précis)} \\ \hline
1. Réponse à l'urgence historique (décolonisation) & La littérature naissante (années 30-60) est indissociable de la lutte pour l'indépendance. Écrire, c'est exister politiquement. & \textbf{Senghor / Césaire / Damas (Négritude)} : \emph{Anthologie de la nouvelle poésie nègre et de Madagascar}. \textbf{Césaire} : \og Ma négritude n'est pas une tour d'ivoire \fg (\emph{Cahier d'un retour au pays natal}). \\ \hline
2. Dénonciation du colonialisme et du néocolonialisme & Fonction démythifiante : montrer la réalité de l'oppression, réhabiliter l'histoire. & \textbf{Mongo Beti} : \emph{Le Pauvre Christ de Bomba}, \emph{Main basse sur le Cameroun}. \textbf{Ferdinand Oyono} : \emph{Le Vieux Nègre et la Médaille} (ironie, dévoilement de la médaille comme leurre). \\ \hline
3. Éveil des consciences et éducation civique & L'écrivain comme \og griot moderne \fg ou \og éclaireur \fg (Sartre : \og L'écrivain doit être le révélateur du monde \fg). & \textbf{Sembène Ousmane} : \emph{Les Bouts de bois de Dieu} (grève, solidarité collective). \textbf{Ken Bugul} : \emph{Le Baobab fou} (condition féminine). \\ \hline
\end{tabularx}
\end{center}

\textbf{Partie II : Les limites de l'engagement exclusif : risque de propagande, négation de l'esthétique et de la liberté créatrice (Antithèse)}

\begin{center}
\begin{tabularx}{\linewidth}{|l|X|X|}\hline
\textbf{Argument} & \textbf{Explication / Analyse} & \textbf{Exemples / Citations (Précis)} \\ \hline
1. Risque de pamphlet / littérature à thèse & Si le message prime sur la forme, l'œuvre devient un tract, datée, sans portée universelle. Perte de la valeur artistique. & \textbf{Sartre} (\emph{Qu'est-ce que la littérature ?}) : distinction engagement / propagande. \textbf{Wole Soyinka} : \og Un tigre ne proclame pas sa tigritude \fg (critique de la négritude essentialiste). \\ \hline
2. Droit à l'intime, à l'imaginaire, à l'universel & L'écrivain africain n'est pas un porte-parole permanent. Il explore la psyché, le rêve, le fantastique, l'absurde. & \textbf{Amos Tutuola} : \emph{L'Ivrogne dans la brousse} (imaginaire yoruba). \textbf{Alain Mabanckou} : \emph{Verre cassé} (oralité, humour, bar, pas de thèse politique explicite). \textbf{Véronique Tadjo} : \emph{Reine Pokou} (mythe, réécriture poétique). \\ \hline
3. Refus de l'assignation identitaire & \og Littérature africaine \fg comme catégorie marketing ou cage. Revendication d'une littérature \og monde \fg. & \textbf{Édouard Glissant} : \og Tout-monde \fg, \og Relation \fg. \textbf{Achille Mbembe} : \og Critique de la raison nègre \fg. \\ \hline
\end{tabularx}
\end{center}

\textbf{Partie III : Vers une réconciliation : l'engagement par la forme, l'esthétique comme résistance (Synthèse)}

\begin{center}
\begin{tabularx}{\linewidth}{|l|X|X|}\hline
\textbf{Argument} & \textbf{Explication / Analyse} & \textbf{Exemples / Citations (Précis)} \\ \hline
1. L'engagement \emph{est} une poétique & Le choix de la langue (français \og décolonisé \fg, wolof, créole), de la structure (oralité, rythme), du style \emph{est} un acte politique. & \textbf{Ahmadou Kourouma} : \emph{Les Soleils des Indépendances} (français malinké, syntaxe africaine = résistance linguistique). \textbf{Ken Bugul} : écriture fragmentée = corps morcelé. \\ \hline
2. L'universel passe par le singulier & La grande littérature engagée transcende son contexte par sa force artistique (Shakespeare, Tolstoï, Hugo). & \textbf{Chinua Achebe} : \emph{Things Fall Apart} / \emph{Le Monde s'effondre} (tragédie grecque en pays igbo = universel). \\ \hline
3. L'écrivain citoyen : double allégeance (art + cité) & L'engagement ne s'arrête pas à la publication ; l'écrivain intervient dans la cité (prix Nobel, tribunes), mais son œuvre littéraire reste une œuvre d'art. & \textbf{Wole Soyinka} (Nobel, prison, engagement politique direct). \textbf{Abdourahman Waberi} : engagement pour les migrants \emph{par} la poésie/roman. \\ \hline
\end{tabularx}
\end{center}

\textbf{Conclusion de l'exercice :} Cette matière permet de bâtir un plan détaillé équilibré (3 parties, 3 arguments chacune, exemples variés : poésie, roman, théâtre, essai ; aires : Afrique francophone, anglophone, lusophone, diaspora).
\end{exoresolu}

\begin{methode}[Élaborer un plan détaillé de dissertation (Dispositio)]
    \textbf{Objectif :} Construire l'architecture rigoureuse de la copie : Introduction, Développement (3 parties, 3 sous-parties chacune), Conclusion. Le plan détaillé s'écrit \textbf{au brouillon}, phrases nominales ou infinitives, avec \textbf{citations exactes} et \textbf{transitions} préparées.
    \begin{enumerate}
        \item \textbf{Introduction (schéma obligatoire) :}
            \begin{enumerate}
                \item \textbf{Amorce (Accroche) :} Citation d'auteur, fait d'actualité, définition étymologique, référence culturelle large. \textit{Liée au sujet.}
                \item \textbf{Sujet posé :} Reformulation du sujet du bac (synonymes, explication des termes).
                \item \textbf{Sujet divisé / Débat :} Présentation des deux thèses opposées (Pour / Contre).
                \item \textbf{Problématique :} La question unique formulée à l'étape précédente.
                \item \textbf{Annonce du plan :} \og Nous montrerons d'abord que... (I). Nous verrons ensuite que... (II). Enfin, nous démontrerons que... (III). \fg
            \end{enumerate}
        \item \textbf{Développement : Structure ternaire (Dialectique classique).}
            \begin{itemize}
                \item \textbf{Partie I (Thèse) :} Arguments principaux soutenant l'affirmation du sujet. Titre de partie : phrase nominale affirmative.
                \item \textbf{Partie II (Antithèse) :} Arguments limitant, contredisant, complexifiant la thèse. Titre de partie : phrase nominale nuancée ou opposée.
                \item \textbf{Partie III (Synthèse / Dépassement) :} Réconciliation, niveau supérieur de compréhension. Titre de partie : phrase nominale intégrative.
            \end{itemize}
            \textbf{Chaque sous-partie (A, B, C) = 1 paragraphe = 1 idée force + 1 argumentation + 1 exemple précis (citation/référence) + 1 mini-transition vers le suivant.}
        \item \textbf{Transitions inter-parties :} Rédiger \textbf{au brouillon} la phrase de transition explicite (rappel de l'idée close + annonce de l'idée qui vient + lien logique : \og Cependant \fg, \og Par ailleurs \fg, \og Dès lors \fg).
        \item \textbf{Conclusion (schéma obligatoire) :}
            \begin{enumerate}
                \item \textbf{Bilan (Récapitulatif synthétique) :} Réponse directe à la problématique en reprenant les grands axes (I, II, III).
                \item \textbf{Ouverture :} Prolongement pertinent (autre œuvre, autre art, actualité, question philosophique nouvelle), sans ouvrir un nouveau débat non traité.
            \end{enumerate}
        \item \textbf{Vérification finale :} Cohérence (le plan répond-il à la problématique ?), Équilibre (longueur ~égale des parties), Progression (logique dialectique respectée).
    \end{enumerate}
\end{methode}

\begin{exoresolu}[Élaboration du plan détaillé complet — Sujet : \og La littérature africaine n'a de sens que si elle est engagée \fg]
\textbf{Consigne :} Rédiger le plan détaillé intégral (Introduction, Développement avec sous-parties, citations, transitions) au format brouillon.
\tcblower
\textbf{Solution détaillée :}

\textbf{INTRODUCTION}
\begin{itemize}
    \item \textbf{Amorce :} Citation de \textbf{Aimé Césaire} (\emph{Discours sur le colonialisme}) : \og Une civilisation qui use de ses principes pour ruser, pour tricher, pour étouffer, est une civilisation moribonde. \fg $\rightarrow$ L'écrivain africain ne peut rester neutre face à l'histoire.
    \item \textbf{Sujet posé :} L'affirmation \og La littérature africaine n'a de sens que si elle est engagée \fg postule que la valeur (le sens) de cette littérature est conditionnée par sa fonction militante, politique ou sociale.
    \item \textbf{Sujet divisé :} D'un côté, l'histoire (négritude, décolonisation) semble valider ce lien organique (littérature = arme). De l'autre, la diversité contemporaine (esthétique, intime, universalité) conteste cette assignation réductrice.
    \item \textbf{Problématique :} \og Dans quelle mesure l'engagement constitue-t-il l'unique horizon de sens de la littérature africaine, au risque de nier sa dimension esthétique universelle et la liberté de l'écrivain ? \fg
    \item \textbf{Annonce du plan :} Nous analyserons d'abord comment l'engagement fonde historiquement le sens de cette littérature (I). Nous examinerons ensuite les périls d'un engagement exclusif qui nie l'art et la liberté (II). Enfin, nous verrons comment les grandes œuvres réconcilient engagement et exigence esthétique dans une \og littérature-monde \fg (III).
\end{itemize}

\textbf{DÉVELOPPEMENT}

\textbf{I. L'engagement : fondement historique et fonction sociale vitale de la littérature africaine (Thèse)}

\begin{itemize}
    \item \textbf{A. Une littérature née de la nécessité décoloniale : l'écrivain comme « éclaireur ».}
    \textit{Arg.} : Contexte colonial : l'écriture est acte de résistance, réappropriation de l'histoire.
    \textit{Ex.} : \textbf{Césaire}, \emph{Cahier d'un retour au pays natal} : \og Ma négritude n'est pas une tour d'ivoire \fg. \textbf{Senghor}, \emph{Liberté I} : \og L'émotion est nègre, la raison hellène \fg (mais au service de la libération).
    \textit{Trans. A$\to$B :} \og Si l'engagement est une réponse à l'urgence historique, il structure aussi une esthétique de la dénonciation au quotidien. \fg
    \item \textbf{B. La dénonciation du néocolonialisme et des dictatures : la littérature comme contre-pouvoir.}
    \textit{Arg.} : Post-indépendance : l'écrivain surveille les nouveaux maîtres. Roman à thèse, réalisme critique.
    \textit{Ex.} : \textbf{Mongo Beti}, \emph{Main basse sur le Cameroun} (essai/pamphlet interdit). \textbf{Ferdinand Oyono}, \emph{Le Vieux Nègre et la Médaille} : l'ironie dévoile la médaille coloniale comme symbole du mensonge (\og La médaille, c'est pour les chiens \fg - pensée de Meka). \textbf{Sembène Ousmane}, \emph{Les Bouts de bois de Dieu} : grève = roman choral engagé.
    \textit{Trans. B$\to$C :} \og Cet engagement collectif trouve son accomplissement dans la réhabilitation de la parole africaine elle-même. \fg
    \item \textbf{C. La réhabilitation culturelle : engagement identitaire et linguistique.}
    \textit{Arg.} : Redonner dignité aux langues, aux mythes, aux valeurs bafouées. Négritude = engagement culturel.
    \textit{Ex.} : \textbf{Cheikh Anta Diop}, \emph{Nations nègres et culture} (essai historique engagé). \textbf{Birago Diop}, \emph{Les Contes d'Amadou Koumba} (sauvegarde oralité).
    \textit{Trans. I$\to$II :} \og \textbf{Toutefois}, cette identification totale du sens à l'engagement risque de transformer l'œuvre en simple instrument, niant sa spécificité artistique. \fg
\end{itemize}

\textbf{II. Les limites de l'engagement exclusif : risque de propagande, négation de l'esthétique et de la liberté (Antithèse)}

\begin{itemize}
    \item \textbf{A. Le piège du pamphlet : quand le message tue la forme.}
    \textit{Arg.} : Littérature à thèse, manichéisme, personnages allégoriques, style négligé. Œuvres datées, illisibles hors contexte.
    \textit{Ex.} : Critique de \textbf{Wole Soyinka} : \og Un tigre ne proclame pas sa tigritude, il agit \fg (refus de la négritude essentialiste). Romans \og à thèse \fg oubliés vs chefs-d'œuvre esthétiques.
    \textit{Trans. A$\to$B :} \og Au-delà du risque artistique, c'est la liberté même de l'écrivain qui est menacée par l'assignation à l'engagement. \fg
    \item \textbf{B. Le droit à l'intime, à l'imaginaire, à l'universel : l'écrivain n'est pas un porte-parole.}
    \textit{Arg.} : Exploration psychologique, onirique, fantastique. Littérature \og pour l'art \fg ou \og pour l'homme \fg.
    \textit{Ex.} : \textbf{Amos Tutuola}, \emph{L'Ivrogne dans la brousse} (imaginaire yoruba pur). \textbf{Alain Mabanckou}, \emph{Verre cassé} (bar, oralité, humour, fragmentation, pas de thèse politique). \textbf{Véronique Tadjo}, \emph{Reine Pokou} (mythe réécrit poétiquement).
    \textit{Trans. B$\to$C :} \og Cette liberté conquise permet paradoxalement un engagement plus profond, porté par la forme elle-même. \fg
    \item \textbf{C. Le refus de l'étiquette \og littérature africaine \fg : vers la \og littérature-monde \fg.}
    \textit{Arg.} : Catégorisation réductrice (marketing, essentialisme). Revendication de l'universel par le singulier.
    \textit{Ex.} : \textbf{Édouard Glissant}, \emph{Traité du Tout-Monde} : \og Relation \fg, \og Créolisation \fg. \textbf{Achille Mbembe}, \emph{Critique de la raison nègre}. \textbf{Chinua Achebe} : \og L'écrivain africain doit écrire pour le monde \fg.
    \textit{Trans. II$\to$III :} \og \textbf{Dès lors}, l'opposition stérile entre engagement et esthétique cède la place à leur fusion dans les plus grandes œuvres. \fg
\end{itemize}

\textbf{III. La réconciliation : l'engagement par la forme, l'esthétique comme résistance (Synthèse)}

\begin{itemize}
    \item \textbf{A. L'engagement \emph{est} une poétique : la langue comme champ de bataille.}
    \textit{Arg.} : Décoloniser la langue du colon (français) par la syntaxe, le vocabulaire, le rythme africains. Le style \emph{est} l'engagement.
    \textit{Ex.} : \textbf{Ahmadou Kourouma}, \emph{Les Soleils des Indépendances} : \og Français malinké \fg (syntaxe africaine, proverbes) = résistance linguistique majeure. \textbf{Ken Bugul}, \emph{Le Baobab fou} : écriture fragmentée = identité morcelée.
    \textit{Trans. A$\to$B :} \og Cette invention formelle permet d'atteindre l'universel par le singulier le plus achevé. \fg
    \item \textbf{B. L'universel par le singulier : la grande œuvre engagée transcende son contexte.}
    \textit{Arg.} : Tragédie humaine universelle ancrée dans une histoire précise. Valeur esthétique = pérennité de l'engagement.
    \textit{Ex.} : \textbf{Chinua Achebe}, \emph{Le Monde s'effondre} : destin d'Okonkwo = tragédie grecque (hubris, destin) en pays Igbo. \textbf{Yasmina Khadra}, \emph{Ce que le jour doit à la nuit} : histoire algérienne, roman d'amour/formation universel.
    \textit{Trans. B$\to$C :} \og Enfin, l'écrivain assume sa double allégeance : citoyen et artiste, sans les dissocier. \fg
    \item \textbf{C. L'écrivain citoyen : double allégeance (art + cité) pour une éthique de la responsabilité.}
    \textit{Arg.} : L'engagement ne s'arrête pas au livre. L'auteur intervient dans la cité (prix Nobel, prison, tribunes), mais son œuvre littéraire reste une œuvre d'art aboutie.
    \textit{Ex.} : \textbf{Wole Soyinka} (Prix Nobel, prison sous Abacha, engagement politique direct \emph{et} théâtre/poésie de haute volée). \textbf{Abdourahman Waberi} (Djibouti/France) : engagement pour migrants \emph{par} la poésie (\emph{La Divine Chanson}) et le roman.
\end{itemize}

\textbf{CONCLUSION}
\begin{itemize}
    \item \textbf{Bilan :} La littérature africaine a \textbf{historiquement} puisé son sens dans l'engagement (I : décolonisation, dénonciation, réhabilitation). Cependant, réduire son sens \textbf{uniquement} à l'engagement (II) la condamne au pamphlet, nie la liberté créatrice et l'universel esthétique. Les plus grandes œuvres (III) montrent que l'engagement le plus fort passe par l'exigence formelle (langue, style, mythe) : l'esthétique \emph{est} l'éthique.
    \item \textbf{Ouverture :} Cette tension engage/esthétique traverse toutes les littératures \og périphériques \fg ou \og mineures \fg (caraïbe, maghrébine, amérindienne). Elle pose la question universelle : \textbf{L'art a-t-il une dette envers la société ?} (Réf. \textbf{Albert Camus}, \emph{L'Homme révolté} : \og L'artiste se met au service de la vérité et de la liberté \fg).
\end{itemize}
\end{exoresolu}

\begin{methode}[Rédiger l'introduction de dissertation (Mise en œuvre)]
    \textbf{Objectif :} Produire un paragraphe unique, fluide, sans saut de ligne, qui respecte le canevas imposé et donne envie de lire la suite.
    \begin{enumerate}
        \item \textbf{Rédiger l'amorce (1-2 phrases) :} Éviter les banalités (\og Depuis la nuit des temps \fg). Privilégier : citation courte (auteur, titre, guillemets), définition précise du concept clé, fait marquant, référence artistique.
        \item \textbf{Poser le sujet (1 phrase) :} Reprendre les termes du sujet en les expliquant (paraphrase). \og Le sujet affirme que... \fg / \og La citation de X invite à réfléchir sur... \fg
        \item \textbf{Diviser le sujet / Montrer le débat (1-2 phrases) :} \og Si l'on peut penser que... (thèse), force est de constater que... (antithèse). \fg Utiliser des connecteurs de concession/opposition (\og Certes... cependant \fg, \og Si... toutefois \fg).
        \item \textbf{Formuler la problématique (1 phrase, point d'interrogation) :} La question centrale préparée. Elle doit être \textbf{inscrite visuellement} (souvent en fin de paragraphe ou isolée par des tirets).
        \item \textbf{Annoncer le plan (1 phrase fluide) :} \og Nous analyserons dans un premier temps... (I). Dans un second temps... (II). Enfin... (III). \fg \textbf{Interdiction :} \og Mon plan sera en trois parties \fg, \og Je vais dire que \fg. Utiliser \og Nous \fg (modestie académique) ou l'impersonnel.
        \item \textbf{Relecture :} Vérifier la cohérence (amorce $\rightarrow$ problématique), l'absence de hors-sujet, la correction syntaxique (pas de répétitions lourdes).
    \end{enumerate}
\end{methode}

\begin{exoresolu}[Rédaction de l'introduction — Application au plan détaillé précédent]
\textbf{Consigne :} Rédiger l'introduction complète en un seul paragraphe, à partir des éléments du plan détaillé de l'exercice résolu n°3.
\tcblower
\textbf{Solution détaillée :}

\og \textbf{« Ma négritude n'est pas une tour d'ivoire »,} proclamait Aimé Césaire dans son \emph{Cahier d'un retour au pays natal}, affirmant par là l'indissociabilité de l'acte d'écrire et de l'engagement historique pour la dignité noire. \textbf{Le sujet affirme que la littérature africaine n'a de sens que si elle est engagée,} c'est-à-dire si elle assume une fonction politique, sociale ou culturelle explicite, au risque de nier sa valeur purement esthétique. \textbf{Certes, l'histoire de cette littérature, de la Négritude aux écritures post-indépendance, semble valider ce lien organique : l'écrivain s'est souvent fait le porte-voix des opprimés, le dénonciateur du colonialisme puis du néocolonialisme.} \textbf{Toutefois, la vitalité contemporaine des lettres africaines, marquée par l'exploration de l'intime, du fantastique ou de la pure forme, conteste cette assignation réductrice et revendique le droit à l'universel par l'esthétique.} \textbf{Dans quelle mesure l'engagement constitue-t-il l'unique horizon de sens de la littérature africaine, au risque de nier sa dimension esthétique universelle et la liberté de l'écrivain ?} \textbf{Nous analyserons dans un premier temps comment l'engagement fonde historiquement le sens de cette littérature (I). Nous examinerons ensuite les périls d'un engagement exclusif qui nie l'art et la liberté (II). Enfin, nous verrons comment les grandes œuvres réconcilient engagement et exigence esthétique dans une « littérature-monde » (III).} \fg

\textbf{Analyse de la rédaction (Auto-évaluation) :}
\begin{itemize}
    \item \textbf{Amorce :} Citation Césaire exacte, titre en italique, lien direct avec le sujet (engagement vs tour d'ivoire).
    \item \textbf{Sujet posé :} Reformulation claire (\og n'a de sens que si \fg $\rightarrow$ \og conditionnée par sa fonction \fg).
    \item \textbf{Sujet divisé :} Structure \og Certes... Toutefois \fg efficace. Mots-clés : \og histoire \fg, \og vitalité contemporaine \fg, \og assignation réductrice \fg.
    \item \textbf{Problématique :} Reprise exacte de la fiche méthode, point d'interrogation visible.
    \item \textbf{Annonce du plan :} Verbes variés (\og analyserons \fg, \og examinerons \fg, \og verrons \fg), numérotation (I, II, III) présente, termes du plan respectés.
    \item \textbf{Style :} Un seul bloc, connecteurs logiques, \og Nous \fg académique, vocabulaire précis (\og indissociabilité \fg, \og assignation \fg, \og réconcilient \fg).
\end{itemize}
\end{exoresolu}

\begin{methode}[Rédiger un paragraphe argumenté avec citation et transition (Corps du développement)]
    \textbf{Objectif :} Construire l'unité de base de la dissertation : le paragraphe \cle{argumenté}, \cle{illustré}, \cle{connecté}.
    \begin{enumerate}
        \item \textbf{Phrase d'ouverture (Thèse du paragraphe) :} Une phrase nominale ou verbale qui annonce \textbf{une seule idée force} (le \og sous-titre \fg de la sous-partie A, B ou C). Pas d'exemple ici.
        \item \textbf{Argumentation / Explication (2-3 phrases) :} Développer l'idée : définir, expliquer le mécanisme, montrer la causalité, nuancer. Utiliser le vocabulaire de l'analyse littéraire (focalisation, ironie, intertextualité, style, thématique).
        \item \textbf{Illustration / Preuve (Citation ou référence précise) :} \textbf{Obligatoire.}
            \begin{itemize}
                \item \textbf{Citation courte ( $<$ 3 lignes) :} Intégrée à la phrase (\og ... \fg), guillemets français, référence (Auteur, \emph{Œuvre}, éditeur/année si possible, ou acte/scène/vers pour théâtre).
                \item \textbf{Citation longue ( $>$ 3 lignes) :} Alinéa, retrait, corps plus petit, sans guillemets, référence en bas.
                \item \textbf{Référence précise sans citation :} Résumé d'une scène, description d'un personnage, analyse d'un procédé stylistique (ex: \og L'usage du français malinké chez Kourouma se manifeste par l'inversion syntaxique... \fg).
            \end{itemize}
        \item \textbf{Commentaire de la citation (Lien Idée/Exemple) :} \textbf{Ne jamais laisser une citation parler seule.} Expliquer \textbf{pourquoi} elle prouve l'argument. Analyser un mot, une image, une structure.
        \item \textbf{Mini-transition (ou phrase de chute) :} Clore le paragraphe en annonçant l'idée suivante (sous-partie suivante) ou en résumant l'apport de ce paragraphe à la partie. Utiliser : \og Ainsi, \fg, \og Dès lors, \fg, \og C'est pourquoi \fg + annonce du suivant (\og ... ce qui nous amène à examiner... \fg).
    \end{enumerate}
    \textbf{Règle d'or :} 1 Paragraphe = 1 Idée Force. Sauter une ligne entre chaque paragraphe.
\end{methode}

\begin{exoresolu}[Rédaction d'un paragraphe type — Sous-partie I.B : \og La dénonciation du néocolonialisme \fg]
\textbf{Consigne :} Rédiger le paragraphe complet correspondant à la sous-partie I.B du plan détaillé (Oyono, Beti, Sembène), en respectant la méthode en 5 étapes.
\tcblower
\textbf{Solution détaillée :}

\og \textbf{La dénonciation du néocolonialisme et des dictatures post-indépendance transforme l'écrivain en un contre-pouvoir vigilant, dont l'arme est le réalisme critique.} \textbf{Après les indépendances, la littérature bascule de la lutte anti-coloniale à la critique des nouvelles élites prédatrices, dévoilant la trahison des idéaux nationalistes.} \textbf{C'est ce que met en scène Ferdinand Oyono dans \emph{Le Vieux Nègre et la Médaille} : le protagoniste Meka, après une vie de loyauté envers l'administration coloniale, reçoit une médaille qui se révèle être une insulte dérisoire.} \textbf{La pensée intérieure du vieux homme, \og La médaille, c'est pour les chiens \fg, résume par une métaphore brutale la lucidité tardive du colonisé face à l'illusion de la reconnaissance occidentale.} \textbf{De même, Mongo Beti, dans \emph{Main basse sur le Cameroun}, use du pamphlet documenté pour dénoncer la \og main basse \fg néocoloniale française sur l'économie et la politique camerounaise, tandis que Sembène Ousmane, dans \emph{Les Bouts de bois de Dieu}, donne une dimension épique et collective à la grève de 1947-48, faisant du roman choral le miroir d'une conscience de classe naissante.} \textbf{Ainsi, par l'ironie, le pamphlet ou l'épopée réaliste, ces œuvres font de l'engagement une exigence de vérité historique et sociale, ce qui nous amène à interroger la dimension culturelle de cet engagement : la réhabilitation de l'identité africaine.} \fg

\textbf{Analyse structurelle (Pourquoi ça marche) :}
\begin{enumerate}
    \item \textbf{Phrase d'ouverture :} \og La dénonciation... contre-pouvoir vigilant \fg = Idée force claire (Thèse du §).
    \item \textbf{Argumentation :} \og Après les indépendances... trahison des idéaux \fg = Contexte + Enjeu (explication).
    \item \textbf{Illustration 1 (Oyono) :} Citation courte intégrée (\og La médaille, c'est pour les chiens \fg) + Analyse de la métaphore (\og lucidité tardive \fg).
    \item \textbf{Illustration 2/3 (Beti, Sembène) :} Références précises (titres en italique, résumé d'action/thème) sans citation littérale mais analyse du genre (\og pamphlet \fg, \og roman choral \fg).
    \item \textbf{Commentaire synthétique :} \og Par l'ironie, le pamphlet ou l'épopée... exigence de vérité \fg = Montée en généralité, lien avec la thèse de la partie I.
    \item \textbf{Transition :} \og Ainsi... ce qui nous amène à interroger... \fg = Lien logique (conséquence) + Annonce précise de la sous-partie I.C (\og réhabilitation de l'identité \fg).
\end{enumerate}
\end{exoresolu}

\begin{methode}[Rédiger la conclusion de dissertation (Clôture)]
    \textbf{Objectif :} Refermer la copie sur une note forte, synthétique et ouverte. Deux temps distincts, souvent séparés par un alinéa ou une ligne blanche.
    \begin{enumerate}
        \item \textbf{Le Bilan (Synthèse) :} \textbf{Ne pas faire une liste de courses} (I a dit ça, II a dit ça). Répondre \textbf{directement à la problématique} en fusionnant les trois parties.
            \begin{itemize}
                \item Structure type : \og Au terme de cette étude, il apparaît que [Réponse thèse I] ; cependant [Nuance thèse II] ; finalement [Dépassement thèse III]. \fg
                \item Utiliser des connecteurs de synthèse : \og En définitive \fg, \og En somme \fg, \og Globalement \fg, \og Force est de constater que \fg.
            \end{itemize}
        \item \textbf{L'Ouverture (Prolongement) :} Une seule piste, pertinente, non traitée dans le développement.
            \begin{itemize}
                \item \textbf{Types d'ouvertures validées :} Autre œuvre du même auteur / autre auteur même thème / autre art (cinéma, peinture, musique) / actualité brûlante / question philosophique / perspective historique future.
                \item \textbf{Interdits :} \og Pour conclure, je dirais que... \fg, nouvelle problématique, banalité (\og La littérature est belle \fg), hors-sujet total.
                \item \textbf{Formulation :} \og Cette réflexion invite à se demander si... \fg / \og On pourrait prolonger en étudiant... \fg / \og Ce débat résonne aujourd'hui avec... \fg
            \end{itemize}
        \item \textbf{Soin formel :} Dernière phrase soignée (rythme, vocabulaire). Dernière image marquante.
    \end{enumerate}
\end{methode}

\begin{exoresolu}[Rédaction de la conclusion — Application au sujet]
\textbf{Consigne :} Rédiger la conclusion (Bilan + Ouverture) en deux paragraphes distincts, à partir du plan détaillé.
\tcblower
\textbf{Solution détaillée :}

\textbf{Bilan :}
\og Au terme de cette analyse, il apparaît que l'engagement a constitué, historiquement, \textbf{l'âme même de la littérature africaine} : née dans l'urgence décoloniale, elle a trouvé son sens premier dans la dénonciation de l'oppression et la réhabilitation de l'identité bafouée (I). \textbf{Toutefois}, réduire son horizon à la seule fonction militante, c'est la condamner au pamphlet éphémère et nier la liberté souveraine de l'écrivain face à l'histoire (II). \textbf{En définitive}, les plus grandes œuvres nous enseignent que l'engagement le plus radical et le plus durable est celui qui passe par l'exigence esthétique : en inventant une langue, en réinventant le mythe, en assumant l'universel par le singulier, l'écrivain africain fait de sa beauté une arme et de son art une éthique (III). \fg

\textbf{Ouverture :}
\og Cette dialectique de l'engagement et de l'esthétique dépasse largement le cadre africain : elle traverse toute \og littérature de combat \fg ou \og littérature mineure \fg (pensons à la littérature caribéenne de Glissant ou Chamoiseau, à la littérature maghrébine de Kateb Yacine ou Assia Djebar, ou encore aux littératures autochtones d'Amérique). Elle pose, in fine, la question universelle de la responsabilité de l'artiste : \textbf{l'œuvre d'art doit-elle être utile pour être vraie, ou sa vérité est-elle sa seule utilité ?} \fg

\textbf{Analyse :}
\begin{itemize}
    \item \textbf{Bilan :} Réponse nuancée (\og âme même \fg / \og toutefois \fg / \og en définitive \fg). Fusion I/II/III sans liste. Vocabulaire fort (\og âme \fg, \og pamphlet éphémère \fg, \og liberté souveraine \fg, \og exigence esthétique \fg, \og beauté une arme \fg).
    \item \textbf{Ouverture :} Élargissement géographique/culturel pertinent (Caraïbes, Maghreb, Amérindiens) $\rightarrow$ concept de \og littérature mineure \fg (Deleuze/Guattari/Kafka). Question philosophique finale (Utilité vs Vérité) qui résonne avec Camus/Sartre. Pas de nouveau développement.
\end{itemize}
\end{exoresolu}

\begin{methode}[Assurer cohésion et cohérence : Référents, Progression thématique, Connecteurs (Outils linguistiques)]
    \textbf{Objectif :} Garantir la fluidité et la lisibilité de la copie (critères de notation \og Expression \fg / \og Structure \fg). Mobiliser les savoirs de l'unité \og Langue française \fg.
    \begin{enumerate}
        \item \textbf{Maîtriser la chaîne référentielle (Le référent et ses substituts) :}
            \begin{itemize}
                \item Identifier le \cle{référent} (nom initial : \og L'écrivain \fg, \og Cette littérature \fg).
                \item Varier les \cle{substituts} pour éviter les répétitions lourdes : \textbf{Pronoms} (il, elle, ce dernier, celui-ci), \textbf{Nominales} (synonymes, hyperonymes : \og l'auteur \fg, \og l'intellectuel \fg, \og le romancier \fg, \og cette production \fg), \textbf{Ellipses} (verbe seul).
                \item \textbf{Attention :} Éviter l'ambiguïté (qui fait \og il \fg ?). Répéter le nom propre si changement de référent ou risque de confusion.
            \end{itemize}
        \item \textbf{Structurer la progression thématique (Thème / Rhème) :}
            \begin{itemize}
                \item \cle{Thème} (Point de départ, connu) $\rightarrow$ \cle{Rhème} (Information nouvelle, fin de phrase).
                \item \textbf{Progression linéaire (enchaînée) :} Rhème 1 = Thème 2. (\og L'écrivain engage sa parole. \textbf{Cette parole} libère les consciences. \fg) $\rightarrow$ Fluidité narrative.
                \item \textbf{Progression à thème constant (parallèle) :} Même thème, rhèmes différents. (\og L'écrivain dénonce. \textbf{L'écrivain} témoigne. \textbf{L'écrivain} espère. \fg) $\rightarrow$ Insistance, énumération.
                \item \textbf{Progression à rhème constant :} Thèmes différents, même information nouvelle. (Rare en dissertation).
                \item \textbf{Choix stratégique :} Linéaire pour l'argumentation suivie ; Constant pour la synthèse/l'énumération d'effets.
            \end{itemize}
        \item \textbf{Utiliser les connecteurs logiques (Opérateurs argumentatifs) :}
            \begin{itemize}
                \item \textbf{Addition / Renforcement :} De plus, par ailleurs, en outre, qui plus est.
                \item \textbf{Opposition / Concession :} Cependant, néanmoins, pourtant, si... toutefois, Certes... mais.
                \item \textbf{Conséquence / Conclusion :} Par conséquent, dès lors, ainsi, c'est pourquoi, en définitive.
                \item \textbf{Explication / Cause :} Car, parce que, en effet, puisque.
                \item \textbf{Exemplification :} Par exemple, ainsi, notamment, c'est le cas de.
            \end{itemize}
        \item \textbf{Les sigles et emprunts (Lexicologie) :} Définir au premier emploi (\og L'OUA (Organisation de l'Unité Africaine) \fg), italique pour les emprunts non francisés (\og \emph{négritude} \fg, \og \emph{griot} \fg, \og \emph{tout-monde} \fg), romains pour les francisés (\og laser \fg, \og radar \fg $\rightarrow$ mais \og négritude \fg reste souvent en italique comme concept).
    \end{enumerate}
\end{methode}

\begin{exoresolu}[Application linguistique : Cohésion / Progression / Connecteurs]
\textbf{Consigne :} Sur l'extrait de paragraphe suivant (rédaction brouillon), identifier les défauts de cohésion/cohérence et proposer une version corrigée en appliquant la méthode.
\textbf{Texte brouillon :} \og L'écrivain africain est engagé. L'écrivain africain dénonce le colonialisme. L'engagement est nécessaire. Mais l'écrivain africain a aussi le droit d'écrire pour l'art. L'art est important. La négritude est un mouvement. Senghor a dit que l'émotion est nègre. C'est une citation célèbre. \fg
\tcblower
\textbf{Solution détaillée :}

\textbf{Diagnostic des erreurs :}
\begin{itemize}
    \item \textbf{Répétition massive du référent :} \og L'écrivain africain \fg (3 fois), \og L'engagement \fg / \og L'art \fg (nom seul, pas de progression).
    \item \textbf{Aucune progression thématique :} Phrases juxtaposées, pas de lien Thème/Rhème. Style télégraphique.
    \item \textbf{Absence de connecteurs logiques :} Pas de liaison \og Mais \fg (faible), pas de cause, pas de conséquence.
    \item \textbf{Citation non intégrée, non commentée :} \og C'est une citation célèbre \fg (commentaire nul).
    \item \textbf{Sigles/Emprunts :} \og négritude \fg non mis en valeur (italique).
\end{itemize}

\textbf{Version corrigée (Application de la méthode) :}
\og \textbf{L'écrivain africain}, \textbf{historiquement investi d'une mission de libération}, \textbf{s'est fait le porte-voix des peuples opprimés.} \textbf{Cet engagement} \textbf{ne se limite pas à la dénonciation du colonialisme} ; \textbf{il structure une esthétique de la résistance, comme l'illustre la \emph{Négritude}.} \textbf{Toutefois}, \textbf{le créateur ne saurait être réduit à un militant : il revendique le droit à l'intime et à la pure forme.} \textbf{Ainsi}, \textbf{lorsque Senghor affirme que \og l'émotion est nègre, la raison hellène \fg, il ne fige pas l'identité dans l'essentialisme, mais souligne la spécificité d'une sensibilité mise au service de l'universel.} \textbf{Dès lors}, \textbf{l'engagement le plus fécond est celui qui, par la maîtrise de la langue et de la forme, transcende la circonstance historique.} \fg

\textbf{Analyse des corrections apportées :}
\begin{enumerate}
    \item \textbf{Chaîne référentielle :} \og L'écrivain africain \fg (Référent initial) $\rightarrow$ \textbf{Cet engagement} (Substitut nominal, reprise Thème) $\rightarrow$ \textbf{il} (Pronom, Thème constant) $\rightarrow$ \textbf{le créateur} (Synonyme/Hyperonyme, variation) $\rightarrow$ \textbf{il} (Reprise) $\rightarrow$ \textbf{Senghor} (Nouveau référent précis, cité) $\rightarrow$ \textbf{l'engagement le plus fécond} (Nominalisation abstraite, synthèse).
    \item \textbf{Progression thématique :}
        \begin{itemize}
            \item Ph 1 : Thème \og L'écrivain \fg / Rhème \og porte-voix \fg.
            \item Ph 2 : Thème \og Cet engagement \fg (Rhème 1 nominalisé) / Rhème \og esthétique résistance \fg. \textbf{Progression linéaire.}
            \item Ph 3 : Thème \og le créateur \fg (Nouveau thème, substitution) / Rhème \og droit à l'intime \fg.
            \item Ph 4 : Thème \og Senghor \fg / Rhème \og citation \fg. Commentaire intégré.
            \item Ph 5 : Thème \og l'engagement le plus fécond \fg (Synthèse abstraite) / Rhème \og transcende la circonstance \fg.
        \end{itemize}
    \item \textbf{Connecteurs :} \og ; \textbf{Toutefois} \fg (Opposition), \og \textbf{Ainsi} \fg (Illustration/Explication), \og \textbf{Dès lors} \fg (Conséquence/Synthèse).
    \item \textbf{Citation :} Intégrée (\og affirme que \fg), guillemets, commentée (\og ne fige pas... mais souligne \fg).
    \item \textbf{Emprunt :} \emph{Négritude} en italique (concept étranger/technique).
\end{enumerate}
\end{exoresolu}

\begin{attention}[Les 5 erreurs qui coûtent des points au Bac (Dissertation)]
    \begin{enumerate}
        \item \textbf{Problématique absente, mal formulée ou hors-sujet :} C'est la \cle{clé de voûte}. Sans problématique vraie (question ouverte, tension), le plan est plat (catalogue) ou hors-sujet. \textbf{Sanction :} Note plafond souvent 6-7/20 (hors-sujet apparent) ou 0 si copie blanche. \textbf{Reflexe :} La formuler \textbf{avant} le plan, la faire valider mentalement : \og Est-elle répondable par un plan dialectique ? \fg
        \item \textbf{Plan binaire (Pour / Contre) au lieu de dialectique (Thèse / Antithèse / Synthèse) :} Le \og Pour/Contre \fg est une liste d'arguments, pas une démonstration qui progresse. La Synthèse (III) n'est pas \og mélange des deux \fg mais un \textbf{niveau supérieur de compréhension}. \textbf{Sanction :} Plafond 10-11/20 (structure incomplète). \textbf{Reflexe :} Vérifier que la Partie III \og dépasse \fg le débat I/II.
        \item \textbf{Absence de citations / exemples précis / œuvres au programme :} Dissertation \og hors-sol \fg, généralités (\og Les écrivains africains disent que... \fg). \textbf{Sanction :} Perte de 4 à 6 points sur la note de fond (sur 10-12 pts). \textbf{Reflexe :} Au brouillon, \textbf{chaque sous-partie (A, B, C) doit avoir sa citation/référence exacte} (Auteur, Titre en italique, passage précis).
        \item \textbf{Paragraphe « pâté » (pas de saut de ligne) ou 1 phrase = 1 paragraphe :} Illisible pour le correcteur. Structure invisible. \textbf{Sanction :} Pénalité Expression/Structure (2-3 pts). \textbf{Reflexe :} \textbf{Sauter une ligne} entre chaque grande idée (A, B, C). Aérer la copie.
        \item \textbf{Conclusion baclée ou absente / Ouverture « poubelle » :} Conclusion = résumé liste de courses. Ouverture = \og On pourrait parler de la musique \fg (hors-sujet) ou nouvelle problématique. \textbf{Sanction :} Perte des points de « Maîtrise de la structure » et « Ouverture » (souvent 2-3 pts). \textbf{Reflexe :} Réserver 10-15 min pour la conclusion. Bilan = Réponse à la problématique. Ouverture = Prolongement \textbf{logique et ciblé} (autre art, même thème ; actualité ; philo).
    \end{enumerate}
\end{attention}

\newpage

\coursec{Exercices}

\courssub{Je vérifie mes connaissances}

\begin{exercice}[QCM : les étapes de la dissertation]
Pour chaque question, entoure la bonne réponse.

\textbf{1.} La problématique d'un sujet de dissertation est :
\begin{itemize}
\item[a)] une simple reformulation du sujet ;
\item[b)] la question centrale qui naît de la tension entre les termes du sujet ;
\item[c)] le titre de l'œuvre étudiée.
\end{itemize}

\textbf{2.} Un plan détaillé comporte obligatoirement :
\begin{itemize}
\item[a)] uniquement les titres des parties ;
\item[b)] les parties, les sous-parties, les arguments et les exemples ;
\item[c)] la rédaction complète de chaque paragraphe.
\end{itemize}

\textbf{3.} Dans \emph{Le Vieux Nègre et la médaille} de Ferdinand Oyono, Meka reçoit une médaille :
\begin{itemize}
\item[a)] pour son courage au combat ;
\item[b)] pour sa longue patience et sa collaboration avec l'administration coloniale ;
\item[c)] pour sa richesse.
\end{itemize}

\textbf{4.} Un sigle est :
\begin{itemize}
\item[a)] un mot emprunté à une langue étrangère ;
\item[b)] une abréviation formée des initiales de plusieurs mots (ex. : MINESEC) ;
\item[c)] un synonyme rare.
\end{itemize}
\end{exercice}

\begin{exercice}[Vrai ou faux — justifie ta réponse]
Indique si chaque affirmation est vraie ou fausse, puis justifie en une ou deux phrases.

\begin{enumerate}
\item L'introduction d'une dissertation peut commencer directement par l'annonce du plan.
\item La transition entre deux parties sert à montrer le lien logique entre les idées.
\item Dans un texte, le référent et ses substituts (pronoms, reprises nominales) assurent la cohésion du texte.
\item Le plan dialectique se compose de deux parties seulement.
\end{enumerate}
\end{exercice}

\begin{exercice}[Définitions]
Définis en deux ou trois phrases chacune des notions suivantes, puis donne un exemple :
\begin{enumerate}
\item la problématique ;
\item le plan détaillé ;
\item l'amorce (ou accroche) de l'introduction ;
\item le sigle et l'emprunt lexical.
\end{enumerate}
\end{exercice}

\begin{exercice}[Repérage : référent et substituts]
Lis l'extrait suivant, inspiré de la situation du roman :

« \emph{Meka attendit longtemps sous le soleil. Le vieil homme s'essuyait le front. Il ne comprenait pas pourquoi le Commandant tardait tant. Celui-ci finit par apparaître, entouré de ses gardes.} »

\begin{enumerate}
\item Souligne tous les substituts qui renvoient au référent « Meka ».
\item Relève le substitut qui renvoie au référent « le Commandant ».
\item Explique en trois lignes le rôle de ces substituts dans la cohésion du texte.
\end{enumerate}
\end{exercice}

\courssub{Je m'entraîne}

\begin{exercice}[De l'analyse du sujet à la problématique]
Soit le sujet : « \emph{La médaille remise à Meka est-elle une récompense ou une humiliation ?} »

\begin{enumerate}
\item Dégage le sens de chacun des termes importants du sujet (médaille, récompense, humiliation).
\item Formule deux questions secondaires qui découlent du sujet.
\item Rédige la problématique en une phrase interrogative.
\end{enumerate}
\end{exercice}

\begin{exercice}[Recherche et classement des idées]
Voici une liste d'idées en désordre, relatives au sujet : « \emph{La médaille est-elle une récompense ou un piège ?} »

\begin{itemize}
\item la médaille flatte l'orgueil de Meka ;
\item la médaille divise la communauté de Doum ;
\item Meka est fier devant les Blancs ;
\item après la cérémonie, Meka est arrêté et humilié ;
\item la médaille donne l'illusion d'une reconnaissance ;
\item les anciens du village se méfient du cadeau des Blancs ;
\item la médaille coûte à Meka sa dignité et sa liberté.
\end{itemize}

\begin{enumerate}
\item Classe ces idées en deux colonnes : « récompense » / « piège ».
\item Propose un titre pour chacune des deux parties du plan.
\item Choisis dans chaque partie l'idée la plus forte et justifie ton choix en deux lignes.
\end{enumerate}
\end{exercice}

\begin{exercice}[Du paragraphe narratif au paragraphe argumentatif]
Voici un paragraphe narratif inspiré du roman :

« \emph{Le jour de la cérémonie, Meka revêtit son boubou neuf. Il marcha jusqu'à la place publique, le cœur battant. Quand le Commandant accrocha la médaille sur sa poitrine, la foule applaudit. Meka souriait.} »

Transforme ce paragraphe en paragraphe argumentatif montrant que la médaille est un instrument de domination coloniale. Pour cela :
\begin{enumerate}
\item remplace le passé simple par le présent de vérité générale ;
\item fais passer l'énonciation de la troisième personne à une énonciation argumentative (on, nous, ou tournures impersonnelles) ;
\item insère au moins trois connecteurs logiques (en effet, cependant, par conséquent, d'abord, ainsi...).
\end{enumerate}
\end{exercice}

\begin{exercice}[Sigles et emprunts dans la presse camerounaise]
Voici un extrait adapté d'un article de presse :

« \emph{La MINAT a annoncé un nouveau plan de modernisation des chefferies. Les chefs traditionnels recevront des ordinateurs et une connexion Internet pour gérer les états civils. Le projet, soutenu par le PNUD, prévoit aussi des séminaires de formation.} »

\begin{enumerate}
\item Relève trois sigles et donne leur signification complète (recherche si nécessaire).
\item Relève deux emprunts à la langue anglaise et propose pour chacun un équivalent français.
\item Explique en trois lignes pourquoi la langue administrative camerounaise utilise autant de sigles et d'emprunts.
\end{enumerate}
\end{exercice}

\courssub{Je me prépare au Bac}

\begin{exercice}[Sujet de type Probatoire : lecture méthodique et commentaire]
\textbf{Texte} (adapté de Ferdinand Oyono, \emph{Le Vieux Nègre et la médaille}) :

« \emph{Le Commandant prit la médaille dans l'écrin. Il la fit briller un instant au soleil, comme pour mieux la montrer à la foule assemblée. Meka, immobile, sentait sur lui tous les regards. Quand l'insigne fut épinglé sur son boubou, le Commandant recula d'un pas, le contempla avec un sourire paternel, puis se tourna vers les notables : "Voilà, dit-il, un homme qui a su comprendre la France." La foule applaudit. Meka ne savait pas encore que ce sourire-là avait un prix.} »

\textbf{Travail à faire :}
\begin{enumerate}
\item Présente le texte (auteur, œuvre, nature du passage, situation).
\item Dégage le thème principal du passage.
\item Analyse les rapports de force entre le Commandant et Meka lors de cette cérémonie : gestes, paroles, regards, position des personnages.
\item Montre, en relevant deux procédés précis (champ lexical, ironie, point de vue), comment le narrateur invite le lecteur à se méfier de cette récompense.
\item Rédige une conclusion ouverte sur la signification de la médaille dans l'œuvre.
\end{enumerate}
\end{exercice}

\begin{exercice}[Sujet de type Baccalauréat technique : dissertation sur l'œuvre]
\textbf{Sujet :} « \emph{La médaille est-elle une récompense ou un piège ?} »

\textbf{Travail à faire :}
\begin{enumerate}
\item Analyse le sujet : définition des termes, repérage de la tension entre « récompense » et « piège ».
\item Formule la problématique.
\item Élabore un plan détaillé complet (deux ou trois parties), avec pour chaque partie : l'idée directrice, deux arguments et un exemple précis tiré de \emph{Le Vieux Nègre et la médaille}.
\item Rédige l'introduction complète de la dissertation (amorce, présentation de l'œuvre, problématique, annonce du plan).
\item Rédige la conclusion (bilan répondant à la problématique et ouverture).
\end{enumerate}
\end{exercice}

\begin{exercice}[Cas pratique : construire un plan dialectique à partir d'un article de presse]
\textbf{Document} (extrait adapté d'un article de presse camerounaise) :

« \emph{À Foumban comme à Bafoussam, les chefferies traditionnelles s'adaptent : certains sultans et chefs supérieurs tiennent désormais une page Facebook, règlent les conflits fonciers avec l'aide des sous-préfets et envoient leurs notables suivre des formations en gestion. Pourtant, les rites d'intronisation, les sociétés secrètes et la parole coutumière restent au cœur du pouvoir traditionnel. "Un chef qui oublie les ancêtres n'est plus un chef", confie un notable. Mais un jeune cadre rétorque : "Un chef qui refuse le développement condamne son village."} »

\textbf{Sujet de dissertation :} « \emph{Modernité et tradition sont-elles compatibles dans la gouvernance locale ?} »

\textbf{Travail à faire :}
\begin{enumerate}
\item Dégage du document deux positions opposées et reformule-les chacune en une phrase.
\item Construis un plan dialectique détaillé en trois parties (thèse / antithèse / synthèse), en précisant pour chaque partie l'idée directrice, deux arguments et un exemple tiré du document ou de la réalité camerounaise.
\item Rédige la transition entre la première et la deuxième partie.
\item Rédige un paragraphe du développement comportant au moins une citation du document correctement introduite (verbe de parole + guillemets) et commentée.
\end{enumerate}
\end{exercice}

\newpage

\coursec{Corrigés des exercices}

\coursec{Corrigés des exercices — Leçon 13 : Produire un plan détaillé d'une dissertation}

\courssub{Analyse du sujet, problématique et outils linguistiques (Exercices 1 à 5)}

\begin{corrige}[Exercice 1 — QCM : les étapes de la dissertation]
\textbf{1. Réponse b).} La problématique n'est pas une simple reformulation du sujet : elle naît de la \cle{tension} entre les termes du sujet (par exemple entre « récompense » et « humiliation ») et s'exprime par une question centrale qui oriente toute la réflexion.

\textbf{2. Réponse b).} Un plan détaillé présente les \cle{parties}, les \cle{sous-parties}, les \cle{arguments} et les \cle{exemples} (avec références précises). Les titres seuls constituent un plan simple ; la rédaction complète des paragraphes relève de la copie, pas du plan.

\textbf{3. Réponse b).} Dans \emph{Le Vieux Nègre et la médaille}, Meka reçoit la médaille de l'amitié franco-camerounaise pour sa longue patience et sa collaboration avec l'administration coloniale (il a donné ses terres à la mission et ses fils à la guerre), et non pour un acte de guerre ou pour sa richesse.

\textbf{4. Réponse b).} Un \cle{sigle} est une abréviation formée des initiales de plusieurs mots : MINESEC (Ministère des Enseignements secondaires). Un mot emprunté à une langue étrangère est un \cle{emprunt lexical}, ce qui est une autre notion.
\end{corrige}

\begin{corrige}[Exercice 2 — Vrai ou faux]
\begin{enumerate}
\item \textbf{Faux.} L'annonce du plan est la \emph{dernière} étape de l'introduction. Elle doit être précédée de l'amorce, de la présentation de l'œuvre ou du sujet et de la problématique. Commencer par le plan prive l'introduction de sa mise en contexte.
\item \textbf{Vrai.} La transition est un court passage (une ou deux phrases) qui conclut la partie précédente et annonce la suivante : elle rend visible la \cle{progression logique} du devoir et assure la cohérence de l'ensemble.
\item \textbf{Vrai.} Le référent (la réalité désignée) est repris par des substituts : pronoms (\emph{il}, \emph{celui-ci}) et reprises nominales (\emph{le vieil homme}). Ces reprises évitent les répétitions et tissent la \cle{cohésion} du texte.
\item \textbf{Faux.} Le plan dialectique comporte \emph{trois} parties : la thèse, l'antithèse et la synthèse. C'est le plan thématique ou le plan « pour/contre » (confrontation) qui peut se limiter à deux parties.
\end{enumerate}
\end{corrige}

\begin{corrige}[Exercice 3 — Définitions]
\begin{enumerate}
\item \textbf{La problématique} : question centrale, généralement formulée à l'interrogatif, qui naît de la tension entre les termes du sujet et qui commande l'ensemble du devoir. \emph{Exemple} : « La médaille remise à Meka est-elle un signe de reconnaissance ou un instrument de domination ? »
\item \textbf{Le plan détaillé} : schéma complet du devoir avant rédaction, présentant les parties, les sous-parties, les arguments et les exemples précis (avec références aux œuvres). \emph{Exemple} : « I. La médaille comme récompense apparente — 1. Elle flatte l'orgueil de Meka (ex. : la fierté de Meka le jour de la cérémonie) ».
\item \textbf{L'amorce (ou accroche)} : première phrase de l'introduction, qui part d'une idée générale (constat, citation, fait historique) pour conduire naturellement au sujet. \emph{Exemple} : « La colonisation a souvent habillé la domination des habits de la gratitude. »
\item \textbf{Le sigle} : abréviation formée des initiales de plusieurs mots (\emph{ex.} : PNUD, MINAT). \textbf{L'emprunt lexical} : mot pris à une langue étrangère et intégré au français (\emph{ex.} : « ordinateur » est français, mais « Internet », « planning » sont des emprunts à l'anglais).
\end{enumerate}
\end{corrige}

\begin{corrige}[Exercice 4 — Repérage : référent et substituts]
\begin{enumerate}
\item Substituts renvoyant au référent « Meka » : \textbf{Le vieil homme} (reprise nominale), \textbf{Il} (pronom personnel), \textbf{lui} implicite dans « s'essuyait » (le pronom réfléchi \emph{se} renvoie aussi à Meka).
\item Substitut renvoyant au référent « le Commandant » : \textbf{Celui-ci} (pronom démonstratif).
\item \textbf{Rôle} : ces substituts évitent la répétition des noms propres et assurent la \cle{chaîne de référence} : le lecteur suit sans rupture le même personnage d'une phrase à l'autre. Ils garantissent ainsi la cohésion du texte et la clarté de la progression thématique (le thème « Meka » est maintenu en tête de phrase avant d'être remplacé par « le Commandant »).
\end{enumerate}
\end{corrige}

\begin{corrige}[Exercice 5 — De l'analyse du sujet à la problématique]
\begin{enumerate}
\item \textbf{Sens des termes} :
\begin{itemize}
\item \emph{médaille} : insigne métallique décerné par une autorité ; ici, la médaille de l'amitié franco-camerounaise remise à Meka par l'administration coloniale ;
\item \emph{récompense} : bien accordé en échange d'un mérite, signe de reconnaissance qui honore celui qui la reçoit ;
\item \emph{humiliation} : atteinte à la dignité, rabaissement public d'une personne.
\end{itemize}
\item \textbf{Questions secondaires} : Que doit Meka à l'administration coloniale pour être ainsi honoré ? Que se passe-t-il pour lui après la cérémonie ?
\item \textbf{Problématique} : « La médaille décernée à Meka marque-t-elle une véritable reconnaissance de ses mérites, ou n'est-elle que l'instrument d'une domination qui le conduira à l'humiliation ? »
\end{enumerate}
\end{corrige}

\courssub{Recherche d'idées, plan détaillé et rédaction argumentative (Exercices 6 à 7)}

\begin{corrige}[Exercice 6 — Recherche et classement des idées]
\begin{enumerate}
\item \textbf{Classement} :
\begin{center}
\begin{tabularx}{\linewidth}{|X|X|}
\hline
\rowcolor{popblueL} \textbf{Récompense (apparence)} & \textbf{Piège (réalité)} \\
\hline
La médaille flatte l'orgueil de Meka & La médaille suscite la division à Doum \\
\hline
Meka est fier devant les Blancs & La cérémonie est suivie de l'arrestation et de l'humiliation de Meka \\
\hline
La médaille donne l'illusion d'une reconnaissance & Les anciens du village se méfient du cadeau des Blancs \\
\hline
& La médaille coûte à Meka sa dignité et sa liberté \\
\hline
\end{tabularx}
\end{center}
\item \textbf{Titres possibles} : I. « La médaille, une récompense qui flatte et rassure » ; II. « La médaille, un piège qui divise et détruit ».
\item \textbf{Idées les plus fortes} : dans la partie I, « la médaille donne l'illusion d'une reconnaissance », car elle contient déjà en germe la critique (le mot « illusion ») et prépare la transition ; dans la partie II, « la médaille coûte à Meka sa dignité et sa liberté », car elle résume le bilan tragique de l'intrigue (arrestation, perte du respect de soi) et donne sa force à la démonstration.
\end{enumerate}
\end{corrige}

\begin{corrige}[Exercice 7 — Du paragraphe narratif au paragraphe argumentatif]
\textbf{Transformation attendue} (modèle de réponse) :

« \emph{D'abord}, la cérémonie de remise de la médaille met en scène la soumission du colonisé : Meka revêt son boubou neuf et marche vers la place publique, le cœur battant, comme on se rend à une consécration. \emph{En effet}, quand le Commandant accroche l'insigne sur sa poitrine, c'est l'autorité coloniale qui consacre publiquement son serviteur, et la foule applaudit un ordre qu'elle ne conteste pas. \emph{Cependant}, ce sourire de Meka ne trompe personne : on comprend que la récompense n'honore pas l'homme, mais célèbre sa docilité. \emph{Par conséquent}, la médaille apparaît moins comme un hommage que comme un instrument de domination coloniale. »

\textbf{Vérification des consignes} : présent de vérité générale (revêt, marche, accroche, applaudit) ; énonciation argumentative (tournures impersonnelles : « on comprend que », « c'est l'autorité coloniale qui... ») ; quatre connecteurs (\emph{d'abord}, \emph{en effet}, \emph{cependant}, \emph{par conséquent}).
\end{corrige}

\courssub{Étude de la langue : sigles, emprunts et cohérence textuelle (Exercice 8)}

\begin{corrige}[Exercice 8 — Sigles et emprunts dans la presse camerounaise]
\begin{enumerate}
\item \textbf{Trois sigles relevés dans le texte support} :
\begin{itemize}
\item MINAT = Ministère de l'Administration territoriale ;
\item PNUD = Programme des Nations unies pour le développement ;
\item MINESEC = Ministère des Enseignements secondaires.
\end{itemize}
(« États civils » n'est pas un sigle mais une locution nominale.)
\item \textbf{Deux emprunts à l'anglais relevés dans le texte support} :
\begin{itemize}
\item \emph{Internet} (équivalent français : « la toile », « le réseau mondial ») ;
\item \emph{planning} ou \emph{e-mail} selon l'article (équivalents : « emploi du temps », « programme » ; « courriel », « courrier électronique »).
\end{itemize}
\item \textbf{Explication} : la langue administrative camerounaise utilise de nombreux sigles parce que les institutions sont nombreuses et que le sigle permet de les nommer rapidement et sans ambiguïté. Les emprunts à l'anglais s'expliquent par le bilinguisme officiel du Cameroun, par l'influence des technologies (nées dans le monde anglophone) et par le prestige de l'anglais dans l'administration et les médias.
\end{enumerate}
\end{corrige}

\courssub{Intégration : commentaire composé et dissertation type Probatoire / Baccalauréat (Exercices 9 à 10)}

\begin{corrige}[Exercice 9 — Sujet de type Probatoire : lecture méthodique et commentaire]
\textbf{Corrigé rédigé (éléments attendus)} :

\begin{enumerate}
\item \textbf{Présentation} : extrait de \emph{Le Vieux Nègre et la médaille} (1956), roman de l'écrivain camerounais Ferdinand Oyono. Il s'agit d'un passage narratif et descriptif : la cérémonie de remise de la médaille de l'amitié franco-camerounaise au vieux Meka, au village de Doum.
\item \textbf{Thème} : la cérémonie coloniale comme théâtre de la domination déguisée en reconnaissance.
\item \textbf{Rapports de force} : le Commandant mène toute la scène : il « prend » la médaille, la « fait briller », l'« épingle », « recule » pour contempler son œuvre ; Meka, lui, est « immobile », objet des regards, réduit à un support de l'insigne. La parole appartient au seul colonisateur (« Voilà... un homme qui a su comprendre la France ») : Meka ne dit rien. La position verticale du Commandant (il agit, se tourne vers les notables) s'oppose à la passivité de Meka. Le « sourire paternel » installe une relation de supériorité condescendante.
\item \textbf{Deux procédés} : (a) le \emph{champ lexical du spectacle et de l'éclat} (« écrin », « briller au soleil », « montrer à la foule ») montre que la médaille est un objet de mise en scène destiné à éblouir, non à honorer ; (b) l'\emph{ironie} de la dernière phrase (« ce sourire-là avait un prix »), énoncée en focalisation interne proleptique, avertit le lecteur que la récompense sera payée très cher : elle invite à se méfier.
\item \textbf{Conclusion ouverte} : la médaille apparaît ainsi comme le symbole d'une reconnaissance trompeuse, instrument de la propagande coloniale ; on peut s'interroger, au-delà du roman, sur tous les « cadeaux » du colonisateur et sur le prix de la dignité perdue.
\end{enumerate}

\textbf{Barème indicatif (sur 20)} : présentation du texte 3 pts ; thème 2 pts ; analyse des rapports de force 6 pts ; deux procédés relevés et commentés 6 pts ; conclusion ouverte 3 pts.

\textbf{Points de vigilance} : citer le texte à l'appui de chaque analyse (guillemets) ; ne pas résumer le roman à la place de l'analyse ; distinguer clairement le point de vue du narrateur de celui des personnages ; la conclusion doit répondre au thème dégagé et s'ouvrir, non répéter l'introduction.
\end{corrige}

\begin{corrige}[Exercice 10 — Sujet de type Baccalauréat technique : dissertation sur l'œuvre]
\textbf{1. Analyse du sujet.} \emph{Récompense} : bien décerné pour un mérite, qui honore. \emph{Piège} : appât qui attire pour mieux nuire. La tension porte sur la nature réelle de la médaille : signe de gratitude ou instrument de manipulation coloniale ?

\textbf{2. Problématique.} « La médaille remise à Meka constitue-t-elle une authentique récompense, ou n'est-elle qu'un piège tendu par l'administration coloniale pour asservir davantage le colonisé ? »

\textbf{3. Plan détaillé (plan de confrontation en deux parties).}

\textbf{I. La médaille, une récompense apparente qui honore et flatte.}
\begin{itemize}
\item Idée directrice : la médaille donne à Meka l'illusion d'être reconnu et l'élève aux yeux de tous.
\item Argument 1 : elle couronne une vie de services rendus à la France (terres données à la mission, fils morts à la guerre). \emph{Exemple} : le discours du Commandant lors de la cérémonie de Doum.
\item Argument 2 : elle flatte l'orgueil de Meka et lui donne une importance nouvelle dans le village. \emph{Exemple} : la fierté de Meka en boubou neuf, les préparatifs et la fête organisée en son honneur.
\end{itemize}

\textbf{II. La médaille, un piège qui divise, humilie et détruit.}
\begin{itemize}
\item Idée directrice : derrière l'honneur, la médaille sert la propagande coloniale et perd Meka.
\item Argument 1 : elle divise la communauté de Doum et isole Meka. \emph{Exemple} : la méfiance et les jalousies des notables et des anciens du village.
\item Argument 2 : elle conduit Meka à l'humiliation et à la perte de sa liberté. \emph{Exemple} : après la cérémonie, Meka, surpris par la pluie et le couvre-feu, est arrêté, battu et jeté en prison par les hommes du même Commandant qui l'a décoré.
\end{itemize}

\textbf{4. Introduction (modèle).} « La colonisation a souvent habillé la domination des habits de la gratitude (amorce). Dans \emph{Le Vieux Nègre et la médaille} (1956), le Camerounais Ferdinand Oyono raconte l'histoire de Meka, vieillard de Doum décoré de la médaille de l'amitié franco-camerounaise pour sa fidélité à l'administration (présentation de l'œuvre). Mais cet honneur est ambigu : la médaille est-elle une véritable récompense ou un piège tendu au colonisé (problématique) ? Nous verrons d'abord que la médaille apparaît comme une récompense qui flatte et honore, puis qu'elle se révèle un piège qui divise, humilie et détruit (annonce du plan). »

\textbf{5. Conclusion (modèle).} « La médaille, qui devait couronner la vie de Meka, n'a fait que précipiter sa chute : récompense en apparence, elle est en réalité l'instrument d'une domination qui use de l'honneur pour mieux asservir (bilan répondant à la problématique). Cette leçon dépasse le roman : ne faut-il pas toujours interroger le prix caché des cadeaux du pouvoir (ouverture) ? »

\textbf{Barème indicatif (sur 20)} : analyse du sujet 3 pts ; problématique 2 pts ; plan détaillé (cohérence, arguments, exemples précis) 7 pts ; introduction 4 pts ; conclusion 4 pts.

\textbf{Points de vigilance} : les exemples doivent être précis et exacts (l'arrestation de Meka \emph{après} la cérémonie) ; l'introduction comporte les quatre étapes dans l'ordre ; la conclusion répond explicitement à la problématique avant l'ouverture ; éviter la paraphrase du roman : chaque exemple sert un argument.
\end{corrige}

\courssub{Cas pratique : plan dialectique et intégration documentaire (Exercice 11)}

\begin{corrige}[Exercice 11 — Cas pratique : plan dialectique à partir d'un article de presse]
\begin{enumerate}
\item \textbf{Deux positions opposées} :
\begin{itemize}
\item Position 1 (le notable) : un chef qui abandonne les rites et la parole des ancêtres perd sa légitimité — la tradition est le fondement du pouvoir traditionnel.
\item Position 2 (le jeune cadre) : un chef qui refuse le développement et les outils modernes condamne son village au retard — la modernité est une nécessité.
\end{itemize}

\item \textbf{Plan dialectique détaillé} :

\textbf{I. Thèse : la tradition reste le fondement de la gouvernance locale.}
\begin{itemize}
\item Idée directrice : la légitimité du chef vient des ancêtres et des rites.
\item Argument 1 : les rites d'intronisation et les sociétés secrètes consacrent le pouvoir. \emph{Exemple} : « les rites d'intronisation... restent au cœur du pouvoir traditionnel » (document).
\item Argument 2 : la parole coutumière règle les conflits mieux que l'administration. \emph{Exemple} : la parole du notable : « Un chef qui oublie les ancêtres n'est plus un chef. »
\end{itemize}

\textbf{II. Antithèse : la modernité s'impose aux chefferies.}
\begin{itemize}
\item Idée directrice : les chefs doivent s'adapter aux outils et aux exigences du développement.
\item Argument 1 : les nouvelles technologies modernisent la gestion. \emph{Exemple} : les chefs de Foumban et Bafoussam tiennent une page Facebook et gèrent les états civils.
\item Argument 2 : la coopération avec l'administration renforce l'efficacité. \emph{Exemple} : les conflits fonciers réglés « avec l'aide des sous-préfets », les formations en gestion des notables.
\end{itemize}

\textbf{III. Synthèse : une gouvernance locale peut concilier tradition et modernité.}
\begin{itemize}
\item Idée directrice : la modernité est un outil, la tradition une légitimité ; l'une peut servir l'autre.
\item Argument 1 : le chef moderne reste légitime s'il respecte les rites tout en utilisant les outils modernes. \emph{Exemple} : un sultan qui gère l'état civil par ordinateur sans renier les cérémonies d'intronisation.
\item Argument 2 : le Cameroun offre des modèles de coexistence. \emph{Exemple} : les chefferies Bamoun qui conjuguent musée, patrimoine et administration numérique.
\end{itemize}

\item \textbf{Transition I $\rightarrow$ II (modèle)} : « Si la tradition fonde ainsi la légitimité des chefs, elle ne saurait à elle seule répondre aux besoins d'une population tournée vers le développement : c'est pourquoi la modernité s'impose aujourd'hui aux chefferies. »

\item \textbf{Paragraphe avec citation (modèle)} : « La modernisation des chefferies n'est plus un projet mais une réalité observable. À Foumban comme à Bafoussam, les chefs supérieurs utilisent déjà les outils numériques pour administrer. Mieux, la jeunesse elle-même réclame ce changement : \og Un chef qui refuse le développement condamne son village \fg{}, affirme un jeune cadre cité par l'article. Cette parole vigoureuse montre que la modernité n'est pas vécue comme une trahison, mais comme une condition de survie pour les communautés rurales : le chef qui s'y refuse perdrait, aux yeux de ses administrés, jusqu'à sa raison d'être. »
\end{enumerate}

\textbf{Barème indicatif (sur 20)} : reformulation des deux positions 3 pts ; plan dialectique en trois parties avec arguments et exemples 9 pts ; transition 3 pts ; paragraphe avec citation introduite et commentée 5 pts.

\textbf{Points de vigilance} : la synthèse ne doit pas être un simple compromis mais un dépassement (la modernité au service de la tradition) ; la citation s'introduit avec un verbe de parole (« affirme », « confie », « rétorque ») et des guillemets, puis se commente (on ne la laisse jamais « parler seule ») ; les exemples doivent être tirés du document ou de la réalité camerounaise, pas inventés hors sujet.
\end{corrige}

\newpage

\coursec{Fiche bilan}

\begin{popfigure}
\begin{tikzpicture}[mindmap, grow cyclic, every node/.style={concept, font=\small},
concept color=popblue!60, text=popdark,
level 1/.append style={level distance=3.6cm, sibling angle=51.4},
level 2/.append style={level distance=2.2cm, sibling angle=45},
root concept/.append style={concept color=popblue!80, font=\bfseries\small}]
\node[root concept]{Dissertation : plan détaillé}
[clockwise from=90]
child[concept color=poporange!60]{ node{Analyser le sujet}
  child{ node{Mots-clés} }
  child{ node{Type de sujet} }
}
child[concept color=popgreen!60]{ node{Problématique}
  child{ node{Question directrice} }
}
child[concept color=poppurple!60]{ node{Recherche des idées}
  child{ node{Brainstorming} }
  child{ node{Arguments + exemples} }
}
child[concept color=poppink!60]{ node{Plan détaillé}
  child{ node{2 ou 3 parties} }
  child{ node{Thèse / antithèse} }
}
child[concept color=popgold!70]{ node{Cohérence}
  child{ node{Progression thématique} }
  child{ node{Référent et substituts} }
}
child[concept color=popblueL]{ node{Rédaction}
  child{ node{Intro / développement / conclusion} }
  child{ node{Citations et transitions} }
}
child[concept color=popmuted!70]{ node{\emph{Le Vieux nègre et la médaille}}
  child{ node{Paratexte, lectures méthodiques} }
};
\end{tikzpicture}
\legende{Carte mentale de la leçon : les sept compétences clés autour de l'élaboration d'un plan détaillé de dissertation.}
\end{popfigure}

\begin{definition}[Définitions clés]
\begin{itemize}
\item \cle{Dissertation} : exercice écrit qui développe de manière argumentée et organisée une réflexion à partir d'un sujet et d'une problématique.
\item \cle{Problématique} : question centrale (souvent formulée en une ou deux questions) qui naît de l'analyse du sujet et oriente tout le devoir.
\item \cle{Plan détaillé} : schéma écrit du devoir, rédigé au brouillon, indiquant pour chaque partie la thèse défendue, les arguments et les exemples (citations).
\item \cle{Thèse / antithèse / synthèse} : point de vue admis, point de vue opposé, puis dépassement qui concilie les deux ou tranche.
\item \cle{Référent} : réalité désignée par un mot ; ses \cle{substituts} (pronoms, synonymes, hyperonymes) évitent les répétitions.
\item \cle{Progression thématique} : enchaînement des thèmes et des propos qui assure la \cle{cohésion} (liens formels du texte) et la \cle{cohérence} (logique du sens).
\item \cle{Sigle} : abréviation formée d'initiales (ONU, MINESEC) ; \cle{emprunt} : mot venant d'une langue étrangère (\emph{week-end}, emprunté à l'anglais).
\item \cle{Texte explicatif} : texte qui rend compte de causes et de conséquences (connecteurs : \emph{parce que, c'est pourquoi, ainsi}).
\end{itemize}
\end{definition}

\begin{methode}[Les cinq étapes de la dissertation]
\begin{center}
\begin{tabularx}{\linewidth}{|l|X|}
\hline
\rowcolor{popblueL} \textbf{Étape} & \textbf{Conduite à tenir} \\
\hline
1. Analyser le sujet & Souligner les mots-clés, définir chaque terme, repérer le type de sujet (question directe, citation, affirmation). \\
\hline
2. Formuler la problématique & Transformer la tension du sujet en une question directrice précise. \\
\hline
3. Chercher les idées & Brainstorming, puis tri : garder 2 ou 3 axes, chacun avec argument + exemple + citation. \\
\hline
4. Bâtir le plan détaillé & Rédiger au brouillon : introduction complète, intitulés de parties et sous-parties, transitions, conclusion. \\
\hline
5. Rédiger & Introduction (accroche -- sujet -- problématique -- annonce du plan), paragraphes argumentés, conclusion (bilan + ouverture). \\
\hline
\end{tabularx}
\end{center}
\end{methode}

\begin{aretenir}[Rappels de langue]
\begin{itemize}
\item Citation insérée entre guillemets, introduite par un verbe de parole (\emph{Oyono écrit : « \dots »}).
\item Connecteurs utiles : \emph{d'abord, ensuite, enfin} ; \emph{cependant, néanmoins} ; \emph{donc, ainsi, c'est pourquoi}.
\item Une idée = un paragraphe ; chaque paragraphe s'ouvre par une phrase d'annonce et se clôt par un bilan partiel.
\end{itemize}
\end{aretenir}

\begin{attention}[Pièges fréquents à l'examen]
\begin{itemize}
\item Hors-sujet : chaque partie du plan doit répondre directement à la problématique, pas au sujet en général.
\item Citation approximative ou inventée : une citation fausse est sanctionnée ; mieux vaut une paraphrase exacte.
\item Plan déséquilibré : une partie très longue et une autre réduite à une ligne fragilise l'ensemble du devoir.
\item Répétitions du même nom : utiliser les substituts du référent (pronoms, synonymes) pour assurer la cohésion.
\end{itemize}
\end{attention}

\begin{aretenir}[Checklist avant le Bac]
\begin{itemize}
\item $\square$ J'ai souligné et défini tous les mots-clés du sujet.
\item $\square$ J'ai identifié le type de sujet et le type de plan attendu.
\item $\square$ Ma problématique est une question claire, issue du sujet.
\item $\square$ Mon plan comporte 2 ou 3 parties équilibrées, sans hors-sujet.
\item $\square$ Chaque partie contient au moins un argument et un exemple précis.
\item $\square$ Mes citations sont courtes, exactes et entre guillemets.
\item $\square$ Mon introduction suit les 4 étapes : accroche, sujet, problématique, annonce du plan.
\item $\square$ Mes transitions relient explicitement les parties entre elles.
\item $\square$ J'ai évité les répétitions grâce aux substituts du référent (pronoms, synonymes).
\item $\square$ Ma conclusion fait le bilan et propose une ouverture.
\item $\square$ J'ai relu pour corriger orthographe, accords et ponctuation.
\item $\square$ J'ai géré mon temps : brouillon (plan détaillé) puis rédaction puis relecture.
\end{itemize}
\end{aretenir}

\findecours

\end{document}
